\quellseite{509}
\chapter[B. Auslegung der zehn Gebote]{B. Auslegung der zehn Gebote aus dem 19. und 20. Capitel des zweiten Buchs Mosis, gepredigt zu Wittenberg durch D. Martin Luther.\footnote{Diese Auslegung gehört, wie aus § 22 dieser Schrift ersehen wird, demselben Predigtcyclus an wie die vorhergehende Schrift, und wird der Zeit nach unmittelbar darauf folgen. Der Bauernaufruhr wird in derselben mehrfach, als etwas Vergangenes erwähnt, z. B. Cap. 20, § 120. Im Druck erschien dieselbe zuerst im Jahre 1528 zu Wittenberg, bei Hans Weiß unter dem Titel: „Auslegung der zehn Gebot, durch Mart. Luther gepredigt zu Wittenberg. Aus dem 19. und 20. Capitel des andern buchs Mosi, sampt einer Unterricht wie Moses zu leren ist. Mart. Luther. Wittenberg 1.5.28.“ Eine andere Ausgabe erschien ohne Angabe der Zeit bei Georg Rhau zu Wittenberg; eine dritte zu Erfurt im Jahre 1529. In den Sammlungen: in der Wittenberger (1566), Bd. V, Bl. 203; in der Jenaer (1566), Bd. IV, Bl. 486b; in der Altenburger, Bd. IV, S. 580, in der Leipziger, Bd. III, S. 609 und in der Erlanger, Bd. 36, S. 1. Wir geben den Text nach der Erlanger Ausgabe, welche den Urdruck bringt, unter Vergleichung der Wittenberger und der Jenaer.}}

1526. Gedruckt 1528.

In diesen zweien Capiteln werdet ihr hören, wie GOtt dem Volk Israel die zehn Gebote gegeben, und es damit geübt hat, wie folgt.\footnote{Hier wird Luther den Text des 19. und 20. Capitels des zweiten Buchs Mosis verlesen haben.}

1. Ich will aber\footnote{„aber“ fehlt in der Wittenberger.} treulich gemant und ermahnt haben alle die, so andere Leute lehren wollen, und GOttes Wort führen und predigen, daß sie sich wohl mit allem Fleiß und Ernst versehen, daß sie ja den Mosen recht lehren und dem Volke vortragen; nämlich, wo er Gebote gibt, wo er was fordert und treibt, daß man ihn da lasse einen Meister, Lehrer und Gesetzgeber der Juden bleiben, und ihn nicht treibe auf die Heiden oder Christen, ihr Gewissen damit nicht mehr\footnote{Walch und Erlanger: wieder.} verstricke noch verwirre. Denn er uns nicht weiter angeht, denn als fern er mit dem natürlichen Gesetze überein kommt. Wir haben unsern Meister Christum Jesum, der uns vorgehalten hat, was wir wissen, halten, thun und lassen sollen.

2. Das ist aber wohl wahr, daß Mose über die Gesetze auch schreibt schöne Exempel, beide des Glaubens und Unglaubens, der Strafe der Bösen und Gottlosen, und Erhöhung der Gläubigen. Ueber das alles, welches das Beste in Mose ist, sind die freundlichen, lieblichen Zusagungen und Verheißungen, derer wir uns annehmen sollen, zu welchen wir auch gehören, welche uns betreffen und belangen, als, vom Samen, in welchem alle Heiden sollen den Segen haben etc. [1 Mos. 22, 18.] Deß ein Exempel im Lucä haben wir, Cap. 17, 14., von den zehn Aussätzigen. Daß sie nun Christus heißt zu den Priestern gehen, und sie sich den Priestern erzeigen, geht mich nicht an, denn er redet da nach Mosis Gesetz, welches uns nicht bindet. Das Exempel aber ihres Glaubens geht mich an, daß ich Christo, wie sie, glaube, und, wie der einige Samaritan [Luc. 17, 18. 19.] Christo die Reinigung und Wohlthat zuschreibt, lobt, preist und dankt.

3. Davon ist nun oben, in der Vorrede [über das erste Buch Mosis\footnote{Die eingeklammerten Worte haben wir aus der alten Ausgabe Walchs herübergenommen. Die Jenaer Ausgabe verweist am Rande auf den „Unterricht“, welchen wir Col. 2 dieses Bandes mitgetheilt haben. — In der Wittenberger Ausgabe folgt unsere Schrift unmittelbar nach den Predigten über das erste Buch Mosis.}] genau in die Breite und Länge geredet, und ist wohl zu merken, denn es liegt die Macht daran. Denn es haben viel großer, trefflicher Männer darin gefehlt, und auch jehund viel großer Prediger, die sich nicht recht vorsehen, stoßen sich hart an diesem Ort, wissen nicht den Mosen recht zu predigen, wollen aus Christo einen Mosen, aus dem Evangelio ein Gesetzbuch, aus dem Worte Werke machen; spannen die Pferde hinten an Wagen, wissen selber nicht, wovon sie sagen; sind toll und thöricht, rasen und toben, wie die rasenden Hunde; schreien und plaudern in das Volk: GOttes Wort, GOttes Wort; verführen die armen Leute, und stoßen in die Grube. Es haben viel gelehrter Leute, Heilige und Väter
\quellseite{510}
nicht gewußt, wie weit Moses ginge, wie er sollte gelehrt werden. Origenes, Hieronymus, und ihres Gleichen, haben nicht gewußt, wie weit er dient.

\subsection{Vor GOtt gilt allein der Glaube}

4. Nun wollen wir den Text handeln, darin wir ein fein Exempel sehen, welches uns GOtt vor Augen stellt zum Spiegel, wie der Mensch so ungeschickt ist vor GOtt zu treten oder handeln, wenn er allein hat die Gebote, Vernunft, menschliche Weisheit, Klugheit, Sinn und Witz, und freien Willen, als, so er ihnen gebeut [2 Mos. 19, 10. 11.], sie sollen drei Tage zuvor, ehe denn sie auf den Berg steigen, ihre Kleider waschen, heilig und bereit sein auf den dritten Tag. Und thut Moses noch eines hinzu (daß sie ja äußerlich heilig und rein genug seien), daß sie sich enthalten von ihren Weibern.

5. Wir, die das Evangelium wissen, hören und gelernt haben, müssen, daß zu der gewissen und innerlichen Heiligkeit, die vor GOtt soll bestehen, nicht dient, ob das Kleid rein, besudelt, befleckt, oder unrein sei; doch hat GOtt zur selbigen Zeit wollen, daß sie geschmückt und rein wären, in schönen reinen Kleidern hergingen. Solche äußerliche Reinigkeit war ihre Heiligkeit zur selbigen Zeit. Item, daß sie sich drei Tage enthielten von ihren Weibern.

6. Im neuen Testamente geht es aber anders zu. Da greift Christus tiefer hinein, nämlich in das Herz, Bein und Mark; da gilt nicht mehr äußerliche Reinigkeit; der Schatten ist vorüber, der rechte Leib ist an die Statt gekommen. GOtt fragt nicht mehr darnach, ob das Kleid besudelt sei oder nicht, ob der Leib rein oder unrein, beschmiert oder nicht beschmiert sei. Paulus hatte Ketten an den Füßen, da er im Gefängniß war, Apost. 16, 24., hatte im Kerker nicht fast seine Kleider können waschen. Die Märtyrer waren voller Wunden, blutrünstig, mit Schweiß begossen, gingen in Hadern und bösen Kleidern herein.

7. Christo ward eine Ehebrecherin vorgebracht, die auf\footnote{So die Jenaer. Wittenberger und Erlanger: an.} frischer That ergriffen ward, Joh. 8, 3. f. (will geschweigen, daß sie sollte bei ihrem ehelichen Manne begriffen sein), noch schlägt er sie nicht von sich; der Ehebruch hindert da nicht, sondern er absolviret sie, und sagt sie los und ledig. Desgleichen das jüdische Fräulein bei dem Born, Joh. 4, 9. 26. 27., die von den Juden unrein war geachtet, wie sie sich selber wunderte, daß Christus trinken von ihr begehrt, und die Jünger wunderten sich auch, daß er mit ihr, als mit einer Heidin, redete: hat er sie doch nicht verachtet, sondern freundlich mit ihr geredet und getröstet, und eine schöne Disputation von dem zukünftigen Messia gehalten.

8. Im alten Testamente, in welchem alles in einer Figur geschah, hat er durch solche äußerliche Heiligkeit und Reinigkeit die Augen und das grobe Volk wollen lehren, zwingen und in Zaume halten, und also in die innerliche Heiligkeit weisen und leiten.

9. Daran aber liegt der rechte Knoten, daß Gott damit will angezeigt haben, daß vor ihm nicht hilft noch gilt einerlei\footnote{einerlei – einigerlei, irgend welche.} Heiligkeit, sie sei so groß als sie immer wolle. Denn was er gibt, das will er aus lauter Gnaden geben; er will den Preis allein behalten. Die Juden waren nicht so heilig, stark und mächtig, daß sie mochten GOtt hören reden, sondern waren ganz verzagt, meinten, sie müßten sterben, flohen von GOtt, und sprachen zu Mose: Rede du mit uns, und laß GOtt nicht mit uns reden, wir möchten sonst sterben.

10. Das ist ihr freier Wille, das ist ihre Reinigkeit und Heiligkeit. So GOtt anfähet zu reden, übereilt sie die Furcht, stecken ganz im Verzagen und Todesnöthen; fliehen von GOtt, bei welchem sie allen Trost sollten suchen, und laufen zu einem armen Sünder, der, wie sie, Fleisch und Blut ist; suchen mehr Freundschaft, Trost und Liebe bei einem Menschen denn bei GOtt. Was ist das für eine Heiligkeit und Reinigkeit, die sich nichts Gutes zu GOtt versieht? Wäre sie gerecht, so hielte sie den Stich, scheute sich nicht vor GOtt. Aber diese\footnote{So die Jenaer. Wittenberger und Erlanger: die.} Heiligkeit spricht zu GOtt unverschämt: Ich will dein nicht, mag dich nicht hören, ich will Mosen hören, Mosen habe ich lieber.

11. Aus, zum Teufel mit solcher Heiligkeit! Des Menschen Herz sollte sich ja mehr Gutes versehen zu GOtt, sollte mehr Hülfe und Trost da suchen, denn bei allen Engeln, heiligen Menschen und Creaturen; wie denn ein rechtschaffener Christ thut. Aber die heiligen Heuchler fliehen von GOtt, wollen ihn nicht haben. Was ist das anders geredet, denn zu Gott sprechen:

\quellseite{511}
Du bist ein Stockmeister, Richter und Henker, der hinten her kommt mit einer Keule, wolle sie an den Kopf schlagen. Wenn ich GOtt nicht leiden kann, der so freundlich ist, der unser Heiland und Seligmacher ist, welchen Teufel kann ich denn leiden?

12. Da ist uns abgemalt die Natur, wie sie verdammt ist, und vor GOtt flieht. Wo bleibt nun der freie Wille und die Menschenkräfte? Wo bleiben nun die Menschengesetze und Werkprediger, die zum Menschen sprechen: Thue dein Vermögen, wir müssen traun auch das Unsere dazu thun, schicke dich zu der Gnade GOttes, so wirst du ein Kind der Seligkeit? Ja wohl, ja, es ist unbereit und ungeschickt mit unserm Vermögen; wenn es an das Treffen geht, so hält es den Stich nicht.

13. Mit diesem Texte mögen wir zu Boden stoßen allen freien Willen. Wie wollten sie sich mit GOtt versöhnen, den sie nicht mögen hören, fliehen von ihm zu einem Menschen, versehen sich mehr Treue und Liebe zu einem Sünder denn zu GOtt. Wie ein fein Versöhnen ist mir das! Nimm deß ein Gleichniß: Wenn ein Sohn seinen Vater nicht hören wollte, der ihm doch freundlich wäre, sondern kehrte sich zu einem Knecht, der ein böser Bube wäre, suchte mehr Freundschaft bei dem Knechte, denn bei dem Vater, ohne Zweifel, es würde den Vater sehr verdrießen, und ihm wehe thun, daß der Sohn aller väterlicher Treue vergesse, und kehre sich zu einem bösen Buben. Also ist es hier auch. Aber damit ist es angezeigt, daß Natur und freier Wille vor GOtt nicht kann und vermag bestehen; denn sie fürchten sich, er schlage mit der Keule drein, halten ihn nicht anders denn für einen Teufel, Henker und Stockmeister, der nichts könnte denn zürnen.

14. Daraus schließen\footnote{Erlanger: beschließen.} wir das: Wo der Heilige Geist dem Herzen nicht vorbildet die evangelische Verheißung, so hält der Mensch GOtt für einen Teufel, Henker, Stockmeister und Richter, versieht sich mehr Gutes zu dem ärgsten Buben denn zu GOtt. Wo nun ein solch Herz flieht von GOtt, und kann ihn nicht leiden, so schließ nur frei, daß es GOtt lästert, und des Teufels ist, wie es steht und geht. Es gilt nicht fliehen, sondern hinzutreten, denn er ist ein Vater. Welcher nun fleucht, der wird in die Hölle gejagt, und ist des Teufels Wildpret.

15. Nun, Moses schreibt hier, daß es die Allerheiligsten gethan haben, die sich drei Tage vorher bereitet haben. Wenn es ein Sünder thäte, der voll Sünden steckt, ginge es wohl hin; aber die lieben Heiligen enthalten sich drei Tage vorher von ihren Weibern, daß doch nicht unrecht wär; thaten, was sie mochten, wuschen sich und schickten sich auf das beste, mit allerlei Werken, Fasten, und was sie nur konnten thun. Eben dieselben lieben Heiligen, die nicht mit Hurerei befleckt waren, rein von eigenen Weibern, geschmückt auf das beste und höchste, so die Vermuthung vermag, die konnten GOtt nicht hören, waren verzweifelt und verzagt an ihrer eigenen Heiligkeit. Darum ist das ein trefflicher Text, der zu Boden stößt alle Heiligkeit etc. Damit ist also angezeigt, daß wir Buben in der Haut sind, wo wir am besten und am heiligsten sind etc.

16. Also haben sich auch unsere Geistlichen, die heiligen Leute, mit großer Heiligkeit beläutet und getrieben, daß sie drob toll und unsinnig sind worden; und ist doch kein blöder und verzagter Volk denn sie. [Sie] haben Keuschheit gelobt, härere Hemden angetragen, übel geschlafen, gefastet, daß ihnen der Odem gestunken hat, [sind] in einen Sack gekrochen, dies und das gethan. Das alles aber hat sie nicht geholfen. Deß wollen wir St. Bernhard zum Exempel nehmen, denn ich gerne vor anderen Mönchen anziehe. Denn, ist je ein frommer Mönch gewesen, so ist es St. Bernhard gewesen. Da er einmal so hart krank lag, daß er sich hatte Sterbens verwegen,\footnote{verwegen = ergeben, gefaßt gemacht auf.} der hatte nun alle seine Tage Jungfräulichkeit gehalten, marterte sich mit Fasten und Wachen, daß ihm der Odem stank; also, daß ihn seine Brüder nicht mochten um sich leiden, führte ein schönes Leben, und that viel schöner Predigten, steckte voll guter Werke. Da er nun sahe das Urtheil GOttes, sprach er: O HErr, ich habe mein Leben übel angelegt, und die Zeit schändlich zugebracht.

17. Was hast du gethan, lieber Bernharde, bist du doch ein Mönch gewesen, hast schöne Predigten gethan, bist kein Ehebrecher gewesen, hast niemand das Seine genommen? Wer reißt nun St. Bernhards seine Werke hinweg? Das Gericht GOttes. Denn wenn dies dem Menschen eröffnet wird, hat aller Menschen Gerechtigkeit seine Statt, mag nicht bestehen; wie David Ps. 143, 2. begehrt, daß GOtt nicht mit ihm
\quellseite{512}
in das Gericht trete; denn vor ihm möge kein lebendiger Mensch bestehen noch gerecht erfunden werden.

18. Wenn nun St. Bernhardus nicht bestanden ist vor dem Gerichte GOttes, wie wollten andere Mönche und Pfaffen bestehen? Ich halte St. Bernhardum höher, denn alle Mönche und Pfaffen auf Erden; ich habe seines Gleichen nicht gehört noch\footnote{Wittenberger und Jenauer: und.} gelesen. Darum will GOtt also dadurch anzeigen, daß alle Werke vor seinen Augen verloren sind, nämlich, wenn sie in der Meinung geschehen, daß man damit vor GOtt will bestehen.

19. Also muß auch St. Johannes der Täufer mit aller seiner Heiligkeit allen Werkheiligen zur Schande stehen, denn er bezeugt, er habe es alles von oben herab, und weist auf Christum, der die Sünde hinwegnehme, Joh. 1, 29. 3, 27., nicht sein heilig Leben. Ich sage nicht von den Werken, die dem Nächsten zugute geschehen, und aus dem Glauben folgen; allein, daß man nicht mit solchen Werken vor GOttes Gericht wolle kommen, GOtt damit zu versöhnen und Freund zu machen, wie die Werkheiligen, die sich hervorthun, bringen ihre Werke, sind damit umhängt, als mit großen Hopfensäcken, dürfen zu GOtt sprechen: Zahle mich; als ob er ein Trödler sei, der sein Himmelreich verkaufe.

20. Das ist aber nicht allein von den Juden geschrieben, daß sie nicht hinan durften, sondern auch von allen Werkheiligen, die auf das heute gerüstet sind mit Werken; noch werden sie verjagt, wenn Gottes Gericht herankommt. Wie viel minder dürfen hinzukommen die mit Sünden beladen sind. Es ist viel gesagt. Welcher nicht mehr hat denn die guten Werke, der trabe nur gen Hölle, er wird nichts erlangen; es muß was Größeres sein, denn unser heilig Leben. Folgt der Text des 19. Capitels.\footnote{Die letzten Worte bilden in den alten Ausgaben die Ueberschrift des Folgenden.}

\section{Das neunzehnte Capitel}

\begin{schriftzitat}
V. 1. 2. Im dritten Monden nach dem Ausgange der Kinder Israel aus Egyptenlande, kamen sie dieses Tages in die Wüste Sinai. Denn sie waren ausgezogen von Raphidim, und wollten in die Wüste Sinai, und lagerten sich in der Wüste daselbst gegen dem Berg. Und Mose stieg hinauf zu Gott.
\end{schriftzitat}

1. Moses ist hier ein Mittler, welcher die Worte des Volks zu GOtt trägt, und die Worte GOttes zu dem Volke. Das zeigt auch Paulus an, Gal. 3, 19., da\footnote{Wittenberger und Erlanger: so.} er redet von dem Gesetze: „Und ist überreicht (spricht er) von den Engeln durch die Hand des Mittlers.“ Wir haben aber jetzund GOtt selber zu einem Mittler, wie St. Paulus spricht 1 Tim. 2, 5.

\begin{schriftzitat}
V. 3. Und der HErr rief ihm vom Berge und sprach: So sollst du sagen zum Hause Jacob, und verkündigen den Kindern Israel etc.
\end{schriftzitat}

2. Das ist die Vorrede und der Eingang zu dem Volke, damit GOtt vor dem Gesetze ein Gebing\footnote{Gebing = Uebereinkunft, Vertrag.} macht mit dem Volke; als wollte er sprechen: Liebes Volk, ehe ich mit euch handle, will ich zuvor mit euch reden, ob ihr das wollt annehmen und halten. Wie bei uns, so man in einer Stadt ein Gesetz machen will, berufet man das Volk, fragt die Gemeinde, ob sie das Gesetz wollen annehmen. Also thut hier der HErr auch; er will sie nicht zwingen, daß sie es annehmen, wo sie es aber annehmen, so will er, daß sie es sollen halten, wie im Texte folgt:

\begin{schriftzitat}
V. 4. Ihr habt gesehen, was ich den Egyptern gethan habe, und wie ich euch getragen habe auf Adlers Flügeln, und habe euch zu mir gebracht.
\end{schriftzitat}

3. Es gilt nun: Ich habe euch nicht allein zu mir gebracht, sondern ich habe euch auch auf Adlers Flügeln getragen, ihr habt die Gutthat hin, ihr habt die Mirakel gesehen, wie ich euch aus Egypten geführt habe, in der Wüste ernährt, erhalten, beschützt, mit Schuhen und Kleidern versorget. Nun habt die Wahl, ob ihr
\quellseite{513}
solchen Gott haben wollet, oder einen andern erwähhlen. Folgt nun weiter:

\begin{schriftzitat}
V. 5. Werdet ihr nun meiner Stimme gehorchen, und meinen Bund halten, so sollt ihr mein Eigenthum sein vor allen Völkern etc.
\end{schriftzitat}

4. Diesen Spruch sollte man unterziehen mit rother Tinte, daß man könnte den Rottengeistern antworten. Er spricht: Er wolle sie zu einem Volk haben, zum Eigenthum. Sprich du zum Rottengeistern: Lege die Brille auf die Nase, siehe den Text recht an, er redet die Worte nicht zu den Heiden, sondern allein zum Volke Israel. Er spricht,\footnote{Wittenberger: spricht nicht.} Israel sei sein Volk. Bin ich nun aus dem Volke Israel, so will ich thun, was Moses heißt. Die Worte aber geben da einen Unterschied. Es sind viel verführt worden mit dem Mose; wer nun noch kann wehren, der wehre, es ist hohe Zeit. GOtt gebe, daß die hohen, stolzen Geister lernen Mosesen recht verstehen; sie aber laufen an. GOtt gebe, daß sie nicht weiter einreißen.

\begin{quote}
Werdet ihr nun meiner Stimme gehorchen, und meinen Bund halten.
\end{quote}

5. Das ist, werdet ihr mein Gesetz, das ich euch geben will, annehmen und halten,

\begin{quote}
So sollt ihr mein Eigenthum sein, vor allen Völkern; denn die ganze Erde ist mein. [V. 6.] Und ihr sollt mir ein priesterlich Königreich, und ein heiliges Volk sein.
\end{quote}

6. Das ist, ich will euch zu einem besondern Volke annehmen, das ich äußerlich regieren will mit Gesetzen, und kein ander Volk; wie auch David im Psalter spricht, Ps. 147, 19. 20.: „Er zeiget Jakob sein Wort, Israel seine Sitten und Rechte. So thut er keinen Heiden, noch läßt sie wissen seine Rechte.“ Denn zu derselben Zeit ließ er alle Heiden gehen, wie sie gingen; wiewohl er zu Zeiten aus den Heiden etliche hervor zog, als, in Arabia den Hiob, und den aussätzigen Naeman in Syria etc. Aber er hat nie keinen solchen Haufen abgesondert, wie er mit Israel hat gethan. Darum spricht er: Wenn ihr halten werdet meinen Bund, so sollt ihr mein Eigenthum sein, und ich will euer GOtt und HErr sein, will mitten unter euch wohnen; wiewohl das Erdreich mein ist, so will ich doch euch erwählen, will etwas Besonderes mit euch anrichten, und ihr sollt mir ein priesterlich Königreich und ein heilig Volk sein. So sprach nun GOtt: Mose, gehe hin, sage es dem Volke, beut es ihm an, was es thun wolle. Das war der Rathschlag und Theiding,\footnote{Wittenberger und Jenaer: „teibing“. Erlanger: „Theiding“. Nach der letztern Lesart des Originals liegt die Vermuthung nahe, daß „Gebing“ zu lesen sein möchte. Vgl. § 2 zu Anfang und § 21. Sonst wäre hier Theiding = Rede zu nehmen, wie in der St. Louiser Ausgabe, Bd. VIII, 978, § 23.} ob sie der Sache möchten eins werden.

7. Wie ein feiner Orator und Redner ist GOtt! Zum ersten erzählt er seine Wohltat, die er ihnen bewiesen hat; darnach stellt er es ihnen heim, ob sie ihn wollen annehmen, für einen GOtt erkennen, seinen Bund halten. Hält ihnen vor, wie er sie aus Egypten geführt habe. Das ist die erste Wohltat.

8. Die andere, daß er sie, wie ein Adler, auf Flügeln geführt habe; welches Moses auch anzeucht in seinem Gesange im fünften Buche Cap. 32,\footnote{Wittenberger und Erlanger: Cap. 33.} 11.: „Wie ein Adler aufweckt sein Nest,\footnote{Statt: „aufweckt sein Nest“ bietet die Bibel: „aufführet seine Jungen“; hier ist dasselbe gleich folgend erklärt durch: „seine Jungen lehret fliegen.“} und über seinen Jungen schwebt. Er breitete seine Fittige aus, und nahm ihn, und trug ihn auf seinen Flügeln.“ Es ist eine besondere Weise zu reden: Wie der Adler seine Jungen lehret\footnote{Erlanger: lernet.} fliegen, also habe ich mit euch gehandelt. Ich habe euch aus Egypten geführt, in der Wüste mit Himmelsbrod gespeist, aus dem Felsen Wasser gegeben, die Schlangen vertrieben, eure Schuhe und Kleider unversehrt behalten, erhalten und beschützt vor euren Feinden, habe mit euch gefahren wie ein Vater mit seinem Kinde; das habe ich euch gethan, daß ich euch lehrete fliegen, das ist, daß ihr mir nachfolget, und mir glaubet.

9. Wenn ihr also aus diesem Exempel beweget werdet, mich zu einem GOtt anzunehmen, so will ich euch annehmen zu einem priesterlichen Königreich, das ist, ihr sollt alle Priester und Könige sein. St. Petrus 1. Ep. 2, 9. führt diesen Spruch, und deutet ihn auf alle Christen. Alle Christen sind Könige und Pfaffen, davon ihr genug gehört habt an andern Orten, und in\footnote{„in“ fehlt in der Erlanger.} Büchern gelesen.

\quellseite{514}
\subsection{König sein}

\begin{schriftzitat}
V. 6. Ihr sollt mir ein priesterlich Königreich sein.
\end{schriftzitat}

10. Warum führt aber Moses diesen Spruch allein auf die Juden? Antwort: Dieser Spruch ist auf die Juden gangen, ehe sie das Gesetz empfangen haben: Wenn ihr werdet halten meine Gebote, und nicht zerreißen meinen Bund, so werdet ihr Könige und Priester sein. Also, daß sie noch nicht waren ein Gesetzvolk; waren gleich denen, so nach dem Gesetze gläubig sind, welchen kein Gesetz gegeben ist. Welcher nun GOtt glaubt ohne alles Gesetz, der wird ein Priester und König, er sei Jude oder Heide, vor oder nach dem Gesetz.

11. Sie sind aber nicht so tolle Könige, wie die Könige der Welt sind. Denn dieselbigen sind nicht die rechten Könige, sie sind nur Zahlpfennige und gemalte Könige gegen den Gläubigen, denn sie regieren nur zeitlich und äußerlich. Die Gläubigen aber sind rechtschaffene Könige; nicht, daß sie eine güldene Krone auf ihrem Haupte tragen, ein gülden Scepter führen, in Seide, Sammet und güldenen Stücken und Purpur einher treten, sondern, das viel herrlicher ist, sie sind Herren über Tod, Teufel, Hölle und alles Unglück. Die weltlichen Könige können allein mit Gold, Silber, Geld und Gut umgehen, Reichthum und Gewalt haben, die Leute würgen und plagen, ihre Unterthanen schaben, schinden und schaben, aber ihnen selber können sie nicht helfen, dem allermindsten Geschwär an einem Finger können sie nicht wehren, oder, daß ihnen der Bauch, Haupt und Glieder nicht wehe thun. Viel weniger können sie wehren der Sünde, dem Tode, Teufel, Hölle, Krankheit, Unglück etc. Sind derhalben Könige, wie Zahlpfennige Gulden, und gemalte Könige auf der Karte Könige sind.

12. Denn den rechten gläubigen Königen, die im Reiche Christi sind, ist Schande eine Ehre, die Hölle das Himmelreich, der Tod das Leben, der Teufel ein fröhlicher Mann, die Sünde Gerechtigkeit, Unglück Glück, Armuth Reichthum etc. Denn sie sind Herren über das alles, fragen nach niemand, weil sie GOttes Kinder sind, und GOtt zu einem Freund, ja, lieben Vater haben [Röm. 9, 26.], bei welchem sie Reichthum, große Schätze und alle Güter, und derselbigen die Fülle finden. Darum kann ihnen keine Sünde, Tod, Teufel, Hölle, Hunger, Durst, Kälte, Hitze, Schwerdt, noch alles Unglück nicht schaden; ja, in diesem allem überwinden sie weit, und finden in dem allem das Widerspiel: in Armuth Reichthum, in der Sünde die Gerechtigkeit, in der Schande große Ehre, in Hunger und Durst alle Fülle, wie gesagt ist. Also, weil sie solch innerlich Reichthum haben, achten sie der weltlichen Pracht nicht, verachten die güldenen Kronen, Purpur, Seide, Gold, Silber, Geld und Gut. Eine güldene Krone, Purpur und Schmuck gehört den Königen auf dem Kartenspiel, es ist den Kindern GOttes alles viel zu schlecht.

13. Man soll auch GOtt nicht also abmalen, als sitze er droben im Himmel auf einem güldenen Stuhle, in Purpur und güldenen Kleidern, mit einem güldenen Scepter und Krone, sondern als einen HErrn, der alle Dinge erhalte, alles erfülle, in allen Creaturen gegenwärtig sei, alles in allen Dingen schaffe und thue, alles in seiner Hand habe; der ein HErr sei über Sünde, Tod, Teufel, Hölle, alles Unglück und Uebel. So sind nun alle die, die ihm trauen und glauben, rechte Könige. Darum sie auch gleiche Gewalt haben mit ihrem Vater, erben alle seine\footnote{Wittenberger: „Erben aller seiner“.} Güter, doch noch in der Hofnung, Röm. 8, 17. 1 Joh. 3, 1. 2.

14. So sollen nun die Worte Mosis, da er spricht: Wenn ihr werdet halten meine Worte etc., sollt ihr Priester und Könige sein, nicht verstanden werden von den weltlichen Königen, sondern innerlichen und geistlichen Königen im Glauben, die im Reiche Christi sind, in welchem sie allzumal Könige sind, gleiche Gewalt habe mit dem obersten Könige Jesu Christo, in welchem\footnote{Jenaer: welches.} Reiche ein jeglicher für sich selber König ist. Wenn ihr nun werdet halten, spricht er, meinen Bund, so werdet ihr nicht weltliche Könige, sondern geistliche sein, und dazu Priester, das ist, Bischöfe und Päpste. Gott verzeihe mir, daß ich den Namen hier nenne, ich sollte ihn ja nicht nennen!

\subsection{Priester sein}

15. Was ist nun ein Priester? In welches Munde GOtt sein Wort legt, wie Malachias [Cap. 2, 7.] spricht:\footnote{„spricht“ fehlt in der Erlanger.} Labia sacerdotis custodiunt
\quellseite{515}
scientiam [des Priesters Lippen bewahren die Lehre]; item, er opfert und bittet für die andern. Ein solcher Priester im Glauben darf vor GOtt treten, für das Volk bitten, ihr Wort reden, und ihr Bestes bei GOtt vornehmen, darnach von GOtt herauszugehen zu dem Volk, ihnen GOttes Antwort und Befehl vorlegen. Und eines solchen Priesters Wort soll so viel gelten, als GOttes Wort selber, denn er führt nicht sein, sondern GOttes Wort, das ist denn eine Ehre über alle Ehre, wie St. Peter sagt [1. Ep. 2, 9.].

16. Das ist nun Summa Summarum dieses Spruchs: Wer mein Wort hat, und dem glaubt, der ist ein Priester; seid ihr nun mein Volk, so habt ihr den Glauben. Welcher nun den Glauben hat, der ist ein König und Herr über Sünde, Tod, Teufel, Hölle und alles Unglück, denn der Glaube allein setzt euch in solche Güter und Herrlichkeit. Wer den Glauben hat, der hat alle Dinge, vermag alles, überwindet alles [Röm. 8, 37. 38.], [es] kann ihm nichts schaden, weder Zeitliches noch Ewiges, auch die Pforten der Hölle nicht [Matth. 16, 18.].

17. Danach, welcher das Wort GOttes hat, der ist ein Priester, und wer ihn hört, hört GOtt selber. Also ist Priester und König sein nichts anders, denn den Glauben und den Heiligen Geist haben, die Gnade GOttes den andern zu predigen, und vor GOtt zu treten in einem guten Vertrauen, als ein Kind zum Vater. Es ist gering anzusehen, predigen, bitten und stehen in einem rechten Glauben; es ist aber vor GOtt ein groß, mächtig Ding, daß ein Mensch, ein armer, elender Madensack, zu einer solchen Ehre soll kommen.

18. Solche Ehre verheißt er ihnen, wo sie seinen Bund halten. Ueber das geistliche Reich und Priestertum hat er daneben auch eingesetzt ein leiblich Königreich und Priestertum, davon Moses hernach\footnote{Wittenberger: auch hernach.} handelt. Hier handelt er nur vom geistlichen Priestertum und Königreiche, und von dem Volke, das solchen Bund hielt. Welche aber ungläubig waren, und seinen Worten nicht gehorchten, die ging dieser Spruch nicht an, sondern waren unter dem leiblichen Königreiche. Doch hatte gleichwohl GOtt zu derselben Zeit etliche, die zugleich dem geistlichen und auch dem\footnote{„dem“ findet sich nur in der Jenaer.} leiblichen Königreiche und Priestertum unterworfen waren.

19. Da aber das Evangelium anging, hat er das leibliche Priestertum aufgehoben, und ist das geistliche Priestertum ausgeprediget worden in die ganze Welt durch die Apostel. Also, David war ein geistlicher und leiblicher König, war dennoch dem leiblichen Priestertum unterworfen, und war doch ein rechter geistlicher Priester; von welchem Priestertum er viel schöner Psalmen geschrieben hat.

20. Also ist nun ein jeglicher Christ ein König für sich selber, und ein Priester für andere. Das Priestertum ist höher denn das Königreich, es breitet sich weiter aus. Denn der Priester treibt das Wort nicht allein für sich, sondern für andere; den Glauben aber, durch welchen er erstlich König wird, hat er für sich allein.

\begin{quote}
Das sind die Worte, die du den Kindern Israel sagen sollst.
\end{quote}

21. Das ist das Gedinge und der Vertrag, den GOtt seinem Volk anbeut. Nun wird folgen die Antwort des Volks, wie sie der Sache eins werden.

\begin{schriftzitat}
V. 7. Mose kam und forderte die Aeltesten im Volke, und legte ihnen alle diese Worte vor, die der HErr geboten hatte.
\end{schriftzitat}

22. Droben ist gesagt, am Ende des 18. Capitels, wie Moses aus dem Rath seines Schwähers Jethro verordnet hatte Hauptleute; derer hat er etliche gesetzt über tausend, über hundert, über fünfzig, und über zehn, die heißt er hier die Aeltesten, und sind die Rathsherren. Er nennt sie aber nicht ohne Ursache die Aeltesten, denn die Aeltesten nahm man zu Häuptern. Denn die Jungen sind närrisch,\footnote{Wittenberger: närrisch und.} ungeschickt und unerfahren zum Regiment; wiewohl sie unterweilen fromm sind, so sind sie doch ungeschickt und untüchtig zum Regiment. Denn es ist nicht genug zum Regimente, fromm sein. Ein Esel ist auch fromm. Es gehöret dazu Geschicklichkeit und Erfahrung. Man findet einen, der fromm ist, kann doch kaum fünfe zählen. Wer regieren soll, dem muß [es] nicht fehlen an Vernunft, Weisheit, Klugheit, Sinn und Witz, will er anders nicht großen Schaden thun im Regimente. Denn das Regiment ist der Vernunft unterworfen.

23. Das haben auch erfahren und gebraucht die Heiden, beide Griechen und Römer, wie denn
\quellseite{516}
ihrer beider Sprache ausweist. Denn auf griechisch und lateinisch nennt man die Regenten, die den andern sollen vorstehen, presbyteros und senatores, die Aeltesten. Als wollten sie sprechen: Das Alter soll regieren, die Jugend soll regiert werden; wie auch St. Petrus [1. Ep. 5, 5.] ermahnt die Aeltesten, daß sie sollen vorstehen der Jugend, und die Jungen, daß sie den Aeltesten sollen gehorchen und unterthan sein [1 Tim. 5, 17.].

24. Wiewohl\footnote{Wittenberger und Jenaer: Wiewohl aber.} GOtt zu Zeiten wunderlich einem Jungen mehr Witz und Verstand gibt denn einem Alten, als Timotheo und Salomoni, so ist doch die natürliche Ordnung, daß der Witz bei den Alten ist, Unverstand bei den Jungen. Und GOtt\footnote{Wir halten dafür, daß diese Lesart der Ausgaben schwerlich richtig ist. Statt „Und GOtt“ möchte zu lesen sein: „Und so“ oder: „Und solches“.} hält auch die Welt dafür, daß ja die Alten sinnreicher seien denn die Jungen, daher ein gemein Sprüchwort geht: Man soll sich nicht an alte Seusel kehren, man fähet sonst Rohm. Und ein anderes: Wenn die Jugend flug wäre, und das Alter stark, so könnte man es nicht mit Golde bezahlen. Die Jugend ist stark, und das Alter verständig, sinnreich und witzig; darum soll die Jugend dem Alter gehorsam sein, daß sie der Stärke recht brauchen; so geht es recht zu.

25. Deß höre ein Exempel von einem Fuhrmanne und Pferden. Die Pferde, wiewohl sie stärker sind, müssen dem Fuhrmanne gehorchen, und wenn das geschieht, so geht der Wagen recht fort. Denn es liegt viel mehr am Fuhrmanne, wiewohl er schwächer ist denn die Pferde, denn an den Pferden; er ist der Klügste, weiset Steg und Weg. Wo aber die Pferde nicht hören, laufen über Stauden und Stöcke, so zerscheitern sie sich selber, den Wagen mitsammt dem Fuhrmanne. Also auch in der Welt geht es recht fort, wenn die Jugend dem Alter gehorsam ist. Denn die Jugend läuft leichtlich an mit ihrer Stärke; wo sie nicht in einem Zaum geführt wird, geht also alles zu Trümmern. Darum wird die Welt nicht wohl regiert, wo man mit dem Kopfe hindurch will unbesonnen,\footnote{In den alten Ausgaben: unversunnen.} burdi burdi, nur hindurch, wie die Sau durch den Koben.\footnote{Jenaer: Kobe. Wittenberger und Erlanger: Kober.} Es geräth nicht wohl. Wie es jetzund mit unsern tollen Fürsten geschieht.

26. Moses beruft nun die Aeltesten, denn es ward\footnote{Erlanger: war.} ihm zu viel, ein solch groß Volk, nämlich sechsmal hunderttausend Männer, zu verhören. Darum hat er das Regiment getheilt unter die Häupter, die es dem Volke sollten antragen, und wieder vom Volke Antwort bringen.

\begin{schriftzitat}
V. 8. Und alles Volk antwortete zugleich, und sprachen: Alles, was der HErr geredet hat, wollen wir thun.
\end{schriftzitat}

27. Fahret schön, lieben Gesellen, es ist zu hoch und zu viel vermessen. Das ist nun\footnote{„nun“ fehlt in der Wittenberger.} also zugegangen: Die Hauptleute, jeglicher besonders, haben ihre Unterthanen gefragt: ob sie den Bund wollen annehmen? und als sie alle zugleich verwilligt haben, und der Sache eingegangen, sind die Hauptleute wieder zu Mose gekommen, und ihm zugleich alle angezeigt, wie das Volk willig sei, den Bund anzunehmen. Da geht nun der Bund an; darnach hat es Moses vor GOtt getragen, als ein Mittler.

\begin{schriftzitat}
V. 9. Und Mose sagte die Rede des Volks dem HErrn wieder. Und der HErr sprach zu Mose: Siehe, ich will zu dir kommen, in einer dicken Wolke, auf daß das Volk meine Worte höre, die ich mit dir rede, und glaube dir ewiglich. Und Mose verkündigte dem HErrn die Rede des Volks.
\end{schriftzitat}

28. Moses hat zum Volke geredet: Ihr habt euch nun verbunden, ihr werdet den HErrn selbst hören reden; er wird in einer dicken Wolke kommen, das ist, mit Blitz,\footnote{Wittenberger: „und Donner“. — Da sich in der Wittenberger solcher unbedeutenden Varianten viele finden, werden wir sie forthin nicht mehr anmerken.} Donner, und grausamen Gewitter. Folgt nun der Text:

\begin{quote}
Daß das Volk meine Worte höre, die ich mit dir rede, und glaube dir (Mosi) ewiglich.
\end{quote}

29. Wie das sei zugegangen, wird hernach folgen. Das Volk hört GOtt und Mosen reden auf dem Berge. GOtt redet, und Moses antwortet ihm; und das Volk hörte also die Rede und Widerrede.

\begin{quote}
Daß es dir glaube.
\end{quote}

30. Das ist, daß du, Mose, einen Namen erlangest und ein Ansehen habest vor dem Volk, daß es wisse, du habest es selber von mir gehört
\quellseite{517}
, und nicht aus eigenem Kopfe und Vornehmen erdacht. Also sind sie der Sache eins worden; nun geht es an, wie folgt:

\begin{schriftzitat}
V. 10. Der HErr sprach zu Mose: Gehe hin zum Volke, und heilige sie heute und morgen.
\end{schriftzitat}

31. Das ist die dritte Botschaft GOttes durch Mosen zum Volke. Was aber das Heiligen sei, ist droben, am Anfange dieses Capitels, genug gehandelt. Folgt im Texte weiter:

\begin{schriftzitat}
V. 11. Daß sie ihre Kleider waschen, und bereit seien auf den dritten Tag. Denn am dritten Tage wird der HErr vor allem Volke herab steigen auf den Berg Sinai.
\end{schriftzitat}

32. Das ist, daß sie hübsch gerüstet, geputzt, und geschmückt seien, in ihren besten Kleidern; wie auch im 33. Capitel, V. 4. 5., angezeigt wird, da der HErr, erzürnt um ihrer Halsstarrigkeit willen, sie heißt ihren Schmuck ablegen, und gemeine, tägliche Kleider anziehen: das thaten sie, und trugen Leid, und niemand trug seinen Schmuck, darum, daß der HErr zornig war. Hier aber befiehlt der HErr, daß sie geziert und geschmückt vor ihm erscheinen.

\begin{schriftzitat}
V. 12. Und stecke Zeichen um das Volk her, und sprich zu ihnen: Hütet euch, daß ihr nicht auf den Berg steiget, noch sein Ende anrühret.
\end{schriftzitat}

33. Das ist aber ein Stück: Wenn sie nun geheiligt und geschmückt sind, so stecke ihnen Malzeichen mit Schranken um den Berg, daß sie nicht zu weit gehen. Nicht soll man es aber verstehen, daß das Volk habe fast hinzu geeilt, auf den Berg zu steigen; denn sie fürchteten sich, und wären lieber über alle Berge gewesen; sondern ehe der HErr kam, und herab stieg auf den Berg, wären sie hinzu gegangen. Denn da er herab stieg im Feuer und Donner, da war es nicht vonnöthen, daß er ihnen verbieten ließ, nicht hinzu zu gehen. Da flohen sie selber, wie sie der Teufel sagte; zitterten, waren vor großem Schrecken schier todt, und wären lieber über hundert Meilen gewesen, denn unten am Berge angestanden.

34. So ist nun die Meinung dieses Textes: Du sollst sie heiligen und verschränken, ehe denn ich herab steige auf den Berg; sonst würde das Volk so vorwitzig sein, daß es sehen wollte, was auf dem Berge wäre. So will ich den Berg geheiligt haben drei Tage; am dritten Tage will ich ihnen wohl wehren, daß sie nicht herzu treten. Denn wenn sie die Stimme, Donner, Blitz und die Posaunen werden hören, so werden sie sich nicht herzu dringen; sonst wäre das Volk\footnote{„Volk“ fehlt in der Wittenberger; in der Jenaer steht es als Conjectur am Rande.} wohl so vorwitzig geblieben, daß es jetzt hinauf ging, vor der Zeit. Also hat er ihnen eine Furcht eingejagt, daß er den Berg geehrt wollte haben, als eine Statt oder Ort, da GOtt hinkommen sollte. Gleich als wenn einem Hausvater angesagt würde: Halte dein Haus rein, über drei Tage wird ein Fürst zu dir einziehen, so würde er sein Haus rein halten, und dem Herrn eine schöne Wohnung bereiten, und verwahren, daß kein Schwein, Pferd, Ochs, oder Thier hinein ginge. Also wollte ihm GOtt den Berg heilig gehalten haben, daß niemand vor ihm hinzu käme.

\begin{schriftzitat}
V. 13. Denn wer den Berg anrühret, soll des Todes sterben. Keine Hand soll ihn anrühren, sondern er soll gesteiniget oder mit Geschoß\footnote{In den alten Ausgaben: „schoß“.} erschossen werden; es sei ein Thier oder Mensch, so soll er nicht leben.
\end{schriftzitat}

35. So hart verbeut er ihnen, daß sie nicht zum Berge nahen, daß er ihnen dräuet, daß, wer hinzu naht, soll gesteiniget werden etc. Das geschieht alles darum, daß er sie scheu und furchtsam mache. Sind das nun die heiligen Leute, die sich also geschmückt, geziert und geheiligt haben, ihre Kleider gewaschen, und in anderer äußerlicher Heiligkeit auf das allerbeste zugerichtet? Warum gehen sie nicht hinan? Sie dringen sich nichts darnach.

\begin{quote}
Wenn des Horns Ton tönet, so sollen sie auf den Berg steigen.
\end{quote}

36. Dann so thue die Malzeichen und die Schranken hinweg, und laß sie hinzu treten und hinauf steigen. Vorhin hat er sie mit Worten erschreckt, jetund wird es an die äußerlichen Zeichen gehen. Da seht ihr, wie der HErr das Volk in äußerlicher Weise eingezogen und getrieben hat, daß sie sich auch mußten vor dem Berge entsetzen, der doch nichts anders denn Stein und Erde war. Nun folgt, wie das Volk herzu kommt, und GOtt herab steigt auf den Berg etc.

\quellseite{518}
\begin{schriftzitat}
V. 14–25. Mose stieg vom Berge zum Volke, und heiligte sie, und sie wuschen ihre Kleider, und er sprach zu ihnen: Seid bereit auf den dritten Tag, und keiner nahe sich zum Weibe. Als nun der dritte Tag kam, und Morgen war, da hub sich ein Donner und Blitzen, und eine dicke Wolke auf dem Berge, und ein Ton einer sehr starken Posaune; das Volk aber, das im Lager war, erschrak. Und Mose führete das Volk aus dem Lager, GOtt entgegen, und sie traten unten an den Berg. Der ganze Berg aber Sinai rauchte, darum, daß der HErr hernieder auf ihn stieg, mit Feuer, und sein Rauch ging auf, wie ein Rauch vom Ofen, daß der ganze Berg sehr erschrecklich war. Und der Posaunen Ton ging, und ward stark. Mose redete, und GOtt antwortete ihm laut. Als nun der HErr hernieder kommen war auf den Berg Sinai, oben auf seine Spitze, forderte er Mosen oben auf die Spitze des Berges, und Mose stieg hinauf. Da sprach der HErr zu ihm: Steig hinab, und zeuge dem Volke, daß sie nicht herzu brechen zum HErrn, daß sie ihn sehen, und viel aus ihnen fallen. Dazu die Priester, die zum HErrn sich nahen, sollen sich heiligen, daß sie der HErr nicht zerschmettere. Mose aber sprach zum HErrn: Das Volk kann nicht auf den Berg Sinai steigen, denn du hast uns bezeuget und gesagt: Stecke Zeichen um den Berg und heilige ihn. Der HErr sprach zu ihm: Gehe hin, steig hinab; du und Aaron mit dir sollt herauf steigen, aber die Priester und das Volk sollen nicht herzu brechen, daß sie hinauf steigen zu dem HErrn, daß er sie nicht zerschmettere. Und Mose stieg herunter zum Volke, und sagte es ihnen.
\end{schriftzitat}

37. In diesem Texte wird uns vorgehalten, wie es zugegangen ist, da Moses das Volk hat hinzu gebracht, zu hören die Stimme GOttes vom Himmel, wie vorhin gesagt ist. Nun ist ein Stück insonderheit hier zu handeln von den Priestern, davon der Text redet, daß sie sich zuvor sollen heiligen, daß sie GOtt nicht zerschmettere: Du und Aaron, spricht der HErr, steiget herauf; die andern aber sollen unten bleiben.

38. Möchte einer fragen: Wie geht das zu, daß er der Priester gedenkt, so doch die Priester zur selben Zeit noch nicht waren? Denn es wird in dem dritten Buche Cap. 1, 5. 3, 2. ff. hernach folgen, wie GOtt den Stamm Levi und Aaron erwählet hat, daß sie Leviten und Priester sollten sein, die zu derselben Zeit noch nicht geordnet waren. Wie nennt er sie hier denn Priester? Antwort: Das muß man also verstehen: Im ersten Buch Mosis leset ihr oft, daß Priester sind gewesen, ehe das alte Testament, das ist, das Gesetz Mosis den Juden gegeben ist worden, wie auch jetzund im neuen Testamente Priester sind, nämlich, die das Wort GOttes haben, verstehen, und den andern predigen, wie Adam, Abel, Noah, Abraham, Isaak, Jakob auch gewesen sind, haben gepredigt etc.

39. [Sie] sind aber nicht geschmierte Priester gewesen, wie unsere papistischen Mönche und Pfaffen, die beschoren sind, tragen Platten, treten in Kutten und langen Röcken herein, halten Messe, singen Vigilien, lesen und murmeln ihre sieben Zeiten; bekümmern sich nicht viel um das Predigtamt, das doch allein ihr recht Amt ist; gedenken nur darauf, daß sie große Herren seien und bleiben. Daß aber die armen Seelen verderben, Hunger leiden am Worte GOttes, liegt ihnen nicht große Macht daran; [sie] essen, trinken, und leben nicht wie der gemeine Mann, vermeinen dadurch eines heiligeren und höheren Standes und Verdienstes vor GOtt zu sein denn andere Leute, die ihrem Affenspiel nicht folgen.

40. Adam aber, Noah, Enoch, Abraham, Isaak, Jakob etc. waren nicht solche Götzen und Maulaffen, sondern lebten wie andere Leute, aßen und trunken wie andere Leute, trugen Kleider an wie andere Leute. Aber darum sind sie Priester und Propheten genannt worden, daß sie das Wort GOttes predigten, das Volk führten, lehrten und ermahneten etc. Also sind solche Priester geblieben im Volke, die das Wort GOttes und solch Priesteramt führten, bis auf die Stunde, da sie auf den Berg Sinai stiegen, und eines aus den zwölf Geschlechten Israel, nämlich, das Haus Levi, von GOtt erkoren und eingesetzt ward, das äußerliche Priesterthum zu führen, 4 Mos. 3, 6., welches Priesterthum ist so viel besser denn das päbstliche, daß es einen Befehl hat von GOtt, und einen Grund aus dem göttlichen Worte, jenes aber nicht einen Buchstaben.

41. Das rechte Amt aber der Priester ist, daß sie das Wort GOttes führen; das ist vor dem alten Testamente, im, und auch darnach gewesen. Die andere Weise, die Mosis und des Pabsts Pfaffen brauchen, ist nur eine äußerliche Larve
\quellseite{519}
und Kinderspiel; doch [ist] jenes von GOtt eingesetzt zur Figur des rechten Priesterthums, dieses aber ein lauter Affenspiel.

42. Da sie nun bereit waren und sich gewaschen und geschmückt hatten, wie ihnen Mose befohlen hatte, auf den dritten Tag, kam ein groß, erschrecklich Wetter auf dem Berge Sinai. Ihr wißt, daß es allen Menschen erschrecklich ist, wenn ein groß, ungestümes Wetter ungefährlich kommt, als wolle es alles einreißen, Himmel und Erde zusammen brechen, ein Donner und Blitz im andern hergeht. Da wird einem wohl die Welt zu enge, weiß nicht, wo er sicher ist, wo er hinfliehen soll. Den Juden war es aber gar erschrecklich, weil sie in der Wüste waren, hatten keine Häuser, dazu kam das Ungewitter geschwind und unversehens, mit Donner, Blitzen, und mit solchem Feuer, wie im fünften Buch Mosis Cap. 4, 11. steht, daß das Feuer bis an [den] Himmel schlug, wie es denn zu thun pflegt, wenn solche großen Wetter kommen. Dazu war eine große Finsterniß, Wolken und\footnote{„und“ fehlt in der Wittenberger und in der Jenauer.} Dunkel, und donnerte mit zu, und fiel ein großer Platzregen. In Summa, es hatte ein Ansehen, als wollte Himmel und Erde in einen Haufen fallen, also, daß sie meinten, und sich nicht anders versahen, denn daß sie alle auf daselbenmal sollten untergehen. Er thut noch etwas Größeres hinzu, und spricht:

\subsection{Und ein Ton einer sehr starken Posaune \&c}

43. Sie hörten, daß es etwas mehr und Größeres war, denn ein gemein Wetter, dazu war ein großer, dicker Rauch und Dampf um den Berg, wie ein Rauch vom Ofen, das ist, eine schwarze Wolke, wie ein schwarzer Pechrauch. Da gedachte ein jeglicher bei ihm selber, was für Kurzweil da sei gewesen, in einem solchen grausamen Wetter, mit Donner, Blitz, Rauch und Posaunen. Da ist Lachen zu verbeißen gewesen. Das war nun die Herrlichkeit GOttes, da erzeigte er seine Gewalt, in seiner Zukunft auf dem Berge Sinai.

44. Nun, Mose führt sie hinzu, daß sie selber Gott hörten reden; aber da war kein Herz, kein Muth, keine Liebe noch Lust, daß sie zu GOtt gehen sollten. Und ist der Text, nämlich so er spricht: „Und das Volk, das im Lager war, erschrak“, wohl zu merken, denn Mose hat ihn nicht vergebens geschrieben. Denn GOtt will damit anzeigen, wie wir gegen ihm geschickt sind. Es ist ein lieber, frommer Gott, wenn er schweigt und stille hält. Er redet auch wohl durch die Zungen der Menschen, aber da ist keine Furcht noch Erschrecken, ja, man spottet ihrer wohl, köpft und ertränkt sie, die GOttes Wort reden. Wenn aber GOtt wieder also sollte reden, so wäre kein Kaiser, König, Fürst oder Herr so gewaltig, es würde ihm das Herz zappeln und entfallen, [er] dürfte nicht mucken.

45. Darum, wenn sich GOtt einmal läßt hören, kann niemandes Herz vor ihm bestehen. Die drei Jünger auf dem Berge Thabor, da sie GOttes Stimme hörten, fielen sie nieder, gleich [als] ob sie todt wären [Matth. 17, 6.]. Und die Juden, so Christum im Garten wollten haschen, wurden mit einem Worte Christi gestürzt [Joh. 18, 6.]. Unsere lieben Papisten rühmen auch viel vom freien Willen, und guten Werken und Verdiensten. Sie haben gut rühmen, weil GOtt nicht redet; sollte er aber also mit ihnen reden, wie mit den Juden, sie würden wohl zurück laufen, in einen Winkel kriechen, und ihres Ruhmens wohl vergessen.

46. Die Juden waren rein und bereit, wie gesagt ist; da sie aber hinan sollten, waren sie verzagt und erschrocken. Waren äußerlich bereit auf das allerbeste, mehr denn alle Papisten, Mönche, Pfaffen und Nonnen sich bereiten können: noch war ihnen unmöglich, daß sie vor GOtt konnten stehen bleiben. Ist das nicht eine große Schande und jämmerliche Klage, daß wir vor GOtt sollen fliehen, der uns Leib und Leben und alles gibt, das wir haben? Wo wollen wir denn hin, wenn wir nicht zu dem wollen, zu dem man laufen soll, ja, dem man nicht entfliehen kann? wie David im Psalter anzeigt, Ps. 139, 7–12.: „Wo soll ich hingehen vor deinem Geist? und wo soll ich hinfliehen vor deinem Angesicht? Führe ich gen Himmel, so bist du da. Bettete ich mir in die Hölle, siehe, so bist du auch da. Nähme ich Flügel der Morgenröthe, und bliebe am äußersten Meer, so würde mich doch deine Hand daselbst führen, und deine Rechte mich halten. Spräche ich: Finsterniß mögen mich decken, so ist die Nacht auch Licht um mich. Denn auch Finsterniß nicht finster ist bei dir, und die Nacht leuchtet wie der Tag, Finsterniß ist wie das Licht“ \&c.

47. Darum siehst du, was der Mensch für
\quellseite{520}
eine Creatur ist, daß, wo er ohne Glauben ist, ist lauter Zorn und Haß in ihm wider GOtt. Und ist ja ein wichtiger Text, der da hell anzeigt, wie der Mensch gegen GOtt geschickt ist. Er flieht ihn, ist GOtt feind, lästert ihn. Denn weil er vor ihm fleucht und fürchtet ihn als einen strengen Richter, kann er sich nichts Gutes zu ihm versehen. Darum schließen wir aus diesem Texte, daß alle Menschen GOttes Feinde sind, können ihn nicht leiden, hassen sein Gesetz. Und wiewohl sie es mit dem Munde reden, sie wollen thun und annehmen, was GOtt gebeut, wie hier die Juden, so ist es doch nur Heuchelei, das Herz erfährt es nimmer. Und wie äußerlich das Wetter braust, sie erschreckt und verzagt macht, also war es noch viel ärger in den Herzen der Juden; ihr Herz zappelte, als ob sie alle Stunde erwürget sollten werden und umkommen.

48. Der Berg Sinai zeigt an, und GOtt durch den Berg, was das Herz fühle, so ihm das Gesetz eröffnet, und von dem Gerichte GOttes gerührt wird. Man frage einen Menschen, der ohne Glauben ist und jetzt\footnote{Erlanger: Menschen ohne Glauben, der jetzt.} sterben soll, was er für Gedanken habe, und was ihm vor Augen sei? Er würde freilich antworten: Nichts denn der Tod. Denn gleichwie die Juden sahen Blitz, Donner, Rauch, Platzregen, also, daß ihr Herz stand, als sollten sie jetzund sterben, waren keinen Augenblick sicher, ihnen war kein ander Bild vor Augen denn der Tod, sahen ihren Henker vom Himmel, hielten GOtt für einen Stockmeister und Scharfrichter, der mit einer Keule hinter ihnen stände, und wollte sie an den Kopf schlagen; also ist auch ein Herz, das ihm selber gelassen ist, hält\footnote{Wittenberger und Erlanger: hat.} GOtt für einen Henker und Stockmeister.

49. Ein sterbender Mensch, der ohne Glauben ist, hat nichts anders vor Augen denn Donner, Blitzen, die Hölle, höllisch Feuer, den Zorn GOttes; fühlt nichts anders denn den Tod. Warum? GOtt redet mit ihm, und hält ihm das Gesetz vor. Da erkennt er, was er gethan hat, und das Gewissen weiß, was es auf sich hat, denn GOtt hat ihm das Gesetz eröffnet. Da ist denn keine Ruhe, sondern lauter Fliehen, Zagen und Zappeln; wie Adam im Paradeise, da er GOtt hörte reden, 1 Mos. 3, 9. 10., wäre er gerne über hundert Meilen, über alle Berge geflohen, wenn er gekannt hätte, ihm war das ganze Paradies zu enge; er konnte aber GOtt nicht entlaufen.

50. Wir wissen es und glauben es aber nicht, denn wir haben es nicht erfahren. Die es aber erfahren haben, sind unterweilen im Unglück gesteckt, in großen Nöthen Wassers, Feuers oder des Todes gewesen, die wissen davon zu sagen, was für ein Herz ein solcher Mensch gegen GOtt habe, wie es ihm gar nicht vertraut, sondern flieht vor ihm, als vor dem leidigen Teufel. Aus diesem gewaltigen Texte sollen wir schließen, daß nichts Gutes im Menschen ist.

51. Wo bleibt nun die Lehre von der Bereitung zu der Gnade, da sie sprechen: Wir sollen thun, als viel an uns ist? Sie reden von der Sache, wie ein Blinder von der Farbe, wissen nicht, wovon sie reden; denn so sie sich am besten bereiten und schicken, wenn es an das Treffen geht, so wird ihnen die Welt zu enge. Sich selber bereiten aus eigenen Kräften, kann nicht bestehen vor GOtt, hält den Stich nicht. Wenn uns der Tod unter Augen schnaubt, so fällt es alles dahin, steht wie Butter an der Sonne. Man muß uns anders lehren, wie wir GOtt entgegen gehen sollen, und vor ihm bestehen. Und (wie 1 Joh. 2, 28. geschrieben steht)\footnote{„geschrieben steht“ und gleich folgend: „gar“ fehlt in der Erlanger.} liegt gar viel daran, daß wir mögen bestehen, und nicht zu Schanden werden, wenn er kommen wird. Ich wäre ein schlechter Christ, wenn ich wollte, daß GOtt mit dem jüngsten Tage außen bliebe, so doch Christus Luc. 21, 30. die Seinen tröstet mit dem jüngsten Tage, als mit dem Sommer.

52. Nun, was helfen uns die Werke, Kappen, Platten, der geistliche Stand? Wenn ich Blut schwitzte, und mich zu Tode marterte, so fällt es doch dahin; wenn der Tod hertritt, so hält es die Furcht und Zorn GOttes nicht auf. Was ist es denn, daß ich mich lange und zu Tode martere, wenn ichs nicht gebessert werde? Es hilft dazu nichts, weder unser Wort noch Werk. Es muß allein Christus ausrichten, den ich für einen Bruder und Freund erkennen soll, und zu GOtt sprechen: HErr, ich weiß niemand, weder im Himmel noch auf Erden, zu welchem ich eine tröstliche Zuflucht möchte haben, denn zu dir, durch Christum. Ich muß mich nackt ausziehen, von allen Freunden, Werken und Verdienst
\quellseite{521}
. Herr, ich habe keine Zuflucht, denn zu deinem göttlichen Schooß, darin der Sohn sitzt. Wenn ich die Hoffnung nicht habe, so ist es verloren.

53. Diesen Text hat man unter die Bank gesteckt. Hätte man ihn hervorgezogen, hätten wir gewußt, was der Mensch vermöchte vor GOtt, und wäre die Lehre von Werken und freiem Willen nicht also eingerissen, die doch keinen Stich kann halten. Darum wage es der Teufel auf sie; und so\footnote{Statt „so“ sollte wohl „da“ gelesen werden.} sie nicht hilft, oder besteht, so es an das Treffen geht, so gebe GOtt allen Pfaffen, Mönchen, Platten, Kappen, Werken, Verdiensten, allem guten Leben eine gute Nacht. Denn es muß ja ein ander und viel höher Ding sein, das mich zu GOtt führt, denn eine lausige Kappe oder Platte. Kurzum, es muß ein solcher Muth und Herz da sein, daß den Zorn GOttes, die Sünde, Tod, Teufel und Hölle nicht fürchte; und wo es sich schon davor fürchtet und entsetzt, daß es doch endlich frei hindurch dringe. Einen solchen Muth aber vermag mir keine Kappe etc. zu geben, es muß etwas Höheres und Besseres da sein, denn aller Mönche und Nonnen Verdienst und Werk.

54. Ich habe euch vermahnt, und thue es noch einmal, daß ihr ja den Text wohl anseht, denn er ist gewaltig; und die rechten Christen wissen auch, was er vermag, wie mächtig er ist, als David, der auch darin gesteckt ist, hat es wohl erfahren.

55. Die Epistel an die Hebräer, Cap. 12, 20. 21., zeigt an aus diesem Orte des andern Buchs Mosis, daß Mose auch erschrocken sei gewesen, mitsamt den andern Juden. Denn wie die Epistel daselbst ausweist, hat er gesprochen am Berge Sinai: „Ich bin erschrocken, und zittere.“\footnote{Erlanger: Ich bin fürchtig und zittern.} Das Schrecken aber kam daher, daß ein Rauch aufging, wie von einem Ofen, und ein erschrecklicher Ton einer Posaune gehört ward etc. Summa Summarum: Es ist kein Heiliger, der nicht davon erschrocken wäre, er wäre denn ganz vollkommen und durchgeistet.

56. Das soll nun unsere Arbeit sein, daß wir dahin kommen durch den Glauben und GOttes Wort, daß wir von Tag zu Tage je stärker und fester werden, je länger je mehr Christum lernen erkennen, nicht allein mit den Ohren und der Zunge, sondern daß es auch in das Herz komme. Denn je größer die Erkenntniß ist, je kleiner die Furcht wird; und je kleiner die Erkenntniß ist, je größer die Furcht ist. Da ist verdammt all unser Wesen, mit Werken, Messen, Vigilien etc.

57. Daß hernach folgt, wie der HErr Mosen heißt hinabsteigen, daß er dem Volk zeugete etc., sind noch nicht die zehn Gebote, sondern ist eine Vorrede; und wiewohl das Volk die Sprache nicht hat verstanden, hat es dennoch die Stimme gehört; und da GOtt mit ihnen geredet hat, hat der Donner, Blitz und die Posaune stille gehalten, daß sie die Stimme konnten hören und vernehmen; das Feuer aber und die dicke Wolke hat für und für gestanden, daraus die Stimme der zehn Gebote gehört ward.

\section{Das zwanzigste Capitel}

\begin{schriftzitat}
V. 1. Und der HErr redete alle diese Worte.
\end{schriftzitat}

1. Da nun Moses dem Volke angetragen, daß ihm der HErr befohlen hatte, nämlich, daß sie untenan stehen blieben, da er die Zeichen aufgesteckt hatte etc., schwieg er stille, und GOtt redete mit dem Volke allein. Ihr könnt wohl abnehmen, wie eine helle, große Stimme es hat müssen sein, nicht eines schlechten Mannes Stimme, die ein solch groß Volk hat mögen hören. Denn es sind sechsmal hunderttausend Mann gewesen, ohne Weib, und Kinder unter zwanzig Jahren. Die haben einen großen Raum eingenommen, wohl zwo oder drei Meilen, auf das wenigste ja eine Meile; dazu ist der Berg Sinai ein sehr hoher Berg gewesen. Einer aus uns bringt seine Stimme kaum über sechstausend Menschen; wenn sie gar stark ist, über zehntausend; die aber war deutlich, erscholl über das ganze Volk,
\quellseite{522}
und mochte von jedermann wohl vernommen werden; darum muß es eine gewaltige, starke Stimme gewesen sein. Sie haben nichts gesehen denn die Wasserwollen, und das Feuer, aus welchem sie die Stimme hörten. Das war aber die Stimme, wie folgt, nämlich die zehn Gebote.

\section{Das erste Gebot}

\begin{schriftzitat}
V. 2–6. Ich bin der HErr, dein GOtt, der dich aus Egyptenland, aus dem Diensthause, geführt habe. Du sollst keine anderen Götter neben mir haben. Du sollst kein Bildniß, noch irgend ein Gleichniß machen, weder deß, das oben im Himmel, noch deß, das unten auf Erden, oder deß, das im Wasser unter der Erde ist. Bete sie nicht an, und diene ihnen nicht. Denn ich, der HErr, dein GOtt, bin ein starker Eiferer, der da heimsucht der Väter Missethat an den Kindern bis in das dritte und vierte Glied, die mich hassen. Und thue Barmherzigkeit an viel tausend, die mich lieb haben, und meine Gebote halten.
\end{schriftzitat}

\section{Das andere}

\begin{schriftzitat}
V. 7. Du sollst den Namen des HErrn, deines GOttes, nicht vergeblich führen, denn der HErr wird den nicht unschuldig halten, der seinen Namen vergeblich führet.
\end{schriftzitat}

\section{Das dritte}

\begin{schriftzitat}
V. 8–11. Gedenke des Sabbathtages, daß du ihn heiligest. Sechs Tage sollst du arbeiten, und alle deine Werke schaffen, aber am siebenten Tage ist der Sabbath des HErrn, deines GOttes, da sollst du kein Geschäft thun, noch dein Sohn, noch deine Tochter, noch dein Knecht, noch deine Magd, noch dein Vieh, noch dein Fremdling, der in deiner Stadt Thor ist. Denn in sechs Tagen hat der HErr Himmel und Erde gemacht, und das Meer, und alles, was darinnen ist, und ruhete am siebenten Tage. Darum segnete der HErr den Sabbaththag, und heiligte ihn.
\end{schriftzitat}

\section{Das vierte}

\begin{schriftzitat}
V. 12. Du sollst deinen Vater und deine Mutter ehren, auf daß du lange lebest im Lande, daß dir der HErr, dein GOtt, geben wird.
\end{schriftzitat}

\section{Das fünfte}

\begin{schriftzitat}
V. 13. Du sollst nicht tödten.
\end{schriftzitat}

\section{Das sechste}

\begin{schriftzitat}
V. 14. Du sollst nicht ehebrechen.
\end{schriftzitat}

\section{Das siebente}

\begin{schriftzitat}
V. 15. Du sollst nicht stehlen.
\end{schriftzitat}

\section{Das achte}

\begin{schriftzitat}
V. 16. Du sollst kein falsch Zeugniß geben wider deinen Nächsten.
\end{schriftzitat}

\section{Das neunte}

\begin{schriftzitat}
V. 17. Du sollst dich nicht lassen gelüsten deines Nächsten Haus.
\end{schriftzitat}

\section{Das zehnte}

Du sollst dich nicht lassen gelüsten deines Nächsten Weibes, noch seines Knechts, noch seiner Magd, noch seines Ochsen, noch seines Esels, noch alles, das dein Nächster hat.\footnote{In der Wittenberger und in der Erlanger stehen die beiden letzten Gebote in umgekehrter Reihenfolge.}

2. Das ist die Predigt, die GOtt selber thut zu dem ganzen Volke Israel in der Wolke und Feuer, auf dem Berge Sinai, mit lieblicher Stimme; wiewohl ein Engel, aus Befehl GOttes, die Stimme geführt hat, wie an einem andern Orte der Schrift steht, mit herrlicher Pracht, an der Statt GOttes. Das Volk aber hat nichts gesehen, allein die Stimme gehört etc.

3. Wir wollen die zehn Gebote kurz überlaufen. Zum ersten ist zu merken, daß uns Heiden und Christen die zehn Gebote nicht betreffen, sondern allein die Juden. Das bezeugt und zwingt der Text, so er spricht:

\begin{schriftzitat}
V. 2. Ich bin der HErr, dein GOtt, der dich aus Egyptenland, aus dem Diensthause, geführt habe.
\end{schriftzitat}

4. Das ist ja wahr und klar genug, daß wir Heiden durch GOtt nicht aus Egypten geführt sind, sondern allein das jüdische Volk Israel. Darum deutet Moses die zehn Gebote allein auf das Volk, welches durch GOtt aus Egypten ist geführt. Daß wir aber gleich auch den GOtt, den die Juden ehren, der sie aus Egypten geführt hat, erkennen, anbeten und ehren, haben wir nicht durch Mosen, oder aus den geschriebenen Gesetze, sondern aus andern Schriften, und aus dem Gesetze der Natur. Das rede ich
\quellseite{523}
abermal darum, daß ich den falschen Geistern wehre, die uns Mosen auf den Hals mit Gewalt wollen legen, ihn zu halten mit allen seinen Geboten. Das wollen wir aber lassen, und ihn mit dem allergeringsten Tüttel\footnote{In den alten Ausgaben: Titel.} nicht annehmen, denn so ferne, wo er mit dem natürlichen Gesetze übereinstimmt. Wir wollen ihn wohl lesen, wie einen andern Lehrer, frei und ungezwungen, aber für unsern Gesetzgeber wollen wir ihn nicht haben. Denn wir haben vorhin im Neuen Testamente Gesetze genug; darum wollen wir ihn nicht haben in unserm Gewissen, sondern das Christo allein rein behalten. Also ist es ja klar, daß die zehn Gebote allein den Juden gegeben sind, und nicht uns. Trotz allen Rottengeistern, daß sie mit Wahrheit anders sagen!

\subsection{Was Zeichen, dabei die Juden GOtt sollen ergreifen}

5. Nun GOtt, der HErr, wie seine göttliche Art ist, gibt den Juden ein gewiß Zeichen, dabei sie ihn ergreifen, fassen und erkennen sollen. So thut er allewege, wenn er mit uns handelt, läßt uns nicht gen Himmel gaffen; denn er will nicht, daß wir ihm dienen sollen aus eigener Wahl und Gutdünkel, wie denn die Natur und Vernunft pflegt zu thun, läßt GOttes Befehl liegen, wählt ihr einen eigenen Gottesdienst; GOtt will aber ihm gedient haben, gesagt und erkannt sein, nach seinem eigenen Wort und Befehl. Darum gibt er uns Zeichen, dabei wir ihn gewiß fassen mögen. Den Juden hat er das Zeichen gegeben, das Wort in ihren Mund gelegt, damit sie ihn sollen anrufen. Welches? Ich rufe dich an, mein GOtt und HErr, der du uns aus Egyptenland, aus dem Diensthause, geführt hast. Item, andere Worte mehr, wie hernach folgt: Der du uns geführt hast wunderlich durch das rothe Meer, durch die Wüste, gespeiset mit Himmelsbrod, aus einem harten Felsen getränket, durch den Jordan in das verheißene Land geführt etc. Das war ihr Wort, ihre Zeichen, dadurch sie GOtt erkannten.

\subsection{Der Christen Zeichen, GOtt zu ergreifen}

6. Wobei erkennen und ergreifen ihn die Christen, weil ihnen der Juden Zeichen nicht gilt? Haben sie nicht auch ein Zeichen oder Wort, GOtt damit zu ergreifen? Ja, sie haben auch ihr eigen Zeichen und Wort, nämlich das: O GOtt, ein Schöpfer Himmels und der Erde, der du deinen Sohn, Jesum Christum, für mich in die Welt gesandt hast, daß er für mich gekreuziget würde, stürbe, und am dritten Tage wieder auferstände, gen Himmel führe, daß er da sollte sitzen zu deiner Rechten, und alles in seiner Hand haben, und seinen Geist senden, daß wir sollten warten auf seine Zukunft, zu richten beide, Lebendige und Todte, und also mit ihm\footnote{„ihn“ fehlt in der Wittenberger und in der Jenaer.} erlangen das ewige Reich, unser Erbtheil, das du uns durch ihn willst geben.

7. Dazu, o HErr GOtt! hast du uns gegeben und eingesetzt die Taufe und das Sacrament des Leibes und Blutes Christi, deines Sohnes etc. Denn an diese seine Sacramente hat er uns Christen gebunden und sich uns darinnen geoffenbart; so wir ihn da ergreifen, haben wir ihn gewißlich getroffen. Das andere aber alles, das uns nicht befohlen ist, sollen wir fahren lassen, denn wir würden sonst sehr weit fehlen. Diesen Titel sollen wir also führen, wie die Juden ihren Titel führten, da sie aus Egypten geführt sind, aus dem Diensthause.

8. Wider diesen Titel hat am ersten Jerobeam, der König in Israel, gestritten, nachmals alle Juden. Derselbe Jerobeam machte güldene Kälber, setzte eines gen Bethel, das andere gen Dan, wollte also GOtt daran binden, und ihn dabei ergreifen; denn er sprach 1 Kön. 12, 28. 29.: „Siehe, da ist dein GOtt, Israel, der dich aus Egyptenlande geführet hat.“ Es war aber der Teufel. Warum? Und war doch seine Meinung, daß er dem rechten GOtt mit meinte? GOtt kann es aber nicht leiden; ja, es ist kein GOtt weder im Himmel noch auf Erden, der solchen eigenerwähleten Dienst haben wolle. Nun, dieser König malt GOtt also in seinem Herzen ab, und wählt ein eigen Zeichen aus seinem Gutdünkel, dabei er GOtt ergreifen wollte. Das will GOtt nicht haben, sondern er malt sich selber ab, und gibt den Juden ein Zeichen, und spricht: Da findest du mich, so will ich gesucht und angerufen sein: „Ich bin der HErr, dein GOtt, der dich aus Egypten, aus dem Diensthause, geführt habe“; dabei sollt du mich ergreifen, da bin ich gewißlich, da wirst du mich finden. Daß du mich aber an das Kalb, an
\quellseite{524}
den Altar, an diesen Ort bindest, irrst du; du wirst mich auch nicht finden, ja, wirst mein weit fehlen, denn ich bin nicht da.

9. Also ist es auch mit uns zugegangen; wir haben ein Zeichen, daran wir GOtt nicht fehlen können,\footnote{Wittenberger: kunden; Jenaer: können, Erlanger: kunnten.} treffen ihn gewiß. Das ist aber das Zeichen, wie oben [§ 6] gemeldet: HErr GOtt, himmlischer Vater, der du deinen Sohn für mich hast lassen Mensch werden, sterben, begraben \&c., in demselben Namen rufe ich dich an \&c. Da kann ich nicht fehlen, [ich] muß treffen, mein Gebet ist auch gewißlich erhört, wenn ich anders von Herzen also kann sagen und glauben. Denn auf die Weise hat er befohlen, ihn\footnote{In den alten Ausgaben (nach lateinischer Construction) „sich“ statt: ihn.} anzurufen, will sich also an dem Zeichen lassen finden, ergreifen und fassen, und geben, was wir bitten, bitten wir anders in Christi Namen und im Glauben.

10. Wenn ich mir aber ein besonder Zeichen nach meinem Gutdünken machte, wie Jerobeam und andere gottlose Menschen, und spräche: O GOtt, mein HErr, ich rufe dich an, daß du ansehen wollest, daß ich eine Kirche gebauet, eine Messe, Altar gestiftet, ein Meßgewand in die Kirche gegeben, eine Platte getragen, eine Kappe angezogen, meine Siebenzeit,\footnote{Siebenzeit = horae canonicae; die Namen der „sieben Gezeiten“ folgen hier vollständig: Mette \&c.} Mette, Prim, Terz, Sext, Non, Vesper und Complet mit großem Fleiß gebetet, mich mit Geißeln gepeitscht, wollen und barfuß gegangen, einen härenen Strick und Hemd getragen, auf der Bank, auf der Erde geschlafen, gefastet, gebetet, große Armuth gelitten, Jungfrauschaft, Gehorsam und Armuth gehalten habe, bin ein Mönch, Nonne oder Pfaffe gewesen: was soll er wohl dazu sagen? [Er] wird mir nicht viel darum hören, sondern denn dazu sprechen: Hab ich mich dir also abgemalt? Habe ich dir auch irgend befohlen, daß du mich durch solche Werke erkennen solltest? Ich kenne dich und deine Werke nicht, das Zeichen habe ich dir nicht gegeben, es ist des Teufels Zeichen und Bild, es ist kein GOtt weder im Himmel noch auf Erden, der also gesinnt sei, der dir gnädig wolle sein um solcher selbsterlesenen und erwählten Werke willen; ich habe dir’s nicht befohlen, ich habe nicht also wollen ergriffen werden. Welcher Teufel hat dichs geheißen? Da fällt denn die Platte, Kappe, Möncherei, Pfafferei alles zum Teufel,\footnote{Die Worte: „hat dichs ... zum Teufel“ fehlen in der Wittenberger.} kann das Urtheil GOttes nicht leiden, zerschmelzt wie der Schnee. GOtt spricht: Ich habe dir ein Zeichen und Maß gegeben, daran halte sich, das andere laß fahren, es gilt vor mir nicht; deine eigene Weise hast du selber erdacht, der Teufel hat dir es eingegeben.

11. Wir Christen sind eben bisher gewesen, wie die Alten, die erdachten ihnen eine eigene Weise GOtt zu dienen; einer nahm dies Zeichen, der andere ein anderes, wie ihm nur träumte; fuhren zu und sprachen: Es ist ein alt Ding, daß uns GOtt aus Egypten hat geführt; zu Bethel und Silo aber ist der rechte Gottesdienst, da wollen wir GOtt treffen. Ja, sie trafen ihren Gott, den Teufel. Ueber das baueten sie Kirchen und Altäre auf den Bergen \&c., meinten, sie hätten es wohl getroffen. Dawider schrieen denn die Propheten, sagten, es wäre unrecht. Da hob sich denn erst der Hader: Hinweg mit den Ketern, nur erwürgt und zu Tode geschlagen. Da mußte denn mancher Prophet das Leben lassen. Also ist es mit uns auch zugegangen: da haben unsere lieben Papisten so manches Kindlein, eines nach dem andern erdacht, davon ihnen nichts von GOtt befohlen ist. Wenn man nun dawider predigt, so werden sie toll und thöricht, schelten uns für Ketzer, Verführer, und aufrührisch, die gute Werke verbieten; da sollen wir alle erwürgt und todtgeschlagen werden.

12. Nun, die Juden sollten bei ihrem Zeichen und Befehl geblieben sein, damit sie GOtt gefasset hatte, und damit sie wiederum ihn auch sollten gefaßt haben, sollten ihnen nicht neue Zeichen und Maß erdichtet haben. Nun aber ist dasselbige Zeichen aus, [es] hat gewähret bis auf Christum. Wir Christen aber sollen nicht also beten, auch das Zeichen nicht auf uns führen; es ist nun damit aus. Die Juden möchten also beten: O HErr GOtt, der du uns aus Egypten geführt hast \&c. Wenn er nun zu mir hätte gesagt: Ich habe dich aus Egypten geführt, so treffe ich es recht, wenn ich also bete: O HErr, der du mich hast aus Egypten geführt \&c.

\quellseite{525}
13. Will ich nun nicht fehlen, so muß ich das Wort fassen, das in die ganze Welt erschollen ist, und ist das: Ich habe meinen Sohn für dich gegeben, der sein Blut für dich vergossen hat, gestorben ist, und dich erlöst, und hat dich mir versöhnt, mich dir zum Freunde und Vater gemacht.

14. So mag ich nun wohl also beten: Himmlischer Vater, der du alle Dinge geschaffen hast, der du die Kinder Israel aus Egypten [2 Mos. 12, 51.], durch das rothe Meer [Cap. 14, 22.], durch die Wüsten und durch den Jordan geführt hast, aus der Hand Pharaonis erlöst, mit Himmelsbrod gespeiset [Cap. 16, 15.], mit Wasser aus dem Felsen getränket hast etc. [Cap. 17, 6. 7.], das alles aber geht mich nicht an. Der du mit Noah große Wunder angerichtet hast [1 Mos. 7, 5. f. Cap. 8, 15. f.], geht mich auch nicht an. Der du Petrum auf dem Meere ließest gehen [Matth. 14, 29.], den Aussätzigen Befehl gabst, den Priestern sich zu erzeigen [Luc. 17, 14.], geht mich auch nicht an. Ich muß ihn anrufen und ergreifen mit dem Worte und Zeichen, das mich angeht, nämlich also: HErr, der du mich erlöst hast durch das Blut deines Sohnes Jesu Christi. Das Wort geht mich an, das durchdringet den Himmel, mit dem Worte treffe ich gewißlich GOtt, damit hat er sich angebunden. Also stimmt er dem jüdischen Volke eine sonderliche Weise, daß sie in diesem Werke ihn ergreifen sollen, daran sie ihn gewißlich getroffen haben. Darum setzt Moses in das erste Gebot: „Ich der HErr, dein GOtt, der ich dich aus Egypten geführet habe“ etc. An das Werk sollt ihr Juden gedenken, und mich dabei erkennen. Das Zeichen hat gewährt bis auf Christum, und nicht weiter.

15. Darum will ich abermal hier gewarnet haben alle Prediger. Denn ich sehe, daß vonnöthen ist, daß sie ja wohl lernen den rechten Brauch Mosis, und das Volk mit dem Mose unverworren lassen, und ihn ja nicht weiter lassen gelten, denn zu einem Exempel, und wo er ein Evangelist und Prophet ist. Wenn dich nun ein Prediger will auf Mosen bringen, so frage ihn: ob du auch unter dem Mose aus Egypten geführt seiest? Spricht er Nein, so sprich: Was geht mich denn Moses an, weil er redet mit dem Volke, das aus Egypten geführt ist? Ist es doch mit Mose aus; denn sein Amt währet nur bis auf Christi Zeit.

16. Also zwingt dieser Text gewaltig, daß die zehn Gebote auch nur allein den Juden sind gegeben, und nicht den Heiden, wie auch im dritten Gebot erzwungen wird. Denn die Heiden sind je nicht aus Egypten geführt. Wir haben ein ander Werk, und einen andern Lehrer, der nicht, wie Moses, zwingt und schreckt, sondern Gnade anbeut, tröstet, gibt und hilft und errettet, nämlich, Jesum Christum.

17. Zu Adams und Abels Zeiten ging das Opfer [1 Mos. 4, 4.], das war ihr Zeichen. Aber es geht mich nichts an. Noah hatte sein Wort und Befehl vom Regenbogen und andern Dingen [1 Mos. 9, 13. 14.], [das] geht mich nicht an. Abraham hatte die Beschneidung [1 Mos. 17, 10. f.], welche ihm gegeben war zum Zeichen seines Glaubens [Röm. 4, 11.], sie geht mich aber nicht an. Isaak und Jakob hatten ihren Befehl. Moses, Aaron, Josua hatten auch ihr Wort. Es ist aber mir nicht gesagt, daß ich die Kinder Israel soll aus Egypten, durch das rothe Meer, Wüste und den Jordan führen; und ist doch gleichwohl GOttes Wort und Befehl. Dieweil er nicht spricht: Das befehle ich dir, das und das sollst du thun, soll ich mich nicht annehmen. David hatte ein Wort, daß Christus von ihm sollte geboren werden, wie oft in der Schrift wird angezeigt [2 Sam. 7, 12. 19.] und es Paulus Röm. 1, 3. hoch aufmüßt, und Christus im Evangelio [Matth. 9, 27. 12, 23. 15, 22. 20, 30.] hin und wieder genannt wird ein Sohn Davids; folgt darum nicht, daß ich müsse Christi Vater sein nach dem Fleisch. Ist es doch GOttes Wort, wie die Schwärmer und Rottengeister plaudern. Ist wahr; es ist aber mir nicht gesagt. Moses hatte auch einen Befehl, die Hütte zu bauen [2 Mos. 26, 1. ff.] und das Priestertum einzusetzen [Cap. 28, 1. ff.], mir und dir ist es aber nicht befohlen.

18. Also durch und durch; siehe in der ganzen Schrift alle GOttes Worte und Befehle an, und deute sie nicht auf dich, so lange bis\footnote{Wittenberger und Erlanger: daß.} du gewiß seiest, daß [es] zu dir gesagt sei; so thue es denn, frage nicht, was den andern vorgehalten und geboten sei. Ja, sprichst du mit den Rottengeistern, GOtt hat es mit Mose geredet, darum muß ich es auch thun. Lieber, sprich: Es liegt nicht Macht daran, daß Wort streckt sich nicht weiter, denn auf den, dem es befohlen ist. GOtt heißt
\quellseite{526}
den Fisch im Wasser schwimmen, wie im ersten Buch Mosis steht [Cap. 1, 20. ff.], die Vögel in der Luft fliegen, das Gewürm auf der Erde kriechen, die Sonne scheinen; da ist auch GOttes Wort. Willst du darum ein Fisch werden, und im Wasser wohnen? in der Luft schweben, wie ein Vogel? Willst [du] darum zur Sonne, Mond und Sternen werden? Siehe darauf, wie es dir werde ansehen. Will man denn nicht sehen, daß Mosis einem sonderlichen Volke gegeben ist? Darum zeuch ihn nicht auf die ganze Welt, sondern auf sein Volk.

19. Also haben wir [die] Ursache, warum GOtt diesen Namen führt: „Ich bin dein GOtt, der dich aus Egypten geführet habe“; denn er nimmt sich ihrer an, als seines sonderlichen Volks, das er ihm zum Eigenthum erwählt hat. Es gilt aber jetzund nicht mehr, denn ich rufe jetzund GOtt nicht also an, der mich aus Egypten geführet habe. Damit man den Schwärmergeistern antworte, und das Maul stopfe, die Mosen mit Gewalt auf uns wollen bringen. Man lasse sie nur frisch mit ihrem Mose fahren, und Mosen einen Gesetzgeber bleiben des Volks Israel, und lade ihn den Heiden und Christen nicht auf den Hals. Denn im neuen Testamente hat er ein Ende, und gilt nichts mehr mit seinen Gesetzen, er muß sich vor Christo verkriechen.

20. Daß aber wir Heiden ein Gesetz haben, das lehrt uns unser eigen Gewissen und Vernunft; wie auch St. Paulus Röm. 1, 19–21. spricht, daß die Heiden auch eine Erkenntniß von GOtt haben; denn „GOtt habe ihnen das offenbaret, daß sie GOttes unsichtbares Wesen, das ist, seine ewige Kraft und Gottheit sehen, so man das wahrnimmt bei den Werken, von der Schöpfung der Welt an; aber sie haben ihn nicht als einen Gott gepriesen“ etc. Mit welchen Worten St. Paul anzeigt, daß alle Heiden Erkenntniß von GOtt haben, nämlich, daß er alle Dinge geschaffen habe, alle Dinge gebe, alles ernähre, erhalte; darum dringet sie ihr eigen Gewissen, daß sie GOtt die Ehre geben sollen, und ihm danken für alle Wohlthaten.

21. Derhalben, wenn gleich Mosis das Gesetz nie geschrieben hätte, so haben doch alle Menschen das Gesetz von Natur in ihren Herzen geschrieben. GOtt aber hat den Juden auch ein geschrieben Gesetz, das ist die zehn Gebote, gegeben, zum Ueberfluß; welche auch nichts anders sind, denn das Gesetz der Natur, das uns natürlich in das Herz geschrieben ist. Was nun Mosis geschrieben hat in den zehn Geboten, das fühlen wir natürlich in unserm Gewissen. Denn so die Heiden, spricht der Apostel [Röm. 2, 14. 15.], „die das Gesetz (das ist, Mosis geschriebenes Gesetz) nicht haben, und doch von Natur thun des Gesetzes Inhalt, dieselbigen, dieweil sie das Gesetz nicht haben, sind sie selber ihnen ein Gesetz, damit, daß sie beweisen, des Gesetzes Werk sei geschrieben in ihren Herzen, sintemal ihr Gewissen sie bezeuget, dazu auch die Gedanken, die sich unter einander verklagen, oder entschuldigen“ etc.

22. Nun wiewohl beide, Juden und Heiden, ein Gesetz haben, dennoch haben sie gleichwohl GOttes gefehlt. Denn es kann dem Gesetze nicht genug geschehen ohne GOttes Geist und Glauben, weil GOtt alles unter den Unglauben beschlossen hat, auf daß er sich aller erbarme [Röm. 11, 32.], und alle Welt in Adam gesündigt hat [Röm. 5, 12.].

23. Nun hat GOtt den Juden die Ehre und Vortheil gethan, daß er ihnen die zehn Gebote mündlich und schriftlich gefaßt hat, zum Ueberfluß, um deswillen, daß er von den Juden wollte Mensch werden. Doch sollen wir Heiden, welchen GOtt kein schriftlich Gesetz gegeben hat, nichtsdestoweniger ihn ehren, preisen und ihm danken. Denn er ist gleichsowohl unser GOtt, als der Juden GOtt [Röm. 3, 29.], wie wir nachmals weiter davon sagen werden.

\subsection{Ich bin der HErr, dein Gott}

24. Diese Worte allein, abgesondert von den andern, nämlich von denen: „Der dich geführet hat aus Egyptenlande, aus dem Diensthause“ etc., die müssen wir wohl ausstreichen, und fleißig Acht darauf haben. Denn sie betreffen uns alle, die ganze Welt insgemein, und einen jeglichen Menschen insonderheit, Juden und Heiden; nicht darum, daß es Mosis geschrieben hat, sondern daß GOtt alle Menschen geschaffen hat, erhält, regieret etc. Wie auch Paulus Röm. 3, 29. 30. spricht: „Ist GOtt allein der Juden GOtt? Ist er nicht auch der Heiden GOtt? Ja freilich auch der Heiden GOtt; sintemal es ist ein GOtt, der da rechtfertigt die Beschneidung aus dem Glauben, und die Vorhaut durch den Glauben“.

25. Diese Worte aber: „Ich bin der HErr, dein GOtt“, redet GOtt zu dem ganzen Volke
\quellseite{527}
nicht anders, denn als sei nur Ein Mensch da. Er spricht nicht: Ich bin der HErr, euer GOtt, sondern: „Dein GOtt.“ Das Wörtlein „dein“, das siehe wohl an, denn es liegt die größte Macht an dem Wörtlein. So sagt er nun: „Ich bin der HErr, dein GOtt“, als wollte er sprechen: Ich will mich euer aller, und eines jeglichen insonderheit annehmen, als wäre nur Einer allein, und sonst keiner auf Erden. Und das thue ich darum, daß ihr meine Worte desto fleißiger sollet fassen, wahrnehmen, und behalten, daß ja keiner möge sagen: GOtt hat wohl das Gesetz gegeben, hat sich hoch und viel erboten, er wolle den Menschen gnädig und ihr GOtt sein, bei ihnen thun, wie ein Vater bei seinem Kinde; es geht aber vielleicht mich nicht an, wer weiß, ob er mich auch meint, es geht allein den Haufen an.

26. Denn das ist die Art der Menschen, daß sie das Wort GOttes in den Wind schlagen, nehmen sich deß nicht an, gaffen nur auf andere Leute, gedenken nicht, daß [es] ihnen solle gelten. Dem will nun hier GOtt vorkommen und ihnen wehren, daß sie nicht das Maul aufsperren, und andern zusehen, was sie thun, sondern will sie gewiß machen, daß er einen jeglichen insonderheit meine, da er spricht: „Ich bin der HErr, dein GOtt“; dich, dich meine ich, und keinen andern; als spräche er:\footnote{In den alten Ausgaben: als er spreche.} Siehe nicht, was andere thun, sondern höre du, was ich dir sage, und siehe, wie du es annimmest und glaubest; siehe nicht die andern an, ich will mit dir handeln, ich nehme mich deiner an, und du wiederum nimm dich meiner an.

27. Darum habe ich oft gesagt, daß einer, der selig will werden, soll also gesinnet sein, als sei kein Mensch sonst auf Erden, denn er allein, und daß aller Trost und Zusagung GOttes hin und wieder in der heiligen Schrift ihn allein angehe, sei auch um seinetwillen allein geschrieben, daß ihn ja der Teufel nicht\footnote{„nicht“ fehlt in der Wittenberger.} irre mache, wenn er sterben soll, und ihm die Augen aufsperre, und viel tausend Menschen weise, die alle gelebt und gethan haben, wie er, und werden dennoch verdammt, und spreche zu ihm: Was willst du dich nun vermessen, selig zu werden, weil du gewiß verloren bist? Willst du besser sein, denn die andern, die nicht anders gethan haben, denn du, und sind dennoch gleichwohl zum Teufel gefahren? Meinst du, du wollest ihm entlaufen? Also kann der Teufel einen in Verzweiflung führen, gleich [als] ob kein GOtt sei, der sich seiner werde annehmen, ihm helfen, und aus aller Noth ihn werde erretten.

28. Darum soll man wohl das Wort fassen, daß ein GOtt sei, der da helfen kann und helfen will, wie denn natürlich dasselbe Erkenntniß von GOtt in die Herzen geschrieben ist aller Menschen, wie Paulus spricht Röm. 2, 15., nämlich ein solches, daß GOtt, wie hier der Text anzeigt, dein GOtt sei, habe dich geschaffen etc. Wenn es aber zur Anfechtung kommt, so ist [es] schwer zu bestehen, GOtt dann recht erkennen, ihm glauben und vertrauen, also, daß auch die Frommen, die den Heiligen Geist, über das natürliche Erkenntniß von GOtt, haben, gar kaum hindurch brechen; so ein trefflicher Puff ist es, wenn uns der Teufel also mitführt, wie gesagt, in der Anfechtung, daß er auch den Frommen Arbeit macht, und sie, wenn es ihm verhängt würde durch GOtt, stürzt. Wenn nun die so hart angefochten werden, die in GOttes Schutz sind, wie wollen die bestehen, die allein das natürliche Erkenntniß von GOtt haben, das doch durch die Sünde ganz dunkel und verblichen ist? Der Teufel führt sie, wo er nur hin will, denn sie achten GOttes Wort nicht, ist ihnen ein Traum; der Teufel hat ihnen ihr Herz besessen und verblendet.

29. Also ist wohl zu merken, was GOtt hier zum ersten spricht: „Ich bin der HErr, dein GOtt.“ „Dein, dein GOtt“, der sich eines jeglichen insonderheit annimmt, mehr denn ein Vater seines eigenen Kindes. Aber wenig fassen, wenig glauben es. GOtt ist viel zu groß, gedenkt der Unglaube, daß er auf mich sollte sehen; sollte sich GOtt, die hohe Majestät, also herunter lassen, daß er auf mich armen Madensack sollte Achtung haben? ja, er läßt es wohl, er sitzt droben im Himmel, läßt ihm die\footnote{Erlanger: drei.} Engel dienen; was bin ich gegen GOtt? Eine arme Wasserblase, die von sich\footnote{In den alten Ausgaben: ihr.} selbst zerbricht. Der Glaube aber zweifelt nicht daran, daß GOtt, der alle Dinge geschaffen hat, Himmel und Erde und alles, das darinnen ist, sich\footnote{Erlanger: er sich.} unser annehme. Denn da steht das Wort:

\quellseite{528}
\subsection{Ich bin der HErr, dein GOtt}

30. So ist er nun mein und dein GOtt, daß er sich insonderheit eines jeglichen annimmt: nährt mich, hilft, errettet,\footnote{Wittenberger: redt; Erlanger: rett.} und sorgt für mich, gibt mir alles, was mir noth ist an Leibe und Seele, hat mit mir zu schaffen, wie eine Mutter mit ihrem Kinde, handelt und gebaret\footnote{Erlanger: beredt. — gebaren = handeln, verfahren.} mit mir, als sei kein Mensch sonst auf Erden, denn ich. Das hat St. Augustinus wohl gesehen und verstanden, da er spricht: HErr GOtt, du regierest alles wunderlich, gehst mit einem jeglichen um, als hättest du sonst mit niemand zu schaffen, denn mit ihm; bringst einen jeglichen hindurch sein Lebenlang. Einer wird geboren in dieser Stadt arm, zeucht in eine andere Stadt, da wird er überschwänglich reich. Ein anderer hat genug, zeucht aus in ein ander Land, da wird er zu einem Bettler. Einem gibt er, dem andern nimmt er, und partitret es so seltsam, daß keiner sprechen kann, daß er sein Leben habe hinausgeführt nach seinen Aufschlägen und Gedanken; es geht alleweg anders, denn wir meinen.

31. Es sind viel Teufel um uns, die uns alle Stunden wohl könnten tödten; aber GOtt ist da, der wehrt, daß uns nichts Böses widerfahre, sonst könnten uns alle Fürsten und Herren der ganzen Welt nicht beschützen vor des Teufels Gewalt. Denn er ist ein Herr und Gott der Welt, und hat wohl tausenderlei Weise, dadurch er uns möchte schaden, durch Feuer, Wasser, Luft, Schwert etc. Aber GOtt ist stärker denn er, wehrt ihm, daß er mir kein Bein breche, kein Auge verderbe, kein Leid thue, auch eines Haares breit. Der Teufel wäre sonst zu gewaltig, gäbe mir bald eine Schlappe; da läge ich denn, es könnte mich aller Harnisch und Geschoß vor ihm nicht einen Augenblick erretten, wenn GOtt nicht über mir hielte. Und wo [es] GOtt verhängt, frißt er den mit Feuer, den andern mit Wasser; denn thut er sonst, denn so.

32. Aber GOtt erhält uns vor ihm, und sorgt tausenderlei Weise für uns, gebraucht der Eltern auch wohl dazu, daß mit Sorgen die Mutter darauf sehe, daß das Kind keinen Mangel leide, wacht, schläft desto weniger, daß der Teufel das Kind nicht umbringe.\footnote{So die Wittenberger und die Jenaer. Erlanger: „Daß aber sie das Kind nicht umbringe.“ Diese Worte müssen sich in der Originalausgabe finden, denn die alte Ausgabe Walchs hat beide Lesarten neben einander.} Das ist eine geringe Sorge, gegen der Sorge, die GOtt auf uns hat; [er] sorgt für das Kind mehr, denn die Mutter selbst. Aber niemand sieht solches, niemand glaubt es auch; und ist doch gewiß wahr, daß GOtt also für uns sorgt und dem Teufel wehrt und allen Creaturen, damit uns der Teufel könnte umbringen. Der nun so viel Sorge auf einen Menschen legt, der sorgt nicht weniger für sie alle insonderheit, so wunderlich, daß einer auch gedenkt, er habe mit ihm allein zu schaffen. Das werden wir sehen im ewigen Leben, hier sehen wir es allein im Glauben. Aber wenig glauben es; und ist doch gewiß wahr. Daher berühmt sich GOtt also:

Sorge, gegen der Sorge, die GOtt auf uns hat; [er] sorgt für das Kind mehr, denn die Mutter selbst. Aber niemand sieht solches, niemand glaubt es auch; und ist doch gewiß wahr, daß GOtt also für uns sorgt und dem Teufel wehrt und allen Creaturen, damit uns der Teufel könnte umbringen. Der nun so viel Sorge auf einen Menschen legt, der sorgt nicht weniger für sie alle insonderheit, so wunderlich, daß einer auch gedenkt, er habe mit ihm allein zu schaffen. Das werden wir sehen im ewigen Leben, hier sehen wir es allein im Glauben. Aber wenig glauben es; und ist doch gewiß wahr. Daher berühmt sich GOtt also:

\subsection{Ich bin der HErr, dein GOtt}

33. Wolte GOtt, wir könnten es also fassen, wie er es redet und täglich beweist. Ich sorge für dich, wenn du arbeitest, schläffst etc., daß dir niemand Schaden thue, so lange, bis die Stunde kommt deines Todes. Also ist das Wörtlein: „Ich bin der HErr, dein GOtt“, nicht allein auf die Juden zu deuten, sondern auf alle Menschen in der Welt, denn er sorgt für sie alle, allein ausgeschlossen, daß sie, wie die Juden, nicht sind aus Egyptenland geführt. Man spricht: GOtt ist mit im Schiffe. Es ist wahr und fein gesagt. Ist er nicht da, so geht es unter, und alle Sorge ist verloren, wenn er nicht sorgt, hütet und wacht, Psalm 127, 1.: „Wo der HErr die Stadt nicht behütet, so wachet der Wächter umsonst.“

\begin{schriftzitat}
V. 3. 4. Du sollst keine anderen Götter neben mir haben. Du sollst kein Bildniß, noch irgend ein Gleichniß machen, weder deß, daß oben im Himmel, noch deß, daß unten auf Erden, oder deß, daß im Wasser unter der Erde ist. Bete sie nicht an, und diene ihnen nicht.
\end{schriftzitat}

34. Das ist das andere Stück im ersten Gesetz, darinnen GOtt uns wehrt, daß wir nicht fremde Götter haben sollen. [Er] deutet selber im Texte, wer die fremden Götter sind, nämlich, daß sie kein Bild sollten haben, weder himmlische noch irdische etc. Das ist, ihr sollt nicht malen die Sonne, den Mond, die Sterne, kein Bild eines Menschen, Thieres oder Fisches.

\subsection[Von den Bildern]{Von den Bildern.\footnote{So die Jenaer. Wittenberger und Erlanger: Das ist der Titel von den Bildern.}}

35. Diesen Text haben die Schwärmergeister geführt, und haben ihn auf uns treiben wollen.

\quellseite{529}
Darum wollen wir ihn handeln, erstlich, auf die Streitweise, nachmals auf die einfältige Weise. Unsere Rottengeister, Meister Klügling, die die Schrift gar gefressen haben, sprechen: Hörst du GOttes Wort, das dir sagt: Du sollst nicht fremde Götter haben? schreien also mit diesem Scheine, daß sie Gottes Wort vormenden, die Menschen; treiben sie mit Gewalt dahin, daß man die Bilder solle nicht leiden, weil es GOtt verboten habe; meinen, es sei köstlich Ding, Götzen stürmen. Was wollen wir aber dazu sagen? Es ist ja GOttes Wort, da können wir nicht klein zu sagen.

36. Lieben Christen, ihr habt also gehört, wenn sie mit ihrem Mose herein kommen, wollen euer Gewissen mit seinen Gesetzen binden, so sprechet zu ihnen: Lieber Herr, setz die Brille auf die Nase, und sehet den Text recht an. Wir wissen wohl, daß man GOtt gehorsam soll sein in dem, das er sagt, und daß wir GOtt so wohl angehören, als die Juden. Man muß aber einen Unterschied machen zwischen dem Worte GOttes, und Worte GOttes. Darauf soll ich Achtung haben, wenn GOtt etwas redet, ob dasselbe mich betreffe. Darum, lieber Geselle, willst du mich mit GOttes Worte zwingen, so sage mir einen Text, der mich angeht, sonst kehre ich mich nichts daran, daß du mir viel aus Mose sagst. Denn Moses mit seinem Worte ist uns nicht gesandt; und ob Moses schon nicht gekommen wäre, so hätten wir dennoch gleichwohl dieses natürliche Erkenntniß durch GOtt in unsere Herzen geschrieben gehabt, daß ein GOtt ist, der alle Dinge mache und erhalte. Denn auch die Heiden GOtt angebetet haben, ohne Mosis Lehre, wiewohl sie GOttes, gleich wie auch die Juden, gefehlt haben.

37. Darum kannst du bald also antworten: Lieber Schwärmer, Moses hin, Moses her. Willst du, daß ich dich höre, so sage mir ein Wort, das mich angeht, oder ich halte dich für einen Verführer und Teufels Apostel, denn du predigst, das andern, nicht dir, befohlen ist. Sollte ich alle Worte GOttes annehmen und halten, so müßte ich auch einen Kasten bauen, wie Noah; denn GOttes Wort hat ihm befohlen, daß er einen Kasten bauete [1 Mos. 6, 14.]. Nun ist ja da\footnote{Wittenberger und Jenaer: das.} Gottes Wort; so gehe hin, und thue wie Noah, baue einen Kasten. Item, Christus [Matth. 17, 27.] hieß Petrum an das Meer gehen, und den Angel einwerfen, und einen Fisch fahen, der am ersten heraufführe, in welches Mund er finden würde einen halben Gulden, den sollte er zu Schoß geben für ihn und sich. Da ist auch GOttes Wort. Aber, lieber Schwärmer, gehe hin, und thue auch, wie Petrus that; laß sehen, wie dir’s wird anstehen. Solche blinden Köpfe sind unsere Rottengeister, plumpen hinein in die Schrift, wie ein Bauer in die Stiefel. Wie käme ich dazu, daß sie mir alle GOttes Wort wollen auflegen? Sie beweisen dorhin, daß Gott die Worte wolle mir gepredigt haben.

[Matth. 17, 27.] hieß Petrum an das Meer gehen, und den Angel einwerfen, und einen Fisch fahen, der am ersten heraufführe, in welches Mund er finden würde einen halben Gulden, den sollte er zu Schoß geben für ihn und sich. Da ist auch GOttes Wort. Aber, lieber Schwärmer, gehe hin, und thue auch, wie Petrus that; laß sehen, wie dir’s wird anstehen. Solche blinden Köpfe sind unsere Rottengeister, plumpen hinein in die Schrift, wie ein Bauer in die Stiefel. Wie käme ich dazu, daß sie mir alle GOttes Wort wollen auflegen? Sie beweisen dorhin, daß Gott die Worte wolle mir gepredigt haben.

38. Darum müssen GOttes Worte den Zusatz haben, daß ich wisse, zu wem sie gesagt sind. Die Engel haben auch GOttes Wort; was geht es aber uns an, weil es zu uns nicht gesagt ist? Darum halten wir uns des Worts, das zu uns gesagt ist. Christus trieb Petrum zurück; da er fragte von Johanne: „Was soll dieser thun?“ antwortete Christus: „Was geht es dich an, folge du mir“ [Joh. 21, 21. 22.]. Ich habe es nicht ihm gesagt, sondern dir, dir; du, du folge; er wird seinen Bescheid auch kriegen. Ich sage es nicht ihm, sondern dir. Darum habe ich gesagt, daß GOtt einen jeglichen insonderheit angreift, als habe er sonst mit keinem zu thun, denn mit ihm. Wenn er nun dir ein Wort auflegt, nimm es an, und halte dich deines Befehls, und laß einen andern seines Befehls warten. Es sind heillose Tropfen und rechte Säue, wollen große Doctores sein, schreiben große Bücher, wissen doch keinen Unterschied des Wortes GOttes. Abraham hat von GOtt empfangen die Beschneidung zu einem Zeichen seines Glaubens, wie Paulus auch anzeigt Röm. 4, 11. Da ist GOttes Wort. Es ist aber aus mit der Beschneidung Abrahä, es steht allein da zum Exempel des Glaubens, bindet aber und zwingt mich nicht.

39. Also sage ich hier, daß das Bilderstürmen und Umreißen der Götzen nicht mag erzwungen werden aus diesem Texte. Denn er\footnote{Jenaer: es.} ist den Juden allein gesagt, und nicht uns. Weise mir einen Text, damit mir GOtt verboten hat die Bilder; nicht, daß ich den Bildern hold sei, sondern daß wir gewiß wissen sollen, worauf unser Glaube gegründet sei, daß wir nicht auf den Sand bauen, und unsern Widersachern können
\quellseite{530}
antworten. Denn ein Prediger, ja, auch ein jeglicher Christ, soll und muß seiner Lehre gewiß sein; nicht auf einen Wahn bauen, oder mit Menschendünkeln umgehen, sondern der Sache gar gewiß sein, daß [es] also sei, und nicht anders [Col. 2, 2.], daß Paulus Plerophorian nennt; auf daß er stehen konnte in aller Anfechtung, und dem Teufel und allen seinen Engeln, ja, GOtt selber ohne alles Wanken antworten.

40. Darum müssen sie einen andern Grund aufbringen, damit sie beweisen, man müsse Bilder stürmen. Denn wenn man sie anfechten würde, und hätten keinen andern Grund, so müßten sie zurückfliehen, könnten nicht bestehen, und würden nichts ausrichten. Darum sage ich, so muß man der Sache gewiß sein, wenn man andere\footnote{Erlanger: die Andern.} lehren will mit GOttes Wort. Also lehrt auch Petrus [1. Ep. 4, 11.], daß niemand reden soll, denn er rede es als GOttes Wort, das GOtt geheißen habe und ihm gefalle; und „so jemand ein Amt hat, daß er es thue als aus dem Vermögen, das GOtt darreicht“ [Röm. 12, 7.]. Sie aber können nichts aufbringen, daß es GOtt befohlen habe, Bilder stürmen und umreißen.

41. Zum andern sage ich, daß man die Bilder soll also abreißen und ausrotten, daß man die Herzen davon reiße und abwende. Denn was hindert mich ein Bild, wenn mein Herz nicht daran hängt? Dann aber hängt das Herz nicht daran, wenn ich nicht an die Bilder glaube, mich nicht darauf verlasse, und sie nicht sonderlich anrufe, als wollte ich mit den Bildern GOtt eine große Ehre und Dienst erzeigen, wie bisher geschehen. Denn wir haben bisher in unserer Frauen, St. Anna, Crucifix\footnote{Erlanger: St. Annen Crucifix.} und dergleichen Bilder gemacht, und die Meinung dazu gehabt, daß [sie] besser wären, denn ander Holz und Steine; ja, daß wir daran GOtt einen großen Gefallen thäten, wenn wir sie ehrten; haben also eine Zuversicht dazu gehabt. Da brachten sie uns denn nicht allein um das Geld, sondern auch um die Seele. Nun muß man solchen Bildern nicht Arm und Bein brechen, sie zerschlagen, denn das Herz bliebe gleichwohl unrein, sondern man muß das Volk mit dem Worte dahin bringen, daß sie keine Zuversicht haben zum Bildern, als könnten sie ihnen helfen, oder als wollten sie Gott einen besondern Dienst damit thun; denn das Herz muß wissen, daß ihm nichts frommt noch hilft, denn GOttes Gnade und Güte allein.

42. Daß die Bilderstürmer aber solches lehrten, lassen sie wohl anstehen, fahren viel lieber zu, und reizen das Volk, daß sie die Bilder stürmen. Da bleibt denn das Herz voll, voll Abgötterei, meint nicht anders, denn es thue wohl daran, und GOtt einen Gefallen, daß [es] die Bilder stürme, und fährt zu, richtet die andern, die es nicht thun; so doch kein Wort noch Befehl GOttes da ist. Wo aber das Volk unterwiesen würde, daß vor GOtt nichts helfe denn seine Gnade und Barmherzigkeit, so würden die Bilder von ihnen selber wohl fallen und in Verachtung kommen. Denn sie würden gedenken: Soll es denn kein gut Werk sein, Bilder machen, so mache der Teufel Bilder und gemalte Tafeln; ich will nun fortan mein Geld wohl behalten, oder besser anlegen. Aber die Rottegeister müssen etwas Sonderliches anrichten, sonst würde man nichts von ihnen halten.

43. Die Juden haben zwar ein Gebot, daß sie nicht sollen Bilder haben; aber das Gebot haben sie zu enge gespannt. Denn GOtt verbeut die Bilder, die man aufrichtet, anbetet, und an GOttes Statt setzt. Denn es sind zweierlei Bilder. Darum macht er einen Unterschied, und gibt eine Regel, welche Bilder verboten sind, nämlich, die man aufrichtet, als wären sie GOttes Bilder; wie denn der Text gewaltiglich schließt. Darum ist denen hier das Maul gestopft, die da\footnote{„da“ fehlt in der Erlanger.} sagen: Den Juden sind alle Bilder verboten.

44. Die Juden sind auch zu abergläubig, daß sie diesen Text zu strenge ziehen und machen; wie unsere tollen Sophisten mit dem Sonntage und Feiertage fahren, die bald eine Sünde daraus machen, wenn man am Sonntage Kraut feil hat, oder sonst etwas Geringeres thut. Wenn sie es so enge wollen spannen, so wäre es besser, daß man den ganzen Feiertag über schliefe; denn sonst kann man nicht ohne Arbeit sein, wie es die Juden und Sophisten deuten. Denn es ist ja eine Arbeit, einen Rock anziehen, Schuhe ausziehen, gehen, stehen, aufstehen, essen und trinken. Wenn man aber die Schrift und das Wort GOttes also will auslegen, was will daraus werden? Das ist aber die Meinung
\quellseite{531}
mit der Feier, daß sie kein Werk thun sollten, dadurch GOttes Werk verhindert würde.

45. Also auch hier wird kein Bild verboten, denn die, dadurch der Gottesdienst wird verhindert. So wird nun hier kein ander Bild verboten, denn GOttes Bild, das man anbetet.

46. Man wehrt den Kindern, daß sie nicht sollen auf die Bank steigen, und auf den Tischen sitzen, daß sie nicht herab fallen; item, daß sie nicht zum Wasser gehen, daß sie nicht ersaufen; man läßt ihnen nicht Brodmesser in den Händen, daß sie sich nicht stechen; also wehrt man den Kindern, daß doch die Natur nicht verderbet. Denn weil die Kinder unverständig und schwach sind, möchten sie Schaden nehmen, wo man ihnen nicht wehrte. Also hat auch GOtt das grobe jüdische Volk geführt mit solchen Geboten, wehrt ihnen die äußerlichen Bilder, daß sie ihrer nicht mißbrauchten, nicht in Abgötterei fielen. Welche aber verständig sind, und voll des Heiligen Geistes, dürfen solcher Gebote nicht.

47. Wenn man mir verbieten wollte, ich sollte kein Messer in die Hand nehmen, Brod damit zu schneiden, daß ich mir nicht Schaden thäte, oder sollte über kein Wasser gehen, denn gestern wäre ein Kind im Wasser ersoffen, oder auf keine Bank steigen, daß ich nicht herab fiele: so wäre es ein närrisch, lächerlich Gebot; [ich] möchte sprechen: Du Narr, siehst du mich für ein Kind an, soll ich mich erst wiegen lassen? Also thun unsere Schwärmergeister auch, geben solch närrisch Kinderspiel vor, und wollen noch [für] große Lehrer gehalten sein; aber sie mögen noch wohl eine Weile zur Schule gehen. Moses ist ein Zuchtmeister der Juden gewesen, wie Paulus [Gal. 3, 24.] spricht, welches ein grobes, fleischliches Volk war, dem mußte man äußerlich Gebot geben von Bildern, daß sie sich nicht daran ärgerten. Und es könnte noch heutiges Tages\footnote{In den alten Ausgaben: „heutes tages“. Vgl. 1 Sam. 17, 46., wo sich in der Weimarschen Bibel: „heute Tages“ findet.} geschehen, daß man solche Gebote dem groben Volke gebe; aber wir Christen, die wir GOttes Wort haben, dürfen solches Gaukelswerks nicht, wir gehören nicht in Mosis Schule, wir haben einen bessern Meister.

48. Darum den Juden, als einem groben, tölpischen Volke, waren die Bilder allein verboten, daran sie GOtt wollten eine Ehre thun. Doch, die Juden sind dennoch gleichwohl so närrisch nicht, daß sie die guten Gulden und dicke Groschen, deren sie viel haben, wegwürfen, ob gleich Marien oder St. Johannis Bild darauf geschlagen ist, so sie doch, wenn ihnen alle Bilder wären verboten, müßten auch diese meiden. Aber es ist lauter Narrentheidung. Christus ist wohl so klug und gelehrt als die Schwärmergeister; dennoch fürchtet er sich nicht, daß er Sünde daran thäte, da er den Zinsgroschen angriff, darauf des Kaisers, der ein Heide war, Bild stund; da er auch fragte\footnote{In der Originalausgabe und in der Wittenberger: da er auch gefragt ward.} [Matth. 22, 19.20.], wes das Bild und die Ueberschrift wäre, sprach er nicht zum Juden: Pfui, daß euch dies und\footnote{Erlanger: oder.} jenes angehe, warum weiset ihr mir ein solch Bild, das euch GOtt verboten hat? ihr sollt es beileibe nicht angreifen. Darum greift man es an der Wand, daß nicht alle Bilder verboten sind. Wollten wir aber den Schwärmern folgen, so müßte man kein Geld münzen, in kein Wasser sehen, in keinen Spiegel sehen, und alles, was nur ein Bild hat, hinweg legen; ja, man müßte den Leuten die Augen ausstechen, wenn man gar kein Bild sollte haben, denn man sieht Bilder auf dem Gelde, im Wasser, im Spiegel. Sind nun das nicht feine Prediger? Darum siehst du klar, daß GOtt nicht allerlei Bilder meint.

49. So antworte nun den Schwärmern: Lieber, was hält das erste Gebot in sich? Was ist seine Meinung? Es lehrt ja nicht von äußerlichen Dingen, von Ackerpflügen, von Schuhmachen, sondern wie man GOtt soll ehren und ihm dienen, innerlich und äußerlich. Willst du nun wohl lehren, und die Bibel recht auslegen, so mußt du die rechte Meinung treffen. Nun, das erste Gebot lehrt, daß ich einen GOtt soll haben, demselben allein dienen und ihn ehren; auf den Sinn und Verstand gehen alle Worte in diesem ersten Gebote. Darum werden hier die Bilder verboten, die man braucht wider die Meinung und Verstand dieses Gebots, nämlich, daß man nicht Zuversicht habe zu Bildern, sondern allein zu GOtt sich alles Gutes versehe, und meide alles das, was uns an dieser Zuversicht hindert.

50. So werden nun durch diesen Text: „Du sollst keine anderen Götter neben mir haben“, allein die göttlichen\footnote{göttlichen = die Gott darstellen.} Bilder verboten; dazu den
\quellseite{532}
Juden, und nicht uns. Denn dieses wird allein den Juden gesagt. Die Worte sind wohl zu merken. Er spricht nicht: Du sollst kein ander Bild haben vor dir, sondern spricht: „Du sollst keine anderen Götter haben neben mir“, oder vor mir. Was heißt Götter haben? Hölzerne, steinerne und silberne Bilder, die da Götter sind, wie folgt. Denn die Menschen werden dadurch verrückt, daß sie Abgöttische werden; und wenn sie gleich Holz, Stein, Silber und Gold nicht anbeten, so haben sie doch das Vertrauen, daß GOtt wohlgefalle, ihm zu Ehren Bilder aufzurichten; welches denn wider die Meinung des ersten Gebots ist: „Ich bin der HErr, dein GOtt, du sollst nicht“ etc.

51. Darum habe Acht auf die Meinung dieses Gebots: Ich will, spricht er, dein GOtt sein, will dich selig machen. Ich will dir helfen, und das aus lauter Gnaden; [du] darfst mir es nicht abdienen, keinen Gottesdienst aufrichten aus eigenem Dünkel; du sollst den Ruhm vor mir nicht haben, daß du mir durch dein Verdienst etwas abkaufest, ich will deinen Dienst nicht ansehen. Soll ich dein GOtt sein, so gebe ich dir alles umsonst, Leib und Leben, Weib und Kind, Acker, Wiesen, Reichthum, Ehre und Gut, Vergebung der Sünden und das ewige Leben. Welches sind göttliche Werke und Gaben, die ich dir gebe; und du kannst mir nichts dafür geben noch thun, denn allein mir danken, mich preisen und loben; nicht für das,\footnote{„das“ fehlt in den Ausgaben.} das du noch empfangen sollst, sondern für das, das du jetzund empfangen hast. Denn die Dankagung geht nur auf die empfangenen Güter, nicht auf das, das du noch empfangen solltest. Was dich nun wider diese Meinung führen will, das thue hinweg.

52. Die Bilderstürmer aber fahren zu, reißen die Bilder äußerlich ab. Das wollte ich nicht so fast aufheben. Aber sie setzen hinzu: es müsse sein, und es gefalle GOtt wohl. Damit machen sie nichts anders, denn daß sie die Bilder aus den Augen ziehen, und setzen sie den Leuten in das Herz, verkehren die Meinung dieses Gebots. Damit verleugnen sie GOtt, und rühmen sich noch, sie reißen die Bilder um nach GOttes Befehl und Wort. Den Teufel auf ihren Kopf; wenn sie eines abreißen, so richten sie wohl zwanzig in den Herzen wieder auf; und von denselben falschen Vertrauen, daß der Pöbel meint, er thue GOtt einen Gefallen mit dem umreißen der Bilder,\footnote{So die Jenaer. Wittenberger: „bild umbreißen“; Erlanger: „umbreißen Bild“.} sagen sie kein Wort.

53. So ist das nun der rechte Verstand, daß Heiden und Juden den HErrn für einen GOtt haben, der alles umsonst gebe etc., es sage es Mosis, oder wer da will. Ueber das hat GOtt dem jüdischen Volke geboten durch Mosen, daß sie nicht Bilder haben sollten, damit sie einen Gottesdienst anrichten wollen. Die andern Bilder sind auch den Juden nicht verboten, wiewohl sie so närrisch sind gewesen, und haben es gedeutet, es zieme ihnen gar kein Bild zu haben. So wird nun in Summa im ersten Gebote erfordert ein rechtschaffener Glaube und Zuversicht zu GOtt, und wird nichts Aeußerliches darinnen geboten. Den kann aber niemand haben, der Heilige Geist gebe ihn\footnote{In den Ausgaben: „innen“.} denn zuvor in das Herz.

\subsection{Die Zusagung und Dräuung des ersten Gebots}

54. Das dritte Stück des ersten Gebots ist eine tröstliche Zusagung. Wir haben gehört am ersten die Meinung, was GOtt damit verstanden haben will, da er spricht: „Ich bin der HErr, dein GOtt“, item, daß wir keine fremden Götter neben ihm sollen haben. Zum andern legt er auch aus, was er damit will verstanden haben, so er spricht, daß sie nicht andere Götter neben ihm sollen haben, nämlich die Bilder, die man anbetet; und daß er allein den Juden das vorhält, und nicht den Christen und Heiden; und will dadurch von ihnen abwenden alle Ursache, dadurch sie einen falschen Glauben überkommen möchten. Nun, zum dritten, ist dahinten geblieben die Dräuung und die Verheißung. Denn GOtt thut alles beides in der Schrift, dräuet oder schreckt, und tröstet; schreckt die Uebertreter und Gottlosen, tröstet die Gläubigen, die seine Gebote halten, und spricht, wie folgt:

\begin{schriftzitat}
V. 5. Denn ich der HErr, dein GOtt, bin ein starker Eiferer, der da heimsucht der Väter Missethat an den Kindern, bis in das dritte und vierte Glied, die mich hassen.
\end{schriftzitat}

55. In diesem Stücke ist zu merken, daß GOtt hier zum ersten dräuet, und zum andern tröstet.

\quellseite{533}
[Er] setzt vorhin, er sei ein starker Eiferer, der da heimsuche der Väter Mißthat \&c., darnach setzt er die Barmherzigkeit hernach, und spricht: Und thue Barmherzigkeit an viel Tausenden, die mich lieb haben \&c. Denn das ist GOttes Weise und Natur, wie er denn hin und wieder in der Schrift auch mit der That beweist, daß er am ersten dräuet, schreckt, und das Herz verzagt macht, darnach tröstet und richtet [er] das Herz wiederum auf; tödtet vorhin das Fleisch, nachmals macht er den Geist lebendig. Das pfleget er zu thun; darum ändert er hier\footnote{„hier“ fehlt in der Wittenberger und in der Jenaer.} diese Ordnung nicht. Welchen er nun aufrichten will, den stößt er vorhin zu Boden; welchen er lebendig machen will, den tödtet er zuvor; den er fromm machen will, den macht er vorhin zum Sünder; und wen er reich machen will, den macht er zuvor arm; welchen er gen Himmel haben will, den stößt er vor in die Hölle; daß also das Schrecken allezeit vorher geht, der Trost und die Freude hernach folgt, wie hin und wieder die Schrift anzeigt.

56. Also spricht auch Hanna, die Mutter Samuels, 1 Sam. 2, 6. 7.: „Der HErr tödtet und gibt das Leben, führet in die Hölle und wieder heraus, der HErr macht arm, und macht reich, erniedriget und erhöhet“ \&c. Item, 5 Moi. 32, 39.: „Kein GOtt ist neben mir, ich kann tödten und lebendig machen, was ich zerschlage, das heile ich.“ Das ist sein Ruhm und Titel, den will er allein führen. Er zürnt und schlägt zuvor, darnach ist er freundlich und heilt, gibt sein tröstlich Evangelium, damit tröstet er wiederum, und beut seine Gnade und Freundschaft an. Der Teufel aber verkehret diese Ordnung GOttes, und macht es gleich Widersinns.

57. Zum andern ist hier zu merken, daß uns Heiden und Christen auch nicht angeht (denn man muß den Rottengeistern allenthalben den Weg verlegen) solche Träuung und Verheißung; denn wir haben andere Träuung und Verheißung oder Zusagung; diese aber gehen allein wie andere Stücke, oben [§ 12 ff.] erzählt, die Juden an, nicht uns. Denn Moses ist ein Lehrer des jüdischen Volks, darum sind auch seine Worte alle dahin gerichtet, daß sie allein die Juden betreffen. Wiewohl die Gebote GOttes allen Menschen in die Herzen geschrieben sind, so werden doch die Herzen durch den Teufel so sehr verfinstert, daß man sie nicht sehen noch erkennen kann; GOtt aber erinnert nun die Juden der Gebote, die in ihre Herzen eingedrückt sind, gibt ihnen über das natürliche Licht auch ein geschrieben Gesetz, ja, trägt ihnen dasselbe\footnote{Erlanger: „das selbs“.} mündlich vor, daß sie sehen, wie es im Herzen geschrieben steht.

\section{Das Gesetz ist natürlich im Herzen}

58. Wenn aber das natürliche Gesetz nicht von GOtt in das Herz geschrieben und gegeben wäre, so müßte man lange predigen, ehe die Gewissen getroffen würden; man müßte einem Esel, Pferde, Ochsen oder Rinde hunderttausend Jahre predigen, ehe sie das Gesetz annähmen, wiewohl sie Ohren, Augen und Herz haben, wie ein Mensch; sie können es auch hören, es fällt aber nicht in das Herz. Warum? Was ist der Fehl? Die Seele ist nicht darnach gebildet und geschaffen, daß solches darein falle. Aber ein Mensch, so ihm das Gesetz wird vorgehalten, spricht er bald: Ja, es ist also, [ich] kann es nicht leugnen. Das könnte man ihn so bald nicht überreden, es wäre denn zuvor in seinem Herzen geschrieben.

59. Weil es nun zuvor im Herzen ist, wiewohl dunkel und ganz verblieben, so wird es mit dem Worte wieder erweckt, daß ja das Herz bekennen muß, es sei also, wie die Gebote lauten, daß man einen GOtt ehre, liebe, ihm diene, weil er allein gut ist und Gutes thut, und nicht allein den Frommen, sondern auch den Bösen. Wiewohl der Teufel stark wehrt, daß der Mensch weder fühle, erkenne, noch vollbringe; ja, der Mensch vermag auch der keines zu thun ohne das Werk und Licht des Heiligen Geistes.

60. Also ist das nun eine helle, klare Anzeigung, daß diese Worte, damit GOtt hier dräuet und verheißt, allein die Juden betreffen, wie jezund angezeigt ist. Das ist aber gewißlich wahr, es soll sich auch ein jeglicher leiblich darauf verlassen, dies Gebot gehe die Heiden oder Juden an, es treffe wen es wolle: Welcher GOtt nicht fürchtet, den straft er; welcher aber ihn ehrt und fürchtet, dem ist er gnädig, ehrt ihn wieder, wie 1 Sam. 2, 30. steht: „Wer mich ehret, den will ich ehren; wer aber mich verachtet, wird verschmähet sein.“ Solche Träuung aber hat er den Heiden nicht mündlich lassen verkündigen,
\quellseite{534}
sondern allein den Juden; wie der 147. Psalm, V. 19. f., spricht:\footnote{„spricht“ fehlt in der Erlanger.} „Er zeiget Jakob sein Wort, Israel seine Sitten und Rechte. So thut er keinen Heiden, noch lässet sie wissen seine Rechte.“

\subsection{Ich bin ein starker Eiferer}

61. Daß er nun spricht: „Ich bin ein starker Eiferer“, ist auch den Juden gesagt, nicht uns. Denn er dräuet hier mit einer leiblichen Strafe, daß er den Vater also strafen wolle, wo er wider sein Gebot handele, daß es auch das Kind muß entgelten, bis in das dritte oder vierte Glied. Aber an der Seele straft er nicht also, sondern allein leiblich und äußerlich. Als, wenn er einen straft an Haus, Hof, Acker und Wiesen, Geld und Gut, wie er dem frommen Hiob that, den er auch am Leibe angriff [Hiob 2, 7.], wiewohl anders, denn wenn er die Bösen straft, davon er hier redet. Item, wenn er Wohlthat beweiset in tausend Glied, das ist auch leiblich, und ist so viel: Ich will euch Gutes thun an Kindern und Kindes-Kindern, nicht allein bis in das vierte, sondern in viel Geschlechte und Glieder, daß das Haus in gutem Wesen\footnote{Wittenberger: Wege.} bleibe, und das Geschlecht durch und durch währen soll, wie Davids Geschlecht durch viel Geschlechte bis auf Joseph und Mariam, ja bis auf Christum währt [Luc. 3, 23. f.].

62. Wiederum, der Gottlosen Geschlecht soll\footnote{Erlanger: sollen.} ausgerottet werden, wie Jerobeams und Ahabs Geschlechten geschah, die GOtt also ausgerottet hat, daß von ihnen geschrieben steht, 1 Kön. 14, 10. und 2 Kön. 9, 9., „er wolle sie ausrotten, auch den, der an die Wand pisset, und die Nachkommen des Jerobeam ausfegen, wie man Dreck ausfegt, bis [es] ganz mit ihm aus ist“. Und von Ahab steht, „er wolle sein Haus machen, wie das Haus Jerobeams“ etc. Solche Ehre kann er thun denen, die ihn ehren; verachten und zerschmettern die, so ihn verachten.

63. Im neuen Testamente aber ist diese Strafe aufgehoben. Denn wir sehen, wie Gott zufährt, und läßt die ärgsten, verzweifelsten Buben reich werden, in allem Muthwillen leben und das beste Glück haben, also, daß sich auch die Frommen daran ärgern und stoßen. Wiederum, die Gerechten läßt er schinden und schaben, wohl quälen und martern, schickt ihnen alles Unglück auf den Hals, also, daß sie dünkt, es gehe alles wider den Strom, wie David gar hübsch im 73. Psalm, V. 9. ff., anzeigt. Doch David sieht daselbst auf das Ende beider, der Gottlosen und der\footnote{„der“ fehlt in der Wittenberger und in der Jenaer.} Gottseligen, daß GOtt also die Gottlosen mästet zur Schlachtbank, und die Christen will er mit solchem Unglück im Zaume halten, daß sie nicht von GOtt weichen, nicht löcken\footnote{löcken — übermüthig hinausgeschlagen.} und zu geil werden. Darum sich keiner das Kreuz zu tragen weigern soll, denn unser HErr und Meister, Christus, hat es selber getragen [Matth. 10, 38. 27, 31.]. Nun ist der Jünger ja nicht mehr, denn der Meister, und der Knecht nicht mehr, denn der HErr [Matth. 10, 24.].

64. Im Neuen Testamente aber haben wir andere Dräuung, nämlich, da wird uns der ewige Tod gedräuet, der Zorn GOttes mit dem jüngsten Tage, die Hölle und das ewige Verdammniß, Marc. 16, 16.: „Wer nicht glaubt, wird verdammt werden“; Joh. 3, 18.: „Wer an Christum nicht glaubt, der ist schon verdammt“;\footnote{Erlanger: aber.} Joh. 3, 36.: „Wer dem Sohne nicht glaubt, der wird das Leben nicht sehen, sondern der Zorn GOttes bleibt über ihm“; item, Röm. 2, 5.: „Du sammelst dir selbst einen Schatz des Zorns auf den Tag des Zorns und der Offenbarung des gerechten Gerichts Gottes“ etc. Item, wiederum haben wir auch im Neuen Testamente nicht leibliche, sondern geistliche und ewige Verheißung, Joh. 3, 16.: „Wer da glaubt, hat das ewige Leben“; item, Joh. 7, 37, 38.: „Wen da dürstet, der komme zu mir, und trinke. Wer an mich glaubt (wie die Schrift sagt), von deß Leibe werden Flüsse des lebendigen Wassers fließen“; item, Joh. 4, 14.: „Das Wasser, das ich gebe, wird in ihm ein Brunnen des Wassers werden, das in das ewige Leben quillet.“

65. In Mose aber klinget es anders: die Zusagung will er ihnen gehalten haben, wenn sie auch fromm sind äußerlich. Und ist das die Meinung: Wenn die Juden sich vor den Bildern hüten, und leben in der Furcht GOttes, so will er sie annehmen, und sie an Leib und Seele versorgen. Leiblich sollen sie genug haben; und wenn sie ihm erkennen als einen rechten GOtt von Herzen, soll auch die Seele selig werden.

66. Diesen Text, da GOtt verheißt denen,
\quellseite{535}
die ihn fürchten und ihm allein dienen, Wohlthat, und bräuet Unglück denen, so ihn verachten, haben viel Propheten in ihren Schriften und Weissagungen hoch angezeigt, damit auch dem Volke gedräuet zukünftiges Unglück, um ihrer Sünde willen; und\footnote{Jenaer: auch.} Glück und Seligkeit verheißen,\footnote{Im Original und in der Wittenberger: Verheißung.} wo sie von ihrem gottlosen Leben abließen, und sich zu dem Herrn bekehrten.\footnote{Erlanger: Und.} Wiederum sind auch viel falsche Propheten dadurch gestärkt worden in ihrem Frevel und falschem Vornehmen, haben den Text nicht recht verstanden, haben sich mit den frommen Propheten gebadert, wie etliche Psalmen anzeigen, vornehmlich der 144., in welchem Psalm der Prophet David klagt über dieselbigen falschen Propheten, und spricht, V. 11. ff.: Ihr Mund redet unnütz, trösten das Volk, da sie es erschrecken sollen, und sprechen: „Unsere Söhne wachsen auf wie die Pflanzen in ihrer Jugend, und unsere Töchter wie die ausgehauenen Erker, gleichwie die Paläste“, und alle Winkel sind voll, Haus, Kisten, Kasten, Keller und Boden, sie haben schöne Rinder, Schafe, viel Freunde, Kinder, auf dem Felde sind alle Dinge wohl gerathen, und singet\footnote{Erlanger: singet.} alles, das da steht und geht: „Wohl dem Volke, dem es also gehet.“

67. Das war die Predigt der falschen Propheten im alten Testamente, die sie auf diesen Text gründeten, verführeten damit das arme Volk, verhießen ihnen Frieden, und war doch Unfriede, Segen, und war doch Vermaledeyung. Da haben denn die frommen Propheten dawider gepredigt, haben sie wollen auf einen rechten Verstand führen, aber die falschen Propheten haben sich wider sie gesetzt, wie dem heiligen Jesaia und Jeremia geschah, und sagten dann die falschen Propheten: Siehe, was ist hier geschrieben? GOtt will vom Himmel wohlthun denen, die seine Gebote halten, und übelthun und strafen die Gottlosen. Uns aber thut er wohl, gibt Haus und Hof voll, Geldes und Gutes genug, und schöne Kinder dazu; alles darum, daß wir seine Gebote halten. Wiederum, straft er die mit Armuth und Unglück, die seine Gebote nicht halten.

68. Darauf geht auch das ganze Buch Hiob durch und durch: da urtheilen den\footnote{„den“ fehlt in der Jenaer.} Hiob seine Freunde, und haben sich mit ihm, sprechen [Hiob 4, 7. ff. 36, 6. f.]: Hältst du GOtt für einen ungerechten GOtt, der die Frommen strafe? Kehre es um: Den Frommen thut er wohl, die Bösen aber straft er. (Gleich also urtheilten auch die falschen Propheten; wenn sie sahen, daß es einem übel ging, daß einem ein Rad über ein Bein ging, oder wenn einer einen Arm entzwei fiel, oder sonst ihm ein Unglück widerfuhr, flugs urtheilten sie: Er hat es verdient, GOtt straft ihn, er ist ein Bube. Und das war eine gemeine Regel, aus diesem Texte gezogen; so folgte denn das Volk denselbigen Verführern, urtheilte auch wie sie; wenn sie sahen arme oder kranke Leute, gründig, Sünder \&c., sprachen sie: Rühre mich nicht an, ich bin rein, GOtt ist mir hold, er straft dich um deiner Sünde willen.

69. In diesem Wahn waren auch die Apostel, da sie Christum fragten des Blindern halben, Joh. 9, 1. 2.: „HErr, wer hat gesündigt, dieser, oder seine Eltern?“ Also auch Simon, der den HErrn zu Hause geladen hatte, da die arme Sünderin dem HErrn die Füße wusch, urtheilte auch nach dem äußerlichen Scheine, rümpfte die Nase über die Frau, sie stank vor seinen Augen, und murrte wider Christum, und sprach bei ihm selber [Luc. 7, 38. 39.]: „Wenn dieser ein Prophet wäre, so wüßte er, wer und welch ein Weib das ist, die ihn anrühret; denn sie ist eine Sünderin.“ Also mußte auch Christus den heiligen Pharisäern und Schriftgelehrten nicht allein nicht der rechte, wahrhaftige Messias sein, sondern auch kein Prophet. Ja wohl, soll dieser Bettler Christus sein? Sehe die Brille auf die Nase, er ist vom Teufel; denn er lehrt wider Mosen, hat mit den aussätzigen, blinden, lahmen Leuten zu schaffen, und, das am ärgsten ist, menget sich unter Huren und Buben, mit denen geht er um, isset und trinket mit ihnen; pfui des losen Mannes! In Summa: Sie hielten es\footnote{„es“ fehlt in der Erlanger.} dafür, daß GOtt den Frommen gebe, was sie wollten, und den Bösen, wie sie verdienten, und wollten daraus schließen, weil es ihnen an Leib und zeitlichem Gute wohl von Statten ging, daß sie die wären, denen GOtt um ihrer Frömmigkeit willen so wohl thäte. Es fehlte aber weit.

70. Was ist denn nun der Mangel? Daran fehlt es, sie haben den Text nicht recht verstanden
\quellseite{536}
. Es ist wahr, und fehlt nimmer: Was GOtt zusagt, es sei zeitlich oder ewig, das hält er [Ps. 33, 4.], gibt es auch treulich; aber sie haben den Text unfleißig\footnote{So die Tenger; Erlanger: „allein angesehen“; Wittenberger: den Text allein leiblich verstanden.} angesehen, und leiblich verstanden. Das ist wahr, welche GOtt erkennen, lieben, ehren und glauben ihm, als ihrem lieben GOtt und Vater, durch einen reinen, wahren und ungefärbten Glauben, die sind sicher aller Zusagung GOttes. Hier aber ist nicht hell und klarlich ausgedrückt und geschrieben, daß er wolle den Bauch versorgen und keine Noth leiden lassen, daß er mitten im Hunger wolle ernähren, wenn gleich kein Heller im Vorrathe da ist, und daß er wolle mitten unter den Feinden schützen, ohne alles Schwert, Spieß und Harnisch, und in Krankheit und Anfechtung fröhlich und gesund machen.

71. Die Juden aber gedachten also: Wenn du GOtt dienst, so mußt du nicht des Brods vom Himmel gewärtig sein, dir wird keine gebratene Taube in das Maul fliegen. Es wird also zugehen: Geld genug wird vorhanden sein, Kisten und Kasten voll sein, Taschen und Beutel werden paussen;\footnote{paussen = bauschen, dick sein.} du wirst schöne Frauen und Kinder, ein groß Gesinde, ein sicher verwahrt Haus haben, in Wohllust leben, und Vorrath auf zehn, zwanzig oder dreißig Jahre haben, und nicht alle Stunden vom Himmel warten, wenn du nichts in der Faust hast. In Summa, es wird kein Mangel da sein, sondern überflüssig genug an allen Dingen. Also haben sie wollen gewiß sein, und frisch geschlossen, wo genug vorhanden ist gewesen, und ein großer Vorrath: Da ist GOtt, das ist ein seliger Mann; wo aber Mangel war, hinten und vorne anstand: Der ist nicht fromm, er ist vermaledeyet, und ein Bube in der Haut, GOtt ist nicht mit ihm. Ja, es hat auch den großen Heiligen in diesem Stück gefehlt.

72. Das will aber GOtt hier nicht, sondern das ist seine Meinung, daß er gnädig sein will, und keinen Mangel die Seinen lassen leiden, es sei Vorrath vorhanden oder nicht, es sei der Beutel voll oder leer, und wenn gleich kein Körnlein auf dem Boden wäre, der Wein und das Bier ganz und gar in Keller gelaufen, und der Hagel alles Getreide erschlagen hätte, kein Heller im Beutel wäre, noch müßte kein Mangel da sein, weder Hunger noch Durst schaden.

73. Also will GOtt sie an sich hangen, und nicht an das gegenwärtige, vergängliche Gut. Sie wollten es aber nicht also verstehen, sondern verstanden, daß solche Verheißungen GOttes erfüllt würden, wenn sie keinen Mangel litten, keine Widerwärtigkeit versuchten, sondern nach allem Willen und Lust lebten. Das heißt denn, GOttes nicht bedürfen, und nimmermehr nach ihm fragen. Ja, solche wollten, daß GOtt seines Dinges droben im Himmel wartete, und ließe sie hernieden ihrer Gulden und Mammons warten. Das heißt denn nicht glauben, sondern fühlen; nicht von GOtt gewarten, sondern vorhin in der Faust haben. So führen denn die falschen Propheten zu, plauderten in das Volk, wo es glücklich zuging: Da, da ist Glück und Heil, da wohnt GOtt. Also mußte das Urtheil bei ihnen beständig wahr sein: Wer gesündigt hat, den straft GOtt, daß [es] ihm übel geht, und wer recht thut, dem gibt er die Fülle.

74. GOtt aber hat diesem leiblichen, groben Volke leibliche Verheißung vorgelagen; sie sollten aber darunter verstanden haben die geistliche Zusagung. Er wollte sie also zu sich gewöhnen, und also lehren, daß sie von ihm gewarten sollten, er würde sie gewiß ernähren, und versorgen zeitlich und ewiglich. Also hat er beide Verheißungen, leibliche und geistliche, in einander gemenget, wiewohl es nicht klarlich und hell lautet und ausgedrückt ist. Er hat aber die geistlichen Verheißungen an die leiblichen also geknüpft, daß sie von einander nicht mögen getrennt werden. Sie sollten täglich das Brod von ihm gewarten, und gewiß sein, daß er sie für und für ernähren wollte; sie aber wollten in den Verstand nicht, verließen sich\footnote{Erlanger: ließen sie.} auf ihre Heiligkeit, meinten, GOtt müßte ihnen wohl geben ihre Nothdurft, um ihrer Frömmigkeit willen; machten also aus GOtt einen Abgott, welches sie nicht sehr nöthig dürften. Denn wer keinen Mangel hat, ist fröhlich, und geht ihm nach alle seinem Willen, der darf Gottes Güte, Trost und Hülfe nicht, darum möchte er droben seiner Engel warten, so wollten sie herunten ihrer Gulden warten. Das heißt denn, dem Mammon vertraut und gedient, und nicht Gott, haben
\quellseite{537}
also zweien Herren wollen dienen und auf beiden Seiten hinken.

75. GOtt kann auch wohl einem Buben den Ranzen voll geben; [es] folgt aber darum nicht, daß er fromm sei, denn er will ihn hier bezahlen. Wiederum, läßt er es einem Frommen schwer und blutsauer werden; er ist aber darum nicht sein Feind, ja, er meint es väterlich gut mit ihm. Denn das ist die höchste Strafe, daß GOtt nicht straft, sondern stille hält, und läßt einen nach seinem Muthwillen dahin in Tag leben.

76. Wer nun ohne GOtt lebt, der wird seines Hellers froh, und genießt alles seines Guts wenig, denn er hat ein böses Gewissen, wie die Schrift sagt Jes. 57, 21.: „Die Gottlosen haben keinen Frieden“; und geht ihnen wie dem reichen Manne im Evangelio, davon Lucas am 12. Capitel, V. 16–21., schreibt und spricht: „Es war ein reicher Mann, deß Feld stand wohl; und er gedachte bei ihm selbst, und sprach: Was soll ich thun? Ich habe nicht, da ich meine Früchte hinsammle. Und sprach: Das will ich thun: Ich will meine Scheuern abbrechen, und größere bauen, und will darein sammeln alles, was mir gewachsen ist, und meine Güter, und will sagen zu meiner Seele: Liebe Seele, du hast einen großen Vorrath auf viel Jahre, habe nun Ruhe, iß, trink, und sei fröhlich. Aber GOtt sprach zu ihm: Du Narr, diese Nacht wird man deine Seele von dir fordern, und weß wird es sein, das du bereitet hast? Also geht es, wer ihm Schätze sammelt, und ist nicht reich in GOtt.“ Solche haben kein Herz zu GOtt, darum alle Augenblicke fürchten sie sich vor dem Tode, sind nicht sicher, nicht allein inwendig, sondern auch auswendig; fürchten, das Haus werde abbrennen, die Diebe werden kommen und ihnen die Gulden stehlen; da ist kein fröhlich Herz, keine Freude, keine Ruhe nicht, weder Tag noch Nacht. Es scheint wohl, als sind sie reich, guter Dinge, als sei ihnen GOtt günstig; aber es ist mit ihnen das Widerspiel. Was ist mir aber das für eine Freude und gut Leben, wenn einer nimmermehr gutes Muths ist, sorgt ohne Unterlaß, und gedenkt nur, wie er den Mammon hüte und mehre. Darum sagt Salomon gar fein im Prediger Cap. 9, 9., daß nichts Besseres ist in diesem Leben, denn daß ein Mensch fröhlich sei mit seinem Weibe, und entschlage sich aller Sorgen. Das kann aber niemand, auch kein König thun, wenn er gottlos ist; denn er fürchtet den Tod, seine Feinde \&c. [Von] solcher Furcht kann ihn nichts erretten \&c.

77. GOtt aber will, daß das Herz inwendig recht stehe, es stehe äußerlich, wie es immer kann, es sei Vorrath da oder nicht, wie der Prediger Cap. 9, 7–9. spricht: „Gehe hin, und iß dein Brod mit Freuden, trink deinen Wein mit gutem Muth; denn dein Werk gefällt GOtt. Laß deine Kleider immer weiß sein, und laß deinem Haupte Salbe nicht mangeln. Brauche das Leben mit deinem Weibe, das du lieb hast, so lange du das eitle Leben hast, das dir GOtt unter der Sonne gegeben hat“ \&c. Das kann kein böser Mensch thun; steht, wie gesagt ist, allwege in Sorgen. Denn GOtt ist nicht in seinem Herzen, [er] fürchtet sich vor seinen Feinden, vor Dieben, vor den Motten, daß sie ihm nicht seinen Gott zerbeißen und fressen, und, wie die Schrift [3 Mos. 26, 36.] sagt, er fürchtet sich vor einem rauschenden Blatt. Denn wer GOtt zum Feinde hat, der hat alle Creaturen zu Feinden.

78. Das ist aber wahr, GOtt gibt etlichen Heiligen nicht allein den Geist und Muth inwendig, sondern überschüttet sie auch äußerlich mit Reichthum. Also gab er dem frommen Abraham viel Reichthum; den David machte er zu einem mächtigen Könige; den Hiob machte er reicher denn alle, die gegen Morgen wohnten; aber dabei hatten sie auch den Heiligen Geist, nahmen das von GOtt an als ein Geschenk, waren guter Dinge dabei, ihr Herz hing nicht an den Reichthümern, sondern an GOtt; wie David selber alle Reichen ermahnt Psalm 62, 11. und spricht: „Fället euch Reichthum zu, so hänget das Herz nicht dran.“

79. Also, alle andern frommen Väter und GOttes Freunde sind ihnen gleich gewesen, haben ihre Hoffnung nicht auf zeitliche Güter, sondern auf GOtt gesetzt. Ihr Herz stand also: Was willst du groß pochen auf Reichthum, Ehre, Gesundheit \&c., es bleibt alles nicht, heute oder morgen mußt du alles verlassen und davon. Also sind sie reich gewesen, beide äußerlich und innerlich. Es sind aber solcher gar wenig, der andern aber ist viel.

80. Item, man findet auch viel frommer Menschen, die großen Mangel an der Nahrung haben, doch ein gut, aufrichtig Herz zu GOtt, und sind die, die fest halten an seiner Zusagung, daß er sie wolle ernähren, derselben vertrauen sie von ganzem Herzen, und GOtt macht es auch mit
\quellseite{538}
ihnen also, daß sie bei einem Pfennig oder Heller, ja, mit einem leeren Beutel einen besseren Muth und Gewissen haben denn ein Fürst, König oder Kaiser mit alle seinem Gut und Reichthum, und schmeckt ihnen ein Bissen Brods und Trunk Wasser besser, denn großen, reichen Herren alle ihre Leckerbißlein und der beste Wein.

81. Was macht das? GOtt, der da spricht: Ich will dich versorgen, wirf deine Sorge auf mich, ich will dich ernähren [Ps. 55, 23.]. Item [Matth. 6, 33.]: „Trachtet am ersten nach dem Reiche GOttes und nach seiner Gerechtigkeit, so wird euch solches alles zufallen.“ Und wiewohl es vor Augen nicht vorhanden ist, so soll es doch euch reichlich widerfahren, mehr denn die es hundertfältig im Vorrath haben. Das hat der heilige Paulus wohl erfahren, darum darf er es auch von sich und andern Christen schreiben, 2 Cor. 6, 10.: Wir sind „als die Traurigen, aber allezeit fröhlich; als die Armen, aber die doch viel reich machen; als die nichts haben, und doch alles inne haben“; und zum Philippern am vierten, V. 12.: „Ich weiß nichts zu sein, und weiß auch hoch herzufahren; ich bin allenthalben und in allen Dingen geschickt, beide satt sein und hungern, beide übrig haben, und Mangel leiden.“

82. Man findet auch viel Gottlose, denen alles Unglück zuhanden kommt, haben dazu kaum das liebe Brod im Hause; die können aber nicht so guter Dinge sein, wie die Gottseligen, sondern fahen an zu schelten, fluchen und lästern, werden ungeduldig und thöricht, sprechen: Hat mich der Teufel in das Leben geführt; wofür hält mich GOtt, daß er mir auch nicht so viel als meinen Nachbarn gegeben hat? Er hat mich zu unrechter Zeit geschaffen. Wo fehlt es denen? Da fehlt es, daß ihnen ihr Herz unrein ist, daß sie keinen Glauben haben, nichts von GOtt wissen, in Summa, ohne GOtt sind, wie Paulus spricht Eph. 2, 12.

83. Das ist nun der rechte Verstand der Worte, so Moses spricht: „Ich bin ein starker Eiferer, der da heimsucht der Väter Missethat an den Kindern, bis in das dritte und vierte Glied, und thue Barmherzigkeit an viel tausend, die mich lieb haben und meine Gebote halten.“ Er verheißt den Juden leibliche Wohlthat. Er will es aber also verstanden haben, daß er sie an sich will binden, als wollte er sagen: Wenn du gleich nichts im Kasten, im Keller oder Beutel hast, so sollst du es doch bei mir haben. Warum? Ich bin der HErr, dein GOtt.“ Bin ich dein GOtt, so kann ich dir geben ohne allen Vorrath, meine Hand ist nicht verkürzt [4 Mos. 11, 23. Jes. 50, 2.]. Etliche haben ihren Schatz an Geld und Gut, etliche an mir allein. Was ist daran gelegen, du habest an Geld oder Gut Vorrath,\footnote{So die Jenaer. Wittenberger: bei dir Vorrath; Erlanger: an ihn Vorrath.} oder nicht: er\footnote{Wittenberger und Jenaer: es.} ist dir bei mir ja so gewiß, und wohl gewisser, denn wenn du ihn liegen hättest auf einem Haufen. Denn der Vorrath, den du bei dir hast, der kann verbrennen, ersaufen, gestohlen, oder von Motten und Würmern gefressen werden und umkommen. Aber bei mir bleibt er gewiß; und ob du gleich nichts hast, so bist du dennoch reich und versorgt, denn du hast mich, und in mir alle Dinge.

84. Wiederum, ein Gottloser hat nichts; und ob er gleich vor Augen reich ist, so hat er doch Mangel, hat Unglück gar genug; und wiewohl er alles im Vorrath mehr denn zu viel hat, so ist er dennoch ärmer denn arm. Denn was schadet Lazaro seine Armuth? Was nützte dem reichen Manne sein großer Reichthum? [Luc. 16, 24. 25.] Darum ist niemand reich, er heiße Kaiser oder Papst etc., denn wer in GOtt reich ist [Luc. 12, 21.].

85. So sind nun das die Verheißungen: Wenn ihr GOtt dienen, und ihn für einen GOtt erkennen werdet, so sollt ihr genug haben, ihr habt es im Vorrath oder nicht. Es wird ihnen aber grob und kindisch vorgelegt. Denn wie das Gesetz sie wie ein Zuchtmeister äußerlich geführt hat, also verheißt ihnen auch GOtt äußerliche Wohlthat; es ist eine Kinderschule und Kinderlehre.

86. Das neue Testament aber kehret es um, und sähet es am Geistlichen und Innerlichen an, hält vorhin den Glauben vor, und spricht: Wenn ihr glaubt und fromm seid, so sollt ihr genug haben, Matth. 6, 33.: „Trachtet am ersten nach dem Reiche GOttes und nach seiner Gerechtigkeit“, so wird der Bauch wohl versorgt. Das haben auch beide Testamente, aber das alte Testament sähet am Leiblichen an, und ist doch das Geistliche darunter verborgen, nämlich, daß man im Glauben müsse genug haben. Es ist ein einerlei Ding im alten und neuen Testamente; es ist aber eine andere Ordnung. Das neue
\quellseite{539}
fähet am Geistlichen an, schleift das Leibliche mit sich hernach. Als: Ich glaube, GOtt könnte mir den Bauch nähren, wenn ich gleich nichts im Vorrath habe; und wenn ich ihn auf einen Tag vertraue, so vertraue ich ihm auch auf ein, zwei oder drei Jahre etc. Und bringt also der innerliche Glaube mit sich, dadurch die Seele reich ist, daß ich auch glaube, ich werde, was den Leib betrifft, reichlich haben.

87. Abraham ward zugesagt [1 Mos. 12, 7. 13, 15.], er sollte das Land Canaan einnehmen; und hat es doch nicht eines Fußes breit eingenommen, Apost. 7, 5. Er traute aber GOtt, und glaubte, er würde es einnehmen, starb also dahin, und glaubte bis an sein Ende,\footnote{In der Jenaer stehen hier die Worte: „und hatte doch kein Erbtheil drinnen, auch nicht eines Fußes breit“. Die Erlanger berichtet, dies sei eine Note Walchs.} er würde es einnehmen, und ist doch leiblich nicht dahin gekommen. Im Glauben aber hat er es eingenommen, denn seinem Samen ward es gegeben. Denn GOtt hat ihm also zugesagt, er und sein Same mit ihm sollte das Land einnehmen. Also ist unter der leiblichen Verheißung eine geistliche Zusage mit eingelaufen, auf welche auch Abraham mehr hat gesehen, denn auf die leibliche Verheißung, wie die Epistel zum Hebräern anzeigt, Cap. 11, 9. Davon redet auch St. Stephan in den Geschichten der Apostel, Cap. 7, 5. f. Also ist diese Zusagung nicht zu den Heiden, sondern zu den Juden gesagt worden.

88. Also auch, daß er zürnt in das dritte oder vierte Glied, gehört auch in das Gesetz, zeigt einen leiblichen Zorn und Strafe an, nicht eine ewige. Im neuen Testamente aber ist es anders: das Kind wird nicht für den Vater gestraft. Es geht aber also zu: Ein jeglicher glaubt für sich selber, der Vater für sich, der Sohn für sich, die Mutter und Tochter für sich; wie auch in Ezechiel [Cap. 18, 20.] geschrieben steht: Ich will das aufheben, daß der Sohn nicht soll leiden für den Vater, und der Vater nicht für den Sohn. Weil er aber den Juden leibliche Zusage verhieß, Träumung träuete, begnadete oder strafte er sie auch leiblich.

89. Im neuen Testamente aber geht die Verheißung an dem Geistlichen an, denn er verheißt die ewige Seligkeit. Also auch, mit der Strafe geht es an der geistlichen Strafe an: „Wer nicht (spricht er Marc. 16, 16.) glaubt, ist verdammt.“ Er spricht nicht: Ich will den Sohn um des Vaters willen verdammen; sondern, der Vater muß selber herhalten. Es ist ein gestraft Urtheil gefällt, und nicht auf die Güter, sondern auf die Person. Wenn man nun die Glieder zählt, so geht es auf die Juden, die er leiblich gestraft hat, als, an den Gütern, Haus, Hof, Acker, Wiesen, Weib und Kindern, Knecht und Magd; wie denn in den Historien bewiesen wird.

\subsection{GOttes natürliches Werk, wohlthun}

90. Das ist auch wohl zu merken, daß GOtt im Zorn und Dräuen nur bis auf das vierte Glied zählt und straft, in der Liebe aber beweist er seine Barmherzigkeit, nicht allein in das vierte, fünfte, zehnte oder zwanzigste Glied, sondern in viel tausend, und ist ein recht feiner und lieblicher Spruch. Der Pabst hat es GOtt auch nachgethan, wollte ihm gleich werden, verfluchte bis in das neunte Glied. Er hatte es aber nicht von GOtt, sondern von seinem Vater, dem Teufel, gelernt, der nichts kann denn lügen, morden und fluchen. GOtt aber macht es nicht also; er erzeigt sich, daß wir ihn sollen lernen erkennen für einen wahren GOtt, der natürlich tausendmal mehr geneigter sei zur Güte denn zum Zorn.

91. Darum ist sein eigen Werk, wohlthun. Zürnen aber heißt sein fremdes Werk, Jes. 28, 11. Es ist uns auch wohl vonnöthen, daß er mehr wohlthue, denn strafe. Denn wir Menschen können von Natur wohl leiden, daß uns GOtt wohlthue zehn, zwanzig, dreißig Jahre; wenn aber ein böse Jahr kommt mit Theurung, Krieg oder Pestilenz, so wollen wir verzagen, so ist kein GOtt mehr daheim, da ist nur lauter Zetergeschrei. Wenn uns GOtt ein Jahr oder zehn gütlich thut, so kennt es niemand, niemand dankt ihm dafür. Also kann unsere Natur die Gutthat wohl leiden, aber die Strafe will sie nicht leiden, murrt von Stunden an, und hätte doch lauter Zorn wohl verdient. GOtt aber hält uns das zugute, daß er mehr Güte denn Zorn erzeigt.

92. Also hier auch; [er] straft nur in das vierte Glied; aber seine Barmherzigkeit läßt er gehen bis in viel tausend. Wenn er so hart wollte strafen, als viel er Güte und Barmherzigkeit erzeigt, wer könnte vor ihm bleiben? Darum, wenn man die Güte will ansehen, so ist seine göttliche Plage und Strafe zu groß, ist
\quellseite{540}
auch keine zu schätzen gegen seiner Güte. Wir sehen es aber nicht, sondern sind blind. Allein, wenn das Unglück über unsern Hals kommt, dann sehen wir erst, was wir gehabt haben. Wenn GOTT so viel zürnte und strafte, so viel er freundlich und barmherzig ist und wohlthut, wie könnten wir es ertragen? Wer würde nicht verzweifeln?

93. Dem GOTT, der so barmherzig ist, der uns so viel Gutes erzeigt, der uns so mit unzähligen Wohlthaten überschüttet, können wir nicht zu gute halten, wenn er uns eine böse Stunde zuschickt; was wollten wir thun, wenn er uns so viel böse Zeit, Krankheit und allerlei Unglück zuschickte, so viel er uns mit guter Zeit, Gesundheit und allem Glück begnadet? Mancher geht dahin zwanzig Jahre, hat keinen Anstoß; wenn einmal ein Fieber kommt, das über drei Tage währt, so will er aus der Haut fahren.

94. Also ist es klar und hell genug, daß GOTT mehr zu\footnote{Wittenberger und Erlanger: „auf“ statt: zu.} Barmherzigkeit denn zu Zorn geneigt ist. Ihr wißt, wie viel Tausend erschlagen sind aus den Bauern; da erscheint GOTTes Zorn und Strafe. Aber dieser Zorn ist nicht zu schätzen gegen der Güte GOTTes. Denn dagegen haben viel Tausend das Leben, die gleiche Strafe verdient hätten, daß sie GOTT auch erwürgen hätte sollen lassen.\footnote{So die Wittenberger und die Jenaer. Erlanger: auch erwürget hätte lassen.} Solches ist nun den andern zum Exempel geschehen und zur Warnung, und ein Fuchsschwanz gewesen. Und GOTT zeigt an, was für Güte er damit uns beweise, daß er uns nicht alle so läßt umkommen, die wir es so wohl als die Bauern verdient haben. Wir aber erkennen es nicht, so lange, bis uns die Strafe heimkommt.

95. Wenn er mir ein Auge nähme, ein Bein bräche, einen Arm ab ließe hauen, schickte mir eine Krankheit zu acht Tage, daß ich in einem halben Jahre nicht arbeiten könnte, da würde ich ein Heulen und Klagen heben. Ich will aber nicht daneben gedenken, daß ich wohl zwanzig Jahre gute und gesunde Tage gehabt habe. Wenn er gleich das eine Auge nimmt, so sind doch die andern Glieder gesund. Also straft er nicht den hundertsten Theil, gegen den Wohlthaten zu rechnen, die er uns schenkt. Wenn wir aber sterben, so geht es ganz mit einander dahin, denn wir sind ihm einen Tod schuldig. Also, wenn er straft, sollen wir wissen, daß es nur ein Fuchsschwanz ist, denn es wird uns nicht der hundertste Theil vergolten, wie wir wohl verdienten.

\subsection{Frage: warum der Sohn für den Vater leide?}

96. Diese Frage, warum der Sohn für den Vater leide, hat der Prophet Ezechiel gehandelt, und spricht [Cap. 18, 2.]: „Ihr habt ein Sprüchwort unter euch in Israel, die Väter haben saure Weinbeere\footnote{Jenaer: „Heerling“.} gegessen, und uns sind die Zähne stumpf worden.“ Und Jeremias [Cap. 31, 29.] spricht: „Unsere Väter haben gesündigt, und sind dahin; wir aber müssen ihrer Sünde entgelten. Und [es] geht noch heutiges Tages [so]: wir sündigen und verdienen, daß es unsere Nachkommen entgelten müssen. Das soll man aber nicht also verstehen, daß [es] eine Strafe der Seele sei, daß das Kind um des Vaters willen verdammt werde; wie es denn Ezechiel selbst anzeigt am zehnten genannten Orte. „Die Seelen“, spricht GOTT durch Ezechiel [Cap. 18, 4.] „sind alle mein, des Vaters und des Sohnes, welche aber sündiget, die wird sterben“; sondern man soll es von der leiblichen Strafe verstehen. Er straft die Kinder um der Väter willen, läßt sie sterben, die doch sonst gleichwohl sterben müssen, und straft also die Väter, daß er zuweilen ein ganz Geschlecht ausrottet. Also segnet er die Väter in den Kindern, daß [es] viel Nachkommen genießen (leiblich), daß die Väter fromm gewesen sind.

97. Wenn er nun einen Vater strafen will, so nimmt er ihm die Kinder, wie er David seinen Sohn ließ sterben [2 Sam. 12. 18.]. Denn er straft uns an den Gütern, die er uns gegeben hat, und zeigt dabei an, daß die Kinder das allerbeste und liebste Gut unter allen Gütern sind; und wenn er nicht der Kinder schont, die das Beste und Liebste sind, wie viel weniger wird er schonen der Ochsen, Esel, Schafe und anderer Güter?

98. Das ist aber zu merken, daß GOTT eher zu den Frommen greift, denn zu den Gottlosen und verzweifelten Buben, also, daß ein Sprüchwort ist worden davon: Je ärger Schalk, je besser Glück. Also nahm er dem Hiob die Kinder und alle sein Gut, strafte ihn leiblich; dennoch versah er sich Gutes zu GOTT [Cap. 13,

\quellseite{541}
15.]. Darum mußte er genug haben, denn zuletzt ward ihm wieder zweifältig alles, das er verloren hatte [Hiob 42, 10.].

99. Das sei kürzlich gesagt von dem Verstande des ersten Gebots, darin er die Juden ermahnt, daß sie gedenken, daß sie einen solchen GOtt haben, der sie aus Egypten geführt habe, aus dem Diensthause; damit er beweiset, daß er mächtig sei, sie vor allen Feinden zu behüten.

\subsection{Epilogus oder Summarium}

100. Also habt ihr das erste Gebot, das zweierlei Verstand oder Brauch hat. Zum ersten, den äußerlichen, groben Verstand, als, so man Stein und Holz anbetet, wie im Texte klarlich ausgedrückt ist: „Du sollst dir kein Bildniß, noch irgend ein Gleichniß machen, weder deß, das oben im Himmel, noch deß, das unten auf Erden, oder deß, das im Wasser“ etc. Das ist die grobe und tölpische Abgötterei.

101. Der andere Verstand und Mißbrauch, wenn man innerlich nicht mit dem Herzen GOtt vertrauet in einem rechten Glauben, sondern das Herz an andere Dinge hänget, und verläßt sich auf ein Ding, das nicht GOtt ist. Das ist denn die rechte, innwendige Abgötterei, und der rechte, lebendige Teufel. Und wiewohl äußerliche Bilder von Holz, Stein, Gold oder Silber anbeten ein großer Mißbrauch ist, so ist es doch nur ein Kinderspiel, und ein Zeichen der innwendigen Abgötterei des Herzens, darin die ganze Welt steckt und ganz ersoffen ist. Denn es kann niemand GOtt glauben und vertrauen, der Heilige Geist erleuchte denn vorhin das Herz. Es sind wohl viele, die sich enthalten, äußerlich die Bilder anzubeten, es ist aber keiner, der nicht den Teufel im Herzen anbete, dem sein Herz nicht ersoffen sei im Unglauben, daß er Gott nicht vertrauet, ihn nicht für wahrhaftig hält, setzt viel mehr seine Zuversicht auf sein Gut, Weisheit, Vermögen, Frömmigkeit und Heiligkeit, denn auf GOttes Güte und Barmherzigkeit. Das ist denn die rechte Abgötterei.

102. Also ist das erste, höchste und alleredelste gute Werk in diesem Gebot der Glaube an\footnote{In den alten Ausgaben: in.} Gott. Denn in diesem Werke müssen alle Werke geben, und ihrer Gutheit Einfluß, gleich wie ein Leben, von ihm empfangen. Und wo der Glaube nicht ist, so ist den Werken der Kopf ab, und alle ihr Leben und Gutes ist nichts; wie Paulus lehrt Röm. 14, 23.: „Alles, was nicht aus oder im Glauben geschieht, das ist Sünde.“ Von dem Glauben, und keinem andern Werke, haben wir den Namen, daß wir Christgläubige heißen. Denn alle anderen Werke mag ein Heide, Jude, Türke, Sünder auch thun; aber GOtt vertrauen festiglich, ist nicht möglich denn einem Christen, mit GOttes Gnade erleuchtet.

103. In diesem Glauben werden auch alle Werke gleich, und wird eines wie das andere ohne allen Unterschied, sie seien groß, klein, kurz, lang, viel oder wenig. Denn die Werke sind nicht von ihretwegen, sondern von des Glaubens wegen GOtt angenehm, welcher einig und ohne Unterschied in allen und jeglichen Werken ist, wirkt und thut sie alle, wie viel und unterschiedlich sie immer sind. Gleichwie alle Gliedmaßen von dem Haupt ihr Leben, Wirken, und den Namen haben, und ohne das Haupt kann kein Glied leben, wirken, oder einen Namen haben.

104. Dieser Glaube steht aber nicht allein darin, daß ich zum ersten glaube, es gefalle GOtt wohl, daß ich esse, trinke, schlafe oder wache, 1 Cor. 10, 31., oder andere dergleichen geringe Werke thue, sondern auch, zum andern,\footnote{„zum andern“ fehlt in der Wittenberger und in der Jenaer.} wenn mir es übel geht an Leib, Gut, Ehre etc., daß ich dennoch glaube, GOtt meine es väterlich mit mir, stehe bei mir, und werde mich nicht verlassen. Solches zu glauben ist die größte Kunst, nämlich, zu GOtt, der sich zornig stellt nach unserm Sinn und Verstand, eine gute Zuversicht zu haben, und Besseres sich zu ihm versehen, denn sich’s empfinden läßt; denn hier ist er verborgen, wie die Braut spricht im Hohenliede [Cap. 2, 9.]: „Siehe, er steht hinter der Wand, und sieht durch die Fenster, und gucket durch das Gitter.“

105. Zum dritten, ist der höchste Grad des Glaubens, wenn GOtt nicht mit zeitlichem Leiden, sondern mit dem Tode, Hölle und Sünde das Gewissen trifft und straft, und gleich Gnade und Barmherzigkeit absagt, als wollte er ewig verdammen und zürnen, welches wenig Menschen erfahren, wie David Ps. 6, 2. klagt: „HErr, strafe mich nicht in deinem Grimme“ etc. Hier zu glauben, daß GOtt einen gnädigen Wohlgefallen über uns habe, ist das höchste Werk,
\quellseite{542}
das geschehen mag von und in der Creatur. Diesen guten Willen und Wohlgefallen, darauf unsere Zuversicht steht, haben die Engel vom Himmel verkündigt, da sie sangen in der Christnacht [Luc. 2, 14.]: Gloria in excelsis Deo, „Ehre sei GOtt in der Höhe, Friede auf Erden, und ein gnädiges Wohlgefallen den Menschen.“

106. So ist nun das die Meinung dieses Gebots: „Du sollst nicht andere Götter haben“: weil ich allein dein GOtt bin, sollst du zu mir allein deine ganze Zuversicht, Vertrauen und Glauben setzen, und auf niemand anders. Denn das heißt nicht einen GOtt haben, so du äußerlich mit dem Munde GOtt nennst, oder mit den Knieen und Gebärden anbetest, sondern, so du herzlich ihm vertraust, und dich alles Gutes zu ihm versiehst, es sei im Leben oder Sterben, in Liebe oder Leid, wie Joh. 4, 23. von rechten Anbetern steht. Und dieser Glaube und Zuversicht des Herzens ist die wahrhaftige Erfüllung dieses ersten Gebots, ohne welche sonst kein Werk ist, das diesem Gebot möge genug thun.

107. Und wie dieses Gebot das allererste, höchste und beste ist, aus welchem die andern alle fließen, in ihm gehen, und nach ihm gerichtet werden, also ist auch sein Werk (das ist, der Glaube oder Zuversicht zu GOttes Huld) das allererste, höchste, beste Werk, aus welchem alle andern fließen, gehen und gerichtet werden; und andere Werke gegen diesem sind eben, als ob die andern Gebote wären ohne das erste, und kein GOtt wäre. Derhalben spricht wohl St. Augustinus, daß des ersten Gebots Werke glauben, hoffen und lieben sind. Denn solche Zuversicht bringt mit sich Liebe und Hoffnung.

\section{Das andere Gebot der ersten Tafel}

\begin{schriftzitat}
V. 7. Du sollst den Namen des HErrn, deines GOttes, nicht vergeblich führen. Denn der HErr wird den nicht unschuldig halten, der seinen Namen vergeblich führet.
\end{schriftzitat}

\subsection{Summa}

108. Das andere Gebot lehrt, wie sich der Mensch halten soll gegen GOtt äußerlich in Worten vor den Leuten, oder auch innerlich vor ihm selbst, nämlich, daß er GOttes Namen ehre. Denn niemand kann GOtt weder vor ihm selbst, noch vor den Leuten anzeigen nach der göttlichen Natur, sondern bei seinem Namen.

\begin{schriftzitat}
Du sollst den Namen des HErrn etc.
\end{schriftzitat}

109. Im ersten Gebote habt ihr gehört eine Dräuung und Verheißung oder Zusagung, so sie hielten seine Gebote. Hier setzt er allein eine Dräuung und keine Verheißung. Also haben die zwei Gebote vor den andern allen den Zusatz, daß er mit Zorn dräuet, weil darob gehalten haben mehr, denn bei keinem andern. Denn je höher und größer die Gebote sind, je leichter, geringer und schändlicher sie verachtet werden. Die nachgehenden Gebote, als, tödten, stehlen, ehebrechen, und falsch Zeugniß geben, wiewohl sie auch läufftig\footnote{läufftig = gangbar.} sind und mißhandelt werden, so sind sie doch nicht so täglich und gemein, als die zwei ersten Gebote. Denn es sind ihrer viel, die äußerlich nicht todtschlagen, nicht ehebrechen, nicht stehlen, nicht falsch Zeugniß geben; aber hier ist niemand, der nicht im Herzen mit Abgötterei befleckt und besudelt wird, und den Namen GOttes nicht mißbrauche.

110. Darum, weil man es so leichtlich verachtet und in den Wind schlägt, hält GOtt desto höher darüber, setzt eine Dräuung dabei, daß man ja desto mehr darauf halten soll. Denn er will es ja nicht verachtet haben, ja, er will, daß wir mehr Fleiß und Achtung darauf legen, denn auf die andern. Denn die andern werden leichtlich gehalten, wo die zwei ersten recht gefaßt werden. Wir aber kehren das Blatt um, halten am wenigsten, die wir am meisten sollten halten.

111. Das Gebot hat auch zween Mißbräuche, wie das erste. Der erste ist, so man unnütz schwört und flucht bei GOttes Namen, und sein mißbraucht.\footnote{Erlanger: mißbrauchet den Namen Gottes.} Also haben es auch die Juden verstanden. Solches Schwören bei Gottes Namen etc. ist jetzt ganz gemein und in täglichem Brauche. Denn wir haben unter allen Gliedern kein leichtfertiger\footnote{Wittenberger: rechtfertiges.} Glied denn die Zunge, die schnappet daher, achtet gering, daß sie also leichtfertig den Namen GOttes nennt in leichtfertigen Dingen, den sie doch in großen Ehren halten sollte. Darum, weil man ihn so gering achtet, setzt er eine Dräuung hinzu, und spricht: „GOtt wird den nicht unschuldig achten“ etc., es wird ihm nicht wohl gehen, er wird einmal die Strafe sehen, er wird ihm auf die Haube greifen, daß er also schilt und flucht, und seinen
\quellseite{543}
Namen also mißhandelt, er wird die Länge nicht zusehen, daß man ihn also schändet.

112. Diese Sünde hat man fast allein in diesem Gebote gepredigt, und am meisten getrieben, daß wir nicht schwören, fluchen, lügen, trügen, zaubern sollen mit dem Namen GOttes, noch andere Mißbräuche treiben; welches alles grobe, äußerliche Stücke sind und jedermann wohl bekannt. In welchen auch begriffen ist, daß wir wehren sollen den andern, daß sie nicht lügen, trügen, schwören, fluchen, zaubern, und auf andere böse Weise mit GOttes Namen sündigen.

113. Der andere Mißbrauch aber ist also groß, daß der erste ein Kinderspiel gegen dem andern ist; wiewohl der erste greulich und groß genug ist. Dieser aber geht also daher, und ist so subtil und geistlich, daß ihn niemand recht versteht, auch die großen Heiligen; es will ein starker Geist dazu gehören, daß man ihn verstehe und sehe. Ich will des geschweigen, daß man sein ganz los werden soll. Denn die in dem Mißbrauche stecken, achten, sie thun GOtt einen Dienst daran, wenn sie seinen Namen am höchsten schänden und lästern.

114. Wie auch wider das erste Gebot viel sündigen, vornehmlich, die einen Schein haben, daß sie die Frömmisten und Heiligsten sind, werfen ihre eigenen Werke auf, setzen sich an die Statt GOttes, richten mit ihrem eigenen Dünkel eine Abgötterei an, unter welcher sie den Teufel anbeten, achten solches gering und schlagen es in Wind, ja, meinen, sie thun wohl und recht, und haben ein groß Gefallen daran; ja, das viel ärger ist, sie wollen ihr Thun ungetadelt haben, vertheidigen es mit Gewalt; Trotz, der wider sie rede; gedenken nicht, daß GOtt redet, er sei ein starker Eiferer, er lasse ihm nicht im Bart greifen. Und wenn sie es schon lesen, meinen [sie], es gehe sie nicht an, ja, schwören zu GOtt, sie wären die lieben Kinder, die GOtt in dem Schooße sitzen.

115. Also sündigen auch wider dies andere Gebot am meisten und greulichsten, die das größte Ansehen der Frömmigkeit und Heiligkeit haben, nämlich, die da predigen und lehren unter dem Scheine und Namen des Wortes GOttes eitel Teufelslehren und Menschenge-setze; predigen und lehren also nach ihrem Wohlgefallen unter dem Namen GOttes, das ihnen nicht befohlen ist, sind ganz frech und verstocket in ihrem stolzen Sinne, entsetzen sich nicht vor der greulichen Dräuung Gottes: „Der HErr wird den nicht unschuldig halten, der seinen Namen vergeblich führet“, die über alle solche Lästerer gehen wird; ja, meinen noch, sie thun recht, und GOtt einen Gefallen daran, der sie darum auch erhöhen werde. Solche Geister kommen nimmer dahin, daß sie ihre Sünde erkennen, ja, sie achten es für die höchste Gerechtigkeit, lästern also innerlich den Namen GOttes unter einem heiligen Scheine, den die Welt, Vernunft und alle menschliche Weisheit für böse nicht erkennen kann.

116. Der Papst ist in dieser Gotteslästerung ganz ersoffen. Denn er gehet unter dem Namen GOttes, das wider GOtt ist, legt große Bürden auf die Menschen durch seine Gebote; gibt vor, er thue es aus GOttes Befehl, und hat’s ihn doch der leidige Teufel geheißen; macht also seiner Lehre einen Deckel mit dem Namen GOttes: Wir gebieten von wegen des allmächtigen GOttes, des Vaters, und des Sohnes, und des Heiligen Geistes, und aus Gewalt der heiligen Apostel Petri und Pauli etc., wer uns nicht gehorsam ist, der ist GOtt nicht gehorsam. Weiß auch gar hübsch die Sprüche auf sich zu beuten, die Christus hat auf die rechten Prediger gedeutet; als den Luc. 10, 16.: „Wer euch höret, der höret mich, und wer euch verachtet, der verachtet mich.“ Reime dich. Da fürchtet sich denn jedermann, Kaiser, König, Fürsten und Herren, fallen plump hinein,\footnote{So die Jenaer; Wittenberger: fallen hinein, plump; Erlanger: „fallen hinein, plump,“ etc.} der Hölle und dem Teufel zu.

117. Also kann der Teufel die Kunst, daß er sich verstellt zum Engel des Lichts [2 Cor. 11, 14.], und gibt falsche Lehre den falschen Predigern und Rottengeistern ein, unter einem Schein rechter, heilsamer Lehre. Die rufen denn den Namen GOttes aus in der Welt, geben tödlich Ding vor und schreien: GOttes Wort, GOttes Wort! Treiben also unter der gleichenden Farbe ihr böses Gift in die Leute, damit die rechtschaffene Lehre ganz ausgegeilt wird, sind ganz frech, fahren herein ohne alle Scheu,\footnote{Erlanger: ahn allen Scheu.} fürchten weder GOtt noch die Welt.

118. Darum spricht GOtt hier, und dräuet solchen Lästermaulern erschrecklich. Siehe dich vor, ich will dir es nicht zugute halten; das ist,
\quellseite{544}
ich will dich der Tage einst gar greulich strafen und heimsuchen. Es lautet viel heftiger, daß er spricht: Ich will dir es nicht zugute halten; denn daß er spräche: Ich will dich strafen; als könnte er die große Schwere\footnote{Walch und die Erlanger: Größe und Schwere.} der Strafe nicht nennen. Als, so ein Hausvater seinem Sohne oder Knechte bräuete, und spräche: Ich will dir es nicht schenken, so lautet es viel heftiger und zorniger, denn wenn er spräche: Ich will dich treffen; oder: Wenn ich dir es schenkte etc. In solchen Träumungen lautet die negatio viel stärker, denn die affirmatio; die Worte der Träumung lauten gering, aber sie werden einen heftigen, harten und erschrecklichen Nachdruck haben wider [die], die GOttes Namen auf die Weise, wie gesagt ist, lästern.

119. So ist nun der rechte Mißbrauch des Namens GOttes, daß man also die falsche Lehre vertheidigt, und vorgibt, GOtt wolle es also haben, rühmt unverschämt: GOttes Wort, GOttes Wort, so es doch GOtt nicht befohlen hat. Und dieser Mißbrauch geht durch und durch in der ganzen Welt, nämlich bei denen, die den Schein und den Namen haben, daß sie gelehrte und fromme Leute sind, und geben durch denselben Schein vor, ihre Lehre sei die rechte Lehre. Darum ist das gemeine Sprüchwort wohl wahr: In GOttes Namen hebt sich alles Unglück an. Es ist kein größer Unglück in die Welt gekommen, denn daß man unter GOttes Namen alle Abgötterei und falsche Lehre verbirgt, und so einen guten Schein führt, daß man es nicht sehen kann.

120. In einem solchen hübschen Scheine sind auch die Bauern aufgestanden. Was hat sie bewegt? Was gaben sie vor? Den Namen und die Ehre GOttes. Wir wollen, sprachen sie, das Evangelium vertheidigen, beschützen und handhaben. Als wäre GOtt zu schwach, sein Evangelium zu vertheidigen, müßte solche Gehülfen haben, sein Evangelium handzuhaben. Nahmen also Spieße und Schwerter zuhanden, meinten, sie hätten GOttes Wort, meinten, sie könnten auch mit Recht die gottlosen Fürsten todtschlagen, weil solches im Alten Testamente dem jüdischen Volke befohlen war. Also müßte der Name GOttes ihr Schanddeckel sein. Wer richtete es an? Der Teufel führte die Mordpropheten her, die überredeten sie, daß sie meinten, sie thäten recht daran; denen mögen sie danken. Aber Gott vergaß auch seiner Worte nicht, kam bald mit der Strafe, wie er geredet hat: „Ich will den nicht unschuldig achten, der meinen Namen mißbraucht.“ Denn er kann nicht leiden, daß man also mit seinem Namen scherze. Er kam plötzlich über sie, und schlug sie über die Köpfe; sie hoben das Spiel wider GOtt an, darum kam das Urtheil über sie. Ihr habt das Exempel erlebt, und selber gesehen, denkt daran, und vergeßet sein nicht.

121. Also geschieht, und ist geschehen allen, wie denen, die den Namen GOttes mißbrauchen. Also ging es den Ketzern Ario, Manichäo, Pelagio,\footnote{In der Erlanger: „Ariano“ und „Pelagiano“.} und allen, die unter einem solchen Scheine GOttes Namen haben mißbraucht. GOtt hat sie wohl eine Weil lassen lästern, seinen Namen mißhandeln, aber darnach plötzlich ist er über sie gekommen, und sie zerschmettert, daß ihr Name in der ganzen Welt jetzt stinkt.

122. Also auch der Papst hat Gottes Namen gelästert; wiewohl er es eine lange Zeit getrieben hat, so ist doch das Urtheil zuletzt über ihn auch gekommen, daß jetzt bei den rechten Christen nichts Verachteteres ist, denn der Papst mit allen seinen Mönchen und Pfaffen; es stinkt wie ein Widehopfennest bei ihnen, und wird auch je länger je mehr stinken, auch bei denen, die ihn jetzund noch hoch achten. Denn es ist ja wahr, und bleibt wahr, wer den Namen GOttes unehrt und schändet, der muß auch zu Schanden und ungeehrt werden.

123. Das soll nun uns, die wir recht lehren, und den Namen GOttes recht erkennen, heiligen und preisen, nicht lästern und unheiligen, ein großer Trost sein. Und ob wir gleichwohl verachtet sind, sind wir doch gewiß unserer Lehre. Aber die Secten und Rottengeister fahren daher, berühmen sich, sie haben den rechten Geist, schweben empor, und ist doch nichts denn lauter Gotteslästerung; plaudern herein, GOttes Wort, GOttes Wort; aber sie müssen herunter. Denn, je höher sie erhoben sind, je mehr sie zu Schanden müssen werden; und wenn es nicht geschieht, so wollen wir hier den Text auskratzen, und muß erfolgen sein, was hier steht.

124. GOtt sieht wohl eine Weile zu; er ist aber ein Eiferer, er ist neidisch, kann es die Länge nicht erdulden; darum schmeißt er endlich drein, und wahrlich grob. Es ist auch verdrießlich einem
\quellseite{545}
solchen, nämlich GOtt, seinen Namen zu lästern und schänden. Wenn einer einen Bürger in seinem Hause überfiele, lästerte und schmählete ihn, und wollte also Gewalt üben, wer wollte das gerne leiden? Wenn man einen solchen ergriffe an frischer That, und würde drob erschlagen, wer wollte sprechen, daß einem solchen Unrecht geschehen wäre? Denn ein jeglicher hat gerne in seinem Hause Frieden, und wer ihm denselben nehmen will, der hat den Hals verfallen.

125. Können wir nun unter einander nicht leiden, daß einer den andern schelte und lästere, wie viel weniger wird es GOtt leiden, daß man seinen Namen also im Maul unnütz führe, und ihn für einen Götzen halte. Wie könne er dazu? Es steht ihm auch nicht zu leiden. Er ist ein rechter GOtt, darum will er auch seine Ehre retten. Es soll sich auch der niemand unterwinden; thut er es aber, so wird er zu Boden gestoßen werden. Also ist das Gebot auch uns natürlich in das Herz geschrieben, daß es ja nicht billig ist, daß man GOtt seinen Namen also schände.

126. So habt ihr nun zwei Mißbräuche dieses Gebots; den ersten äußerlich und grob, daß man nicht fluche, schelte. Der ist nun also grob, daß ihn jedermann versteht, daß ihn auch weltliche Obrigkeit straft. Wenn sie es auch nur thäte! Der andere ist subtiler, der einen Schein hat, als sei er GOttes Ehre; wie denn die falschen Prediger, und die ihre Lehre von ihnen annehmen, GOttes Namen also mißbrauchen. Die beide straft GOtt; fahren beide, Meister und Schüler, zum Teufel.

127. Wie wird nun das andere Gebot erfüllt? Das erste Gebot, habt ihr gehört, wird erfüllt durch einen wahren, rechten, festen Glauben, und durch ein aufrichtiges Herz und gute Zuversicht gegen GOtt; also, daß sich der Mensch zu GOtt versieht, er werde ihn erhalten, helfen, rathen, und alles Gutes geben, zeitlich und ewiglich. So ist es recht gehalten. Wo aber diese Zuversicht nicht ist, da wird es nicht gehalten.

128. Also, in diesem Gebote, wird der Name GOttes recht geführt, wenn man das Wort Gottes recht\footnote{Wittenberger: nicht.} predigt, und wird von den Zuhörern recht angenommen. Und wiederum der Name GOttes wird gelästert, wenn die Prediger nicht recht predigen, sondern verführen das Volk, doch unter dem Schein göttliches Worts und Namens. Wenn nun der Prediger gewiß ist, daß er das Wort GOttes habe und recht predige, unterwindet sich keiner Ehre, die GOtt allein zusteht, sondern richtet fleißig aus, was\footnote{So die Jenaer. Wittenberger: wie es; Erlanger: wie.} ihm befohlen ist, leugt und treugt nicht, trägt das Wort klar vor, ohne alle Falschheit, der erfüllt das Gebot. Item, die dem Prediger gehorchen, nehmen das Wort an, als ob es GOtt selber redte, fassen es mit Herzen etc., das ist denn der rechte Brauch des Namens GOttes, nämlich, den Namen GOttes predigen, an ihn glauben, ihn bekennen, loben, anrufen, und darob Leib und Leben lassen; denn dies alles folgt aus dem ersten.

129. Also ist das größte und höchste Werk nach dem Glauben, GOttes Namen loben, ehren, preisen, predigen, singen, und auf allerlei Weise erheben und groß machen. Und wiewohl das wahr ist, wie droben [§ 103] gesagt ist, daß kein Unterschied ist unter den Werken, wo der Glaube ist und wirkt, so ist doch solches zu verstehen, wenn die Werke gegen dem Glauben und seinen Werken geachtet werden, so ist kein Unterschied unter ihnen. Wenn sie aber unter einander gemessen werden, so ist ein Unterschied zwischen ihnen, und ist eines höher denn das andere. Gleichwie am Leichnam die Gliedmaßen gegen der Gesundheit keinen Unterschied haben, und die Gesundheit in einem gleich sowohl wirkt, als in dem andern, so sind doch der Gliedmaßen Werke unterschieden, und ist eines höher, edler, nützlicher denn das andere. Also auch hier, GOttes Ehre und Namen preisen, ist besser denn die folgenden Werke der andern Gebote, und muß doch in demselben Glauben gehen, da alle andern in gehen.

130. Das ist auch wohl zu merken bei diesem und andern Geboten, daß, gleichwie das erste Gebot verbeut, wir sollen keine andern Götter haben, also gehet es auch herunter, wir sollen Einem, das ist, dem rechten GOtt, durch\footnote{Erlanger: durch einem.} festen Glauben vertrauen, Zuversicht, Hoffnung und Liebe zu ihm haben; welches die Werke sind, damit man einen GOtt haben, ehren und behalten mag. Also wird auch im andern Gebote verboten: Wir sollen GOttes Namen nicht unnützlich brauchen. Doch will das nicht genug
\quellseite{546}
sein, sondern wird darunter auch geboten: Wir sollen seinen Namen nützlich brauchen, das ist, ehren, anrufen, preisen, predigen und loben. Und zwar ist es unmöglich, daß GOttes Name nicht sollte verunehrt werden, wo er nicht recht geehrt wird. Denn, ob er schon mit dem Munde, Kniebeugen, Küssen, oder andern Geberden wird geehrt, wo aber solche Ehrerbietung nicht im Herzen durch den Glauben an GOtt geschieht, ist es nur ein Schein und Gleisnerei.

131. Dieses Gebot wird auch erfüllt, wenn man den Namen GOttes anruft in aller Noth und Anfechtung. Denn darum läßt uns GOtt noch in vielen Nöthen leiden, mit Furcht des Todes, mit bösen Lüsten angefochten werden, auf daß er uns durch solche Stücke dringe und Ursache gebe, zu ihm zu laufen, zu schreien, und seinen heiligen Namen anzurufen, wie der 50. Psalm, V. 15., spricht: „Rufe mich an in der Zeit der Noth, so will ich dich erretten, so sollst du mich preisen.“ Denn ein Opfer des Lobes will ich haben, und dasselbe ist der Weg, dadurch du magst kommen zu der Seligkeit. Denn durch solche Anfechtung wird der Mensch gewahr und erfährt, was GOttes Name sei, und wie mächtig er ist zu helfen allen denen, die ihn anrufen, und wächst dadurch fast sehr der Glaube und Zuversicht zu GOtt. Davon sagt David Psalm 54, 8. 9.: „Du hast mich erlöset von aller Noth; darum will ich deinem Namen nachsagen,\footnote{Die Jenaer hat hier die Lesart unserer Bibel: „danken“ und „tröstlich“; beides notirt die Erlanger als Bemerkungen Walchs.} und bekennen, daß er so gut ist.“ Psalm 91, 14. spricht GOtt: „Ich will ihn erlösen, darum, daß er auf mich hoffet; ich will ihm helfen, darum, daß er meinen Namen kennet.“

132. Das ist aber das größte und allerschwerste Werk dieses Gebots, daß man den heiligen Namen GOttes schütze wider alle, die sein geistlicher Weise mißbrauchen, und dazu ihn ausbreiten unter allen Menschen. Denn es ist nicht genug, daß ich für mich selbst und in mir selbst [den] göttlichen Namen lobe und anrufe in Glück und Unglück, ich muß auch heraus fahren, und um GOttes Ehre und Namen willen auf mich laden Feindschaft aller Menschen, wie Christus [Matth. 24, 9.] spricht zu seinen Jüngern: „Es werden euch feind sein, um meines Namens willen, alle Menschen.“ Hier müssen mir Vater und Mutter und die besten Freunde erzürnen, wie Christus selbst spricht im Evangelio Luc. 14, 26.: „So jemand zu mir kommt, und hasset nicht seinen Vater, Mutter, Weib, Kind, Bruder, Schwester, auch dazu sein eigen Leben, der kann nicht mein Jünger sein“; und Matth. 10, 35. 36.: „Ich bin kommen, den Menschen zu erregen wider seinen Vater, und die Tochter wider ihre Mutter, und die Schnur wider ihre Schwieger, und des Menschen Feinde werden seine eigenen Hausgenossen sein.“

133. Hier müssen wir auch den Namen haben, daß wir wider die Obrigkeit geistlich und weltlich streben, und ungehorsam gehalten werden. Hier müssen wir die Gelehrten, Heiligen, Reichen, Gewaltigen, und alles, was nur etwas in der Welt ist, wider uns erwecken; das heißt denn, GOttes Freund, und aller Welt Feind. Und wiewohl das sonderlich zu thun schuldig sind, denen GOttes Wort zu predigen befohlen ist, so ist doch auch ein jeglicher Christ dazu verbunden, wo es die Zeit und Noth erfordert.

134. Wenn nun ein Mensch das Wort GOttes, das Evangelium, annimmt, so gedenket er nur nicht anders, denn daß er in derselben Stunde trete in Gefahr aller seiner Güter, Hauses, Hofs, Acker und Wiesen, Weibes, Kinder, Vaters und Mutter, auch seines eigenen Lebens. Wenn ihm denn Gefahr und Unglück zu Hause kommt, so ist es ihm desto leichter, und gedenkt, ich habe es vorhin wohl gewußt, daß [es] also gehen würde. Da gehören denn die Sprüche her, Matth. 10, 24. 25.: „Der Jünger ist nicht über den Meister“ etc., item: „Haben sie den Hausvater Beelzebub geheißen, wie viel mehr werden sie seine Hausgenossen also heißen“; item, Jac. 4, 4.: „Wer der Welt Freund sein will, der wird GOttes Feind sein“; item, Joh. 15, 19.: „Sie werden euch hassen, denn ihr seid nicht von der Welt“; item, Matth. 10, 34.: „Ich bin nicht kommen, Frieden zu senden auf Erden, sondern das Schwert.“

\subsection{Vom Eide}

135. Hier gebühret sich auch zu reden vom Eide. Denn durch einen rechten Eid wird auch äußerlich GOttes Name recht gebraucht und geehrt, und durch einen falschen Eid wird er ungeehrt.

136. Ist nun hier eine Frage: ob man auch möge schwören bei dem Namen GOttes im
\quellseite{547}
neuen Testamente? Denn im Alten Testamente [5 Mos. 6, 13. 10, 20.] war ihnen geboten, wenn sie wollten schwören, daß sie bei dem Namen GOttes mußten schwören; Christus aber, Matth. 5, 34–37., verbeut, „man solle ganz und gar nicht schwören, weder bei dem Himmel, noch bei der Erde, noch bei der Stadt Jerusalem, noch bei dem Haupte etc., sondern unsere Rede soll Ja, Ja sein“, wenn es Ja ist, „und Nein, Nein“, wenn es Nein ist. Es ist oft davon gehandelt worden, doch um der Jungen\footnote{Wittenberger und Jenaer: „Jüden“.} und Einfältigen willen müssen wir es wiederholen.

137. Es geht mit dem Schwören zu, wie mit dem Todtschlagen. Todtschlagen ist hart verboten, wie wir wissen, in den zehn Geboten; wenn es aber geschieht aus Befehl des Worts GOttes, so ist es recht und gut. Also, wenn ein Richter einen Mörder oder Todtschläger den Kopf läßt abschlagen, oder auf ein Rad legen, der thut wohl und recht daran, thut auch nicht wider das fünfte Gebot GOttes, denn GOtt hat es also geordnet und befohlen, die Bösen zu strafen, um eines gemeinen Friedens willen. Also haben todgeschlagen die lieben Heiligen Mosis, Samuel, David etc. aus dem Befehle GOttes, und nicht aus eigenem Lust und Wahl. Wer aber das Schwert nimmt ohne GOttes Befehl, der soll mit dem Schwert gerichtet werden [Matth. 26, 52.]. Wenn ich nun tödte aus Zorn oder Bosheit, ich sei Richter oder nicht, so thue ich wider das Gebot GOttes: Du sollst nicht tödten. Wenn aber der Richter seines Amts halben einen köpfen, hängen, ersäufen, oder verbrennen läßt, so thut es der Richter nicht, sondern das Urtheil GOttes, ja, GOtt selber. Denn wer wollte anders Richter sein, und ein solch Amt führen, wenn einer nicht gewiß wäre GOttes Befehls, und einen solchen Trost hätte, daß es GOttes Wort und Befehl ausrichtet, nicht er. Wo das nicht wäre, müßte das Gewissen eines Richters verzweifeln.

138. Aber also kann der Richter seine Augen aufrichten zu GOtt, und sprechen: Lieber HErr, ich tödte den Menschen nicht gerne, denn vor deinen Augen bin ich nicht frömmer denn er; weil es aber dein Wille ist und Befehl, daß die öffentlichen Uebelthäter gestraft sollen werden, daß man vor ihnen Frieden habe, so bin ich gewiß und sicher, daß ich recht daran thue, und dir Gehorsam darinnen leiste, daß ich den würgen lasse; ja, HErr, du thust es selbst, denn es ist ja dein Befehl. So wird ihm GOtt das Recht zufällen, und sprechen: Du hast recht gehandelt, bist meinem Befehl gehorsam worden, denn ich habe es ja dir befohlen. Den Trost muß ein Richter behalten, will er anders vor GOtt bestehen.

139. Wie nun der Todtschlag, Köpfen, Radbrechen, und dergleichen Richters Werke, nicht aus eigenem Muthwillen, Lust und Wahl geschehen sollen, also geht es auch zu mit dem Eide. Kurzum: der Eid und das Schwören ist verboten, wie das Todtschlagen; wenn aber die Obrigkeit einen Eid von mir erfordert, soll ich schwören. Man findet in der Schrift oft, daß die Propheten geschworen haben; item, Christus und die Apostel, sonderlich Paulus, schwören oft. Warum? Hat er doch das Schwören verboten? Christus hat seine Lehre aus Befehl seines Vaters in die Welt gebracht, und zu bestätigen seine Lehre, die der Welt unbekannt war, und zu erregen und bekräftigen die schwachen Herzen der Menschen, hat er den Eid gebraucht. Also hat auch hin und wieder geschworen Sanct Paul.

140. Also ist nun der Eid verboten, wenn er geschieht aus eigenem Vornehmen, Muthwillen, Leichtfertigkeit und böser Lust. Wenn er aber aus Noth geschieht, ist er nicht verboten, ist auch nicht unrecht. Dann geschieht er aber aus Noth, wenn die Obrigkeit einen Eid erfordert vor Gericht etc. Wie auch geschieht, wenn man den Fürsten und Herren huldiget und schwört; und ist recht. Also kann man auch schwören aus Pflicht der Liebe, nämlich, so ein Nächster dem andern schwört, ihm Glauben zu halten, so es der andere begehrt; nicht (wie sie gesagt haben) in geistlichen Sachen allein, sondern in allen Sachen gegen dem Nächsten, ihn sicher zu machen, wenn er es begehrt, und will dir sonst nicht Glauben geben; oder daß man sich mit ihm versöhne, wie die Epistel an die Hebräer anzeigt, Cap. 6, 16., da sie also spricht: „Der Eid ist das Ende alles Haders, zur Befestigung unter ihnen.“ Also hat Joseph geschworen bei dem Namen Pharaonis, im ersten Buch Mosis Cap. 42, 15. Also schworen die Propheten bei dem Leben der Könige. Denn solche Schwüre sind nicht verboten, wie es die tollen Sophisten ganz
\quellseite{548}
enge gespannt haben, haben selber nicht gewußt, wovon sie redeten. Darum ist das der Beschluß: Wenn es die Obrigkeit heißt, und dein Nächster begehrt es, und nöthigt ihn, so schwöre immerhin, es ist recht; denn du setzest ihm GOtt zum Pfande, und gibst GOtt die Ehre, daß er treu sei, und werde solch Zeugniß sehen. Es ist die höchste Ehre, die man GOtt geben kann, daß man ihn für treu und wahrhaftig hält.

141. Das ist nun vom äußerlichen Eide zu halten gesagt, wenn er geschieht aus Befehl der Obrigkeit. Wenn man aber schwört ohne Gebot der Obrigkeit, und ohne die Pflicht brüderlicher Liebe, sondern aus Leichtfertigkeit und eigenem Muthwillen, so ist es verboten und unrecht, und wider das andere Gebot, das da gebeut, daß man den Namen GOttes nicht unnütz soll brauchen. Da soll man denn (wie Christus Matth. 5, 34. 35. sagt) „nicht schwören, weder bei dem Himmel, noch Erde, noch Jerusalem, noch unserm Haupte“, so lange bis es die Noth erfordert. Dann erfordert es aber die Noth, wenn es dein Nächster haben, und ohne den Eid nicht glauben will, und es die Obrigkeit gebeut. Denn so ist es gleich, als ob man GOtt anruft, daß er der Sache wolle Zeugnüß geben. Wenn es nun nicht unrecht ist, daß man GOtt über einem Trunk Bier oder Bissen Brods anruft, so ist es auch nicht unrecht, daß man den Namen GOttes dem Nächsten zugute braucht.

142. Das sei nun von den zweien Bräuchen dieses Gebots gesagt. Der äußerliche Mißbrauch ist grob und gering; der innerliche aber ist so hoch und groß,\footnote{Walch und die Erlanger: grob.} daß über alle Maße ist. Als nämlich, mit falscher Lehre und falschem Gottesdienste wird Gottes Name umgekehrt und gelästert, durch rechte Lehre aber, die das Wort GOttes recht und klar lehrt, wird GOttes Name gelobt und gepreiset. Daher gehören die Sprüche aus dem Psalter, Ps. 7, 18.: „Ich will loben“, preisen und ehren „den Namen GOttes des HErrn.“ Item [Ps. 99, 3.]: „Wir danken deinem großen und wunderbaren Namen, der da heilig ist.“ Item, Ps. 100, 4. 5.: „Danket und lobet seinen Namen, denn der HErr ist gut.“ Item, Ps. 113, 1–3.: „Lobet den Namen des HErrn. Gelobet sei sein Name, von nun an bis in Ewigkeit. Vom Aufgange der Sonne, bis zu ihrem Niedergange, sei gelobet der Name des HErrn.“ Und Salomon in seinen Sprüchen [Cap. 18, 10.]: „Der Name des HErrn ist ein festes Schloß, der Gerechte läuft dahin und wird beschirmet.“

143. Solch Lob soll aber nicht geschehen allein mit dem Munde; denn also loben Mönche, Pfaffen und Nonnen auch GOttes Namen, sondern, daß man allein GOttes Lob und Preis verkündige und predige, daß er uns nicht um unserer guten Werke, Verdienste und Frömmigkeit willen selig mache, sondern allein aus seiner Gnade und Barmherzigkeit. Also loben ihn nicht unsere Geistlichen. Und daß man dasselbe nicht allein predige und höre, sondern auch bekenne vor der Welt und dem Teufel, und darüber lasse Leib und Leben. Ach HErr, ihrer sind gar wenig, die also GOttes Namen loben, bekennen und preisen.

\section{Das dritte Gebot der ersten Tafel}

\begin{schriftzitat}
V. 8–11. Gedenke des Sabbathages, daß du ihn heiligest. Sechs Tage sollst du arbeiten, und alle deine Werke schaffen, aber am siebenten Tage ist der Sabbath des HErrn, deines GOttes, da sollst du kein Geschäft thun, noch dein Sohn, noch deine Tochter, noch dein Knecht, noch deine Magd, noch dein Vieh, noch dein Fremdlinger, der in deiner Stadt Thor ist. Denn sechs Tage hat der HErr Himmel und Erde gemacht, und das Meer und alles, was darinnen ist, und ruhte am siebenten Tage. Darum segnete er den Sabbathag und heiligte ihn.
\end{schriftzitat}

\subsection{Summa dieses Gebots}

144. Das dritte Gebot lehrt, wie sich der Mensch halten soll gegen GOtt äußerlich in Werken, das ist, in Gottesdiensten.

\begin{schriftzitat}
V. 8. Gedenke des Sabbathages, daß du ihn heiligest. Sechs Tage sollst du arbeiten etc.
\end{schriftzitat}

145. St. Paulus und das ganze Neue Testament haben den Sabbath der Juden aufgehoben, auf daß man ja greifen kann, daß der Sabbath die Juden allein angeht. Darum ist nicht vonnöthen, daß die Heiden den Sabbath halten, wiewohl es ein großes und gestrenges Gebot war bei den Juden. Die Propheten haben es auch angezeigt, wie dieser Sabbath sollte aufgehoben werden. Jesaias am letzten Capitel
\quellseite{549}
, V. 23., spricht: Wenn der Messias kommen wird, so wird eine solche Zeit sein, daß ein Neumond an andern, ein Sabbath am andern sein wird. Als wollte er sprechen: Es wird alle Tage Sabbath, und alle Tage Neumond sein. Also, im neuen Testament liegt der Sabbath danieder, nach der groben äußerlichen Weise. Denn dies Gebot hat auch zweierlei Verstand, wie die andern Gebote, einen äußerlichen und einen innerlichen oder geistlichen.

146. Es ist im neuen Testamente bei den Christen alle Tag Feiertag, und sind alle Tage frei. Darum spricht Christus: „Des Menschen Sohn ist ein HErr auch über den Sabbath“, Matth. 12, 8. Darum Paulus hin und wieder ermahnt die Christen, daß sie sich an keinen Tag lassen binden [Gal. 4, 10. 11.]: „Ihr haltet Tage, und Monden, und Feste, und Jahrzeiten. Ich fürchte euer, daß ich nicht vielleicht umsonst habe an euch gearbeitet.“ Item, zum Colossiern [Cap. 2, 16. 17.] noch klärlicher: „So lasset nun euch niemand Gewissen machen über Speise oder Trank, oder über eines Theils Tagen“, nämlich, den Feiertagen, „oder neuen Monden, oder Sabbather, welches ist der Schatten von dem, das zukünftig war“ etc.

147. Wiewohl nun der Sabbath aufgehoben ist, und die Gewissen frei davon sind, so ist es doch gut, und auch vonnöthen, daß man einen sonderlichen Tag in der Woche halte, um das Wort Gottes\footnote{So von uns herübergenommen aus der alten Ausgabe Walchs. Jenaer und Erlanger: „um das Wort GOttes willen“; Wittenberger: „um des Worts Gottes willen“.} daran zu handeln, hören und lernen. Denn jedermann kanns alle Tage nicht gewarten. Auch fordert es die Natur, daß man in der Woche Einen Tag stille halte, und enthalte sich von der Arbeit, beide, Menschen und Vieh.

148. Wer aber nun ein nöthig Gebot will aus dem Sabbath machen, als ein Werk von GOtt erfordert, der muß den Sonnabend halten, und nicht den Sonntag; denn der Sonnabend ist den Juden geboten, und nicht der Sonntag. Die Christen aber haben bisher den Sonntag, und nicht den Sonnabend gehalten, um deßwillen, daß am Sonntage Christus auferstanden ist. Das ist nun ein gewiß Zeichen, daß uns der Sabbath nicht mehr angeht, ja, der ganze Mose; sonst müßten wir den Sonnabend halten; und ist eine große und starke Bewegung, daß der Sabbath aufgehoben ist. Denn wir finden durch das ganze Neue Testament an keinem Orte, daß der Sabbath uns Christen geboten sei zu feiern.

149. Warum hält man denn den Sonntag auch bei den Christen? Wiewohl alle Tage frei sind, und ist einer wie der andere, so ist es doch nütze und gut, ja, sehr vonnöthen, daß man an Einem Tage Feier halte, es sei am Sabbath, Sonntag, oder an einem andern Tage. Denn GOtt will die Welt äußerlich führen und friedlich regieren; darum hat er sechs Tage zur Arbeit gegeben, am siebenten Tage aber sollen Knechte, Taglöhner und allerlei Arbeiter, ja, auch Pferde, Ochsen und ander arbeitsam Vieh Ruhe haben, wie dies Gebot lautet, auf daß sie sich mit Ruhe wieder erholen. Und sonderlich, daß die, so zu andern Zeiten nicht der Ruhe haben, am Feiertage die Predigt hören, und GOtt dadurch lernen erkennen. Und um solcher Ursache, nämlich, der Liebe und Nothdurft willen, ist der Sonntag geblieben, nicht um des Gebots Mosis,\footnote{Erlanger: willen.} sondern um unserer Noth willen, daß wir ruhen, und das Wort Gottes lernen.

150. Unsere Papisten haben es aber, leider, dahin gebracht, daß wir härter ob unsern Feiertagen gehalten haben, haben ihrer auch mehr gehabt denn die Juden. Es liegt aber nichts daran, wir feiern oder nicht, die Gewissen sind frei; wer nicht will feiern, der arbeite immerhin; wir wollen ihn nicht schelten noch verjagen. Will er aber fromm sein, GOtt lernen kennen, so brauche er den Feiertag um der Ursache willen, daß er das Wort GOttes höre. Will er aber gar frech und ruchlos sein, so fahre er hin zum Teufel, als jene, die nichts anders an Feiertagen ausrichten, denn Fressen, Saufen, Tag und Nacht schwelgen.

151. Das ist der grobe und äußerliche Verstand dieses Gebots, welcher uns nicht bindet, sondern ist uns frei, steht in unserer Macht und Willkür, ob wir wollen\footnote{Wittenberger: sollen.} feiern oder nicht. Wenn es meinem Nächsten noth ist, so lasse ich den Feiertag fallen und diene dem Nächsten; wenn er aber Unterweisung aus dem Worte GOttes bedarf, so entschlage ich mich der Arbeit. Darum hat der Papst nicht Macht, weder Sonntage noch
\quellseite{550}
andere Tage zu feiern, [als] nöthig zu gebieten. Darum feiern wir aber, denn es ist uns noth, GOttes Wort zu lernen. Daß der Sabbath aber frei sei, zeigt Christus hin und wieder im Evangelio an, und da er den Sabbath auslegt.

\subsection{Vom geistlichen Sabbath}

152. Hier muß man wissen, daß Sabbath auf hebräisch heißt Feier\footnote{Wittenberger: frey.} oder Ruhe, darum, daß GOtt am siebenten Tage ruhte, und aufhörte von allen seinen Werken, die er geschaffen hatte, 1 Mos. 2, 2. Darum gebot er auch den Juden, daß sie den siebenten Tag sollten feiern, und aufhören von allen ihren Werken; und derselbe Sabbath ist nun uns in den Sonntag verwandelt, und die andern Tage heißen Werktage. Der Sonntag aber heißt Ruhetag oder Feiertag oder Heiligtag. Es ist wohl verdeutscht Feiertag, darum, daß wir daran feiern und müßig stehen. Die Juden nannten ihn Sabbath, wir aber Feiertag.

153. Darnach ist es auch in die Gewohnheit kommen, daß man ihn auch nennt Heiligtag; und ist hier aus dem Text genommen, so er spricht: Du sollt den Sabbath heiligen, das ist, du sollt ihn für heilig halten, laß ihn dir heilig sein. Das war nun bei den Juden eine grobe, äußerliche Heiligkeit, daß sie nicht daran arbeiteten, aßen und tranken anders, kleideten sich anders. Das ist nun vor GOtt eine schlechte Heiligkeit, daß ich dem Sabbath zu Ehren andere Kleider, Speise und Trank brauche etc.

154. Der geistliche Feiertag aber ist der rechte Sabbath, welcher gehalten wird, wenn das Herz den rechten Sabbath feiert; welches ist das höchste und recht geistliche Werk dieses Gebotes, welches begreift die ganze Natur des Menschen.

155. Diesen Sabbath hat uns Christus im Grabe vorgebildet, da hielt er den Sabbath recht, da liegt er in der Ruhe und Feier, enthält sich von allen Werken. Er sieht nicht, er hört nicht, er schläft nicht, er wacht nicht, er ißet nicht, er trinkt nicht, er bauet nicht, regt weder Zunge noch Adern, weder Hand noch Fuß, er steht nicht, so geht er nicht.

156. Welcher nun den rechten geistlichen Sabbath will halten, muß mit Christo ganz todt sein. Es wird ihn auch niemand recht halten, er sei denn recht todt. Wir sahen aber den rechten Sabbath hier an zu halten, wenn unser alter Adam aufhört von allen seinen Werken, Vernunft, Willen, Begierden, Lust, welches alles todt, und auf soll hören im rechten Sabbath. Und alles, das in uns ist, soll göttlich sein, wie Paulus spricht Gal. 2, 19. 20.: „Ich bin durch das Gesetz dem Gesetze gestorben, auf daß ich GOtt lebe. Ich bin mit Christo gekreuziget. Ich lebe, aber doch nun nicht ich, sondern Christus lebet in mir. Denn was ich lebe im Fleisch, das lebe ich im Glauben des Sohnes Gottes“ etc. Das ist eine rechte Feier der Christen. Item, Gal. 6, 14.: „Durch Christum ist mir die Welt gekreuzigt, und ich der Welt.“ [Es] ist eines dem andern gestorben. Die Welt weiß nicht, was ich thue, und ich nicht, was die Welt thut; [es] kennt keines das andere, es ist alles todt und rein ab. Dieser Sabbath wird aber endlich erst recht gehalten, wenn wir gestorben sind.

157. Wenn wir nun also leben, wie Paulus von sich schreibt, so wird es alles göttlich, Hände und Füße, Zunge, Ohren, Augen, Leib und Seele, alle Gedanken; und was ich dann thue, bin ich gewiß, daß es GOtt in mir thue. Wenn ich mein Handwerk treibe und arbeite, so weiß ich, daß es GOtt wohlgefalle; denn es ist sein göttlicher Befehl. Wenn ich nun also gewiß bin, es gefalle GOtt, so sind es nicht meine, sondern GOttes Werke, denn ich thue sie in GOttes Gehorsam, und thue, was GOtt gefällt, und nicht, was mir gefällt, thue es mit Willen und ganzem Herzen. Wenn ich aber eigene Werke treibe, und meine Glieder, Ohren, Augen, Zunge, Hände und Füße, Herz und Gedanken, gerichtet sind nach meinem, und nicht nach GOttes Willen, so sind es alles Werke, die außerhalb und wider den Glauben gehen. Darum sind dies die rechten Werke des rechten Sabbathes, die aus dem Glauben, im Gehorsam und Befehl GOttes hergehen: da wirst denn Gott den rechten Sabbath.

158. Darum hat auch GOtt so streng und hart gehalten über dem Sabbath im alten Testamente, dadurch dieser geistliche Sabbath bedeutet ist worden. Denn es ist alles in einer Figur geschehen, Col. 2, 16. 17. Der alte Juden-Sabbath ist nur ein Schatten gewesen gegen dem rechten Sabbath der Christen, darum ließ GOtt auch steinigen den, der am Sabbath Holz auflas, wie 4 Mos. 15, 32. 35. geschrieben steht. Das ist aber geschehen um unsertwillen, daß wir ja mit
\quellseite{551}
Fleiß auf unsere Feier sehen, und nichts vornehmen oder ausrichten ohne GOttes Wort oder wider GOttes Wort,\footnote{„Wort“ fehlt in der Erlanger.} sondern was wir thun, daß wir es aus GOttes Befehl thun.

159. Welche aber mit eigener Vernunft und Gutdünkel drein fallen, richten Abgötterei an und predigen unrecht, geben vor falsche Lehre, unter dem Namen GOttes Worts und Schein des Worts GOttes, die verführen denn die Welt, fahren frisch heraus, daß sie nur gehört werden; meinen, sie haben es wohl ausgerichtet, wenn sie nur was Neues auf die Bahn bringen. Aber dieser Text hier vermahnt uns, daß wir müßig stehen, und uns vor solchen Werken hüten, die nicht von GOtt sind geboten; plumpen denn also, wie die Säue, herein, reißen Bilder um, und berühmen sich darnach, GOtt habe sie es geheißen, es sei GOttes Wort, aus GOttes Befehl geschehen. Es ist aber ein eigen erlesen, und nicht GOttes Wort, wie oben [§ 35 ff.] angezeigt ist im ersten Gebote. GOtt spricht zu uns: Halte stille, heilige mir den Sabbath, laß mich arbeiten. Wenn wir nun erwählen unser Werk, und lassen daneben anstehen GOttes Befehl, so halten wir des Teufels Feier, nicht GOttes.

160. St. Peter, 1. Ep. 4, 11., spricht: „So jemand redet, daß er es rede als GOttes Wort; so jemand ein Amt hat, daß er es thue, als aus dem Vermögen, das GOtt darreichet.“ Als wollte er sprechen: Alle eure Worte und Werke sollen gegründet sein in GOttes Wort und Befehl, und wo ihr deß nicht gewiß seid, so lasset es nur anstehen, daß ihr euch und andern nicht schädlich seid. Item, der heilige Paulus, wie oft zeigt er an, daß alle seine Werke, Predigt und Lehre gehen aus GOttes Befehl; wie er spricht Eph. 3, 7.: „Nach der Gabe aus der Gnade GOttes, die mir nach der Wirkung seiner Kraft gegeben ist.“ Item, 1 Cor. 15, 10.: „Von GOttes Gnaden bin ich, das ich bin, und seine Gnade an mir ist nicht vergeblich gewesen, sondern ich habe viel mehr gearbeitet denn sie alle; nicht aber ich, sondern GOttes Gnade, die in mir ist.“ Item, 2 Cor. 13, 3.: „Ihr suchet, daß ihr gewahr werdet deß, der in mir redet, Christus.“ Item, Röm. 15, 17–19.: „Ich kann mich rühmen durch Jesum Christum, daß ich mit göttlichen Sachen umgehe. Denn ich dürfte nicht etwas reden, wo dasselbe Christus nicht durch mich wirkte, die Heiden gehorsam zu machen, durch Wort und Werk, durch Kraft der Zeichen und Wunder, und durch Kraft des Geistes GOttes“ etc.

161. Also genau haben die heiligen Leute darauf gesehen, daß sie sich ja nichts unterwunden, sie wären denn zuvor gewiß, daß es GOtt in ihnen wirkte; wo sie aber nicht gewiß waren, hielten sie stille mit alle ihrem Thun, Singen und Sagen. Das heißt denn der rechte Sabbath, welchen die ganze Welt hat lassen anstehen, und des Teufels Feier angenommen. Der Pabst mit all seinem Anhange, als Mönchen, Nonnen und Pfaffen, haben lauter Menschenwerk gelehrt, und die rechten göttlichen und geistlichen Werke unterwegen gelassen, ja, verachtet und für ketzerisch gehalten, haben ihr Leben zugebracht in eigen erdichteten Werken und Weisen.

162. Also auch jetzund, da der rechte Sabbath wieder auf den Plan durch die Predigt des heiligen Evangelii kam, erhoben sich von neuem die Rotten und Secten, richteten neue Werke nach ihrem eigenen Dünkel auf, zerrissen also den rechten Sabbath, daß GOtt sein Werk nicht in uns mag haben.

163. Denn, soll der rechte Sabbath gehalten werden, so muß der alte Adam mit seinem Dünkel gedämpfet und getödtet sein, daß also inwendig im Herzen nichts geschehe, es thue es denn GOtt selber, auf daß also ein jeglicher in seinem Stande gewiß sei, es gefalle GOtt sein Thun. Ein ehelich Weib soll gewiß sein in ihrem Stande, daß ihr Kindertragen, Gebären, Säugen und Kinderwarten GOtt so wohl gefalle, als hätte er es selber mit ihr geredet, und ihr fleißig befohlen. Also eine Magd, wenn sie das Haus kehrt, kocht, wartet des Viehes, soll auch den Trotz haben, daß sie da den rechten Sabbath halte, und in GOttes Befehl gehe, wenn sie treulich ausrichtet, was ihr befohlen ist. Also durchaus, in allen Ständen, soll man den Trotz behalten, daß sie GOtt in solche Stände und Werke geworfen habe, und soll ein jeglicher ihm seinen Stand wohlgefallen lassen, er sei so schlecht er immer wolle. Da werden denn alle Stände gleich im Glauben. Denn GOtt sieht nicht, wie groß oder klein dein Stand sei, sondern ob dir der Stand gefalle, und du ihn annehmest als von GOtt gegeben. Daher denn ein Hirte so wohl vor GOtt treten kann, als
\quellseite{552}
der Kaiser, und hindert ihn sein geringer, verachteter Stand gar nichts.

164. Darum so soll man allein GOttes Wort, Befehl und Werk ansehen, das allein preisen und loben; das andere alles schelten, das außerhalb dem Worte und Befehl GOttes geht. Wenn nun eine Magd des Viehes wartet, eine Frau Kinder trägt, ein Mann seines Handwerks wartet, ein Knecht die Pferde beschickt, ein Fürst oder Herr einem bösen Buben den Kopf läßt abschlagen, mögen sie alle, ein jegliches in seinem Stande, trotzen, und sagen: GOtt hat es gethan.

165. Also auch, die recht und aus dem Glauben fluchen, können sprechen: GOtt flucht dir. Denn so haben die Propheten den Gottlosen geflucht, ja, GOtt hat es in ihnen gethan. Auf diese Weise aber haben die Propheten geflucht: HErr, laß sie zu Schanden werden, laß sie schamroth stehen, zerstöre und zerschmeiß ihren Rath und Anschlag, und alle ihre Lehre, treib sie zurück etc. Also betet auch Moses wider Korah und seine Rotte [4 Mos. 16, 15.]: „HErr, siehe nicht auf ihr Opfer“, laß sie nicht fortfahren.

166. Wenn du nun merkst bei dir ein Werk, das GOtt nicht in dir wirkt, so tritt es mit Füßen, und bitte GOtt, daß er auch in dir zu Schanden mache alles, das er nicht selber wirkt. Und wenn du mit Einer Predigt könnest die ganze Welt selig machen, und hast den Befehl nicht, so laß es nur anstehen, denn du wirst den rechten Sabbath brechen, und wird GOtt nicht gefallen.

167. Da gehören nun her die Werke, damit man die Natur und den alten Adam tödtet, als, fasten, den Leib kasteien etc. Von welchen, wie weit und lang sie gebraucht sollen werden, ist gehandelt im Büchlein von der christlichen Freiheit, und von guten Werken.\footnote{Gemeint ist Luthers Sermon von guten Werken, Walch, St. Louiser Ausg., Bd. X, 1298.} Von dem Sabbath sagt auch die Epistel an die Hebräer am dritten, V. 11., und im 4. Capitel, V. 5.

168. Das ist aber wohl zu merken, daß solche Sabbaths-Werke sollen im Glauben und guter Zuversicht göttlicher Gnade geschehen, auf daß alle Werke im ersten Gebote und Glauben bleiben, und der Glaube sich in denselben übe und stärke, um welches willen alle anderen Gebote und Werke geschrieben sind.

169. Darum siehe, wie ein hübscher güldener Ring aus diesen dreien Geboten und ihren Werken sich selber macht, und wie aus dem ersten Gebote und Glauben fließt das andere, bis in das dritte; und das dritte wiederum treibt durch das andere, bis in das erste. Denn des ersten Gebots Werk ist glauben, ein gut Herz und Zuversicht zu GOtt haben. Aus dem fließt das andere gute Werk, GOttes Namen preisen, seine Gnade bekennen, und ihm allein alle Ehre geben. Darnach folgt das dritte, Gottesdienst üben mit Beten, Predigt hören, dichten und betrachten GOttes Wohlthat, dazu sich kasteien und sein Fleisch zu zwingen.

170. Wenn nun der böse Geist solches Glaubens, GOttes Ehre und Gottesdiensts gewahr wird, so tobt er, und hebt an zu verfolgen, greift Leib, Gut, Ehre und Leben an, treibt auf uns Krankheit, Armuth, Schande und Sterben, welches auch GOtt also verhängt und verordnet über uns. Siehe, da hebt sich denn das andere Werk, oder die andere Feier des dritten Gebots; dadurch wird der Glaube fast hoch versucht, wie das Gold im Feuer. Denn es ist ein groß Ding, eine gute Zuversicht zu GOtt zu haben, ob er uns schon den Tod, Schmach, Krankheit, Armuth zufügt, und in solchem grausamen Bilde des Zorns ihn für den allergütigsten Vater zu halten; welches geschieht in diesem Werke des dritten Gebots. Da bringt denn das Leiden den Glauben, daß er GOttes Namen muß anrufen und loben in solchem Leiden, und kommt denn also durch das dritte Gebot wiederum in das andere; und durch dasselbe Anrufen göttliches Namens und Lobes wächst der Glaube, und kommt in sich selbst, und stärkt also sich selbst, durch die zwei Werke des dritten und andern Gebots. Und also geht er aus in die Werke, und kommt wieder durch die Werke zu sich selbst; gleichwie die Sonne aufgeht bis an den Niedergang, und kommt wieder zum Aufgange etc.

171. Also haben wir die Summa der drei ersten Gebote. Im ersten ist geboten, wie sich unser Herz gegen GOtt halten soll mit Gedanken; im andern, wie sich der Mund mit Worten; in dem dritten wird geboten, wie wir uns gegen GOtt halten sollen in Werken. Und das ist die erste und rechte Tafel Mosis, in welcher diese drei Gebote beschrieben sind, und den Menschen regieren auf der rechten Seite, das ist, in den Dingen, die GOtt anlangen, und in welchen GOtt mit dem Menschen, und er mit GOtt zu thun hat, ohne Mittel irgend einer Creatur.

\quellseite{553}
\subsection{Die andere Tafel Mosis}

172. Die Gebote in der andern Tafel fahren hiernieder zum Nächsten, und zeigen an, wie ein Mensch gegen dem andern sich halten soll. Folgt also das erste Gebot der andern Tafel:

\begin{schriftzitat}
V. 12. Du sollst deinen Vater und deine Mutter ehren, auf daß du lange lebest im Lande, das dir der HErr, dein GOtt, geben wird.
\end{schriftzitat}

\subsection{Summa}

173. Das erste lehrt, wie man sich halten soll gegen alle Obrigkeit, welche an GOttes Statt sitzt, als sind, Vater und Mutter, Herren und Frauen ehren etc. Darum folgt dies vierte Gebot den andern ersten dreien, die GOtt selbst antreffen.

\begin{quote}
Du sollst deinen Vater und deine Mutter ehren, auf daß du lange lebest etc.
\end{quote}

174. Dieser Zusatz, daß er spricht: „auf daß du lange lebest im Lande, das dir der HErr, dein GOtt, geben wird“, gibt auch Anzeigung, daß die zehn Gebote allein den Juden gegeben sind, und nicht den Heiden. Denn er redet hier zu denen, welchen das Land Canaan verheißen war; nun ist es nur den Juden verheißen, wie er hernach lehrt.

175. Dies Gebot, spricht St. Paulus Eph. 6, 2., ist das erste, das eine Zusagung und eine Verheißung hat von GOtt. Wer nun den Eltern gehorsam ist, Vater und Mutter ehrt, der wird lange leben, und im Lande bleiben. Nun habt ihr droben, in dem ersten Gebote, gehört, daß es auch eine Zusagung und Verheißung habe, und das andere eine Dräuung; wie spricht denn St. Paulus, daß das vierte Gebot sei das erste, das eine Zusagung habe von GOtt? Antwort: Paulus führt die Gebote gewöhnlich ein der andern, und nicht der ersten Tafel. Also auch Christus, da er von der Liebe des Nächsten redet, führt die Gebote der andern Tafel ein. Da ihn der Schriftgelehrte fragte, was er thun sollte, sprach er [Matth. 19, 16, 18.]: „Du sollst Vater und Mutter ehren, nicht tödten, stehlen, ehebrechen, nicht falsch Zeugniß reden“ etc., gedenkt der ersten drei Gebote nicht.

176. So ist nun dies die Ursache auf die Frage: St. Paulus redet von Christen, die jetzt genugsam unterrichtet sind im Glauben, und in dem, das GOtt antrifft; und, wie in den ersten dreien Geboten wird angezeigt, die jetzund schon voll des rechten Gottesdienstes sind. Als wollte St. Paulus sprechen: Ihr dürfet jetzund nicht mehr, denn daß ihr euch gegen euren Nächsten haltet, wie sich GOtt gegen euch gehalten hat. Nun, die ersten Gebote werden nicht erfüllt mit Werken, sondern durch den Geist GOttes im Herzen. Den Glauben aber, den sie im Herzen zu GOtt haben, sollen sie äußerlich beweisen in den Werken, also, daß die Werke anzeigen, wie sie vor GOtt stehen, doch also, daß diese Werke durch den Geist im Herzen gegründet seien. Also führt das Neue Testament die Gebote der andern Tafel; darum es auch Paulus nennt, daß es sei das erste Gebot, verstehe, der andern Tafel, das eine Verheißung habe, nämlich, wenn du Vater und Mutter ehren wirst, sollst du im Lande bleiben; wo nicht, sollst du aus dem Lande gestoßen werden.

177. Da seht ihr, wie GOtt dies Gebot von Vater und Mutter so hoch hält, daß er es schier dem ersten Gebote gleich macht, führt schier solche Worte. Denn wie das erste Gebot eine Verheißung hat, nämlich die: Wer es hält, soll Barmherzigkeit erlangen bis in viel tausend etc., also hat dies vierte auch vor den andern Geboten der andern Tafel seine Zusagung denen, die es halten, nämlich, daß sie lange leben sollen etc. Also will er nach sich nichts mehr gefürchtet und geehrt haben denn Vater und Mutter.

178. Zum andern gebraucht er des herrlichen Worts: „Du sollst deinen Vater und deine Mutter ehren.“ „Ehren“ ist nicht ein geringes Wort; spricht nicht: Du sollst deinen Vater und deine Mutter lieb haben, ihnen gehorsam sein, wohlthun, oder dergleichen, sondern, du sollst sie ehren. Die Ehre streckt sich aber weiter aus denn die Liebe; darum hält GOtt viel von Vater und Mutter. Denn die Ehre gebührt allein GOtt. Nun theilt er die Ehre Vater und Mutter mit; darum auch keine größere Herrschaft auf Erden ist denn der Eltern Herrschaft. Aber sie ist gemein worden, darum achtet ihrer niemand
\quellseite{554}
, wie denn allen GOttes Worten und Werken geschieht; wenn er es der Welt reichlich vorlegt, so verachtet sie es, wird ihrer bald satt, müde und überdrüssig. Die Könige, Fürsten und Herren haben auch ihre Ehre, tragen eine güldene Krone auf dem Haupte, einen güldenen Scepter in der Hand, güldene Ketten am Halse; nun man soll sie auch ehren, Röm. 13, 7. 1 Petr. 2, 17., und fürchten. Es ist aber nicht eine gleiche Ehre, gegen der Ehre, die man Vater und Mutter soll beweisen.

179. Denn der Fürsten und Herren Obrigkeit ist nicht eine liebliche Obrigkeit, sondern erschreckliche, denn sie sind unseres HErrn GOttes Stockmeister, Richter und Henker, mit denen er die bösen Buben straft; aber Vater und Mutter sind nicht also erschrecklich, sondern ganz freundlich. Es steht Röm. 13, 3. ff. von der Obrigkeit, daß sie soll wehren den bösen, wilden Leuten, und die Frommen handhaben; darum nimmt sie Geld, Zinse und Schoß. Aber Vater und Mutter ist eine zarte, feine, lustige Obrigkeit, die nehmen nicht von Kindern, sondern sie wagen Leib und Leben, schlagen in die Schanze um der Kinder willen Hab und Gut, setzen für sie Hals und Bauch und alles, was sie haben. Weltliche Obrigkeit will nur haben; diese kann nichts denn geben, denn die Kinder sind aus ihrem Fleische und Blute. Solches können sich die Fürsten und Herren nicht berühmen, denn wir sind ja nicht ihr Fleisch und Blut.

180. Aber Vater und Mutter, die können sichs wohl rühmen. Sobald das Kind geboren wird, muß [es] ernähret werden von Schweiß, Fleisch und Blut seiner Eltern. Sie warten sich selber, sorgen Tag und Nacht, und ist da lauter Liebe, und gar kein Ernst oder Zorn. Denn wenn es nicht also lieblich zuginge, Vater und Mutter die Kinder nicht also herzlich liebten, wie wollten sie erzogen und ernähret werden? Sie müßten im Dreck verfaulen. Aber da ist lauter Liebe, Dienen und Wohlthun von Vater und Mutter; [sie] strecken Leib und Gut daran, tragen Sorge und Angst; [es] ist lauter Mühe und Arbeit Tag und Nacht vorhanden für die Kinder; ja, es muß den Kindern alles dienen, Gut, Haus und Hof, Magd und Knecht.

181. Also scheidet sich die Obrigkeit und Gewalt der Eltern von der Gewalt der weltlichen Herren. Bei den Eltern ist keine Furcht und Schrecken, sondern lauter Liebe. Bei der Obrigkeit ist nicht viel Liebe, sondern Furcht und Erschrecken, Röm. 13, 3. 4.: „Die Gewaltigen sind den bösen Werken zu fürchten.“ 1 Petr.: „Thust du Böses, so fürchte dich; denn sie trägt das Schwert nicht vergeblich. Sie ist GOttes Dienerin, eine Rächerin zur Strafe über den, der Böses thut.“ Diese Gewalt gibt nicht, sie nimmt von uns, damit sie könnte die Guten schützen, und die Bösen strafen.

\subsection{GOtt wird erkannt und sein abgemalt im Bilde der Eltern}

182. Hier sollen wir auch lernen, daß Vater und Mutter GOtt ganz ähnlich sind in ihrem Amte gegen den Kindern, und ist uns in ihnen sein abgemalt das göttliche und väterliche Herz gegen uns. Denn in Vater und Mutter können wir spüren und erfahren, wie GOtt gegen den Menschen gesinnt ist. Darum er sich auch nicht schämt des väterlichen Namens, und Christus, sein Sohn, nennt sich einen Bräutigam, Matth. 9, 15.: „Wie können des Bräutigams Kinder Leide tragen, so lange der Bräutigam bei ihnen ist?“ und Joh. 3, 29.: „Wer die Braut hat, der ist der Bräutigam“ etc., und nennt die Gemeine seiner Gläubigen die Braut. Und 2 Cor. 11, 2.: „Ich habe euch vertraut einem Manne, daß ich eine reine Jungfrau Christo zubrächte“; und die aus seiner Braut geboren werden durch die Predigt und Glauben, heißen Kinder GOttes, Joh. 1, 12. Matth. 6, 9. Röm. 8, 16.

183. Das ist nun sein abgemalt in der Ehe. Aber vor der Welt ist es ein leichtfertig Ding, die solches alles verachtet, denn sie muß blind und verstockt bleiben, daß sie nicht sehe die großen Wunder und hohen Werke GOttes. Wenn wir anders nichts hätten von GOtt, denn die große, herrliche Güte unserer Eltern, in welcher GOtt seine Güte erzeigt, könnten wir GOtt nicht genug dafür danken. Aber die Welt achtet sein nicht; [es] ist ihr ein schlecht Ding ein solch fein, köstlich GOttes Werk, Vater und Mutter sein, Kinder zeugen, ihrer warten. Ja, unsere Ungeistlichen rümpfen die Nase drüber, wenn man den Ehestand lobt, sehen nicht, daß der Reichthum Gottes daran hangt. Denn GOtt befiehlt Vater und Mutter also das Amt, daß sie der Kinder warten, dabei man lernen und gleichwie in einem Spiegel sehen kann, wie GOtt gegen uns
\quellseite{555}
gesinnt sei, nämlich, wie des Vaters Herz gegen den Kindern, also steht GOttes Herz gegen dir.

184. Daher kommt denn das gemeine Sprüchwort, und ist auch wohl wahr: daß Vater und Mutter können an den Kindern den Himmel verdienen und die Hölle, wenn sie denen wohl oder übel vorstehen. Denn Vater und Mutter müssen sorgen und gedenken, wie sie die Kinder leiblich versorgen mit Essen, Trinken, Schuh und Kleidern, und auch an der Seele, daß sie recht GOtt erkennen lernen durch sein Wort; also sind die Hungrigen, Durstigen, Nackten, Gefangenen, Kranken etc., die Vater und Mutter zu versorgen haben, die Seelen der Kinder. Da macht GOtt aus eines jeglichen Hausvaters Hause, der da Kinder hat, ein Spital, und setzt ihn zu einem Spittelmeister, daß er seiner Kinder warten soll, sie speisen, tränken, und mit guter Lehre und Exempel vorstehen, daß sie lernen GOtt vertrauen, glauben, ihn fürchten und ihre Hoffnung auf ihn setzen, seinen Namen ehren, nicht schwören, noch fluchen, sich lassen mit Beten, Fasten, Wachen, Arbeiten, Gottesdienstes und Worts warten, und ihm den Sabbath feiern, daß sie zeitlich Ding lernen verachten, Unglück mit Sanfte und Geduld tragen, und den Tod nicht fürchten, das Leben nicht lieb haben.

185. Siehe, welch große Lectionen dies sind. Siehe, wie viel gute Werke du vor dir hast in deinem Hause, an deinen Kindern, die solches alles bedürfen, wie eine hungrige, durstige, bloße, arme, gefangene, kranke Seele. Wie eine selige Ehe wäre das, wo solches Ehevolk beisammen wäre, und ständen also ihren Kindlein vor. Fürwahr, ihr Haus wäre eine rechte Kirche, ein auserwählt Kloster, ja, ein Paradies. Denn Vater und Mutter werden GOtt hier gleich, denn sie sind Regenten, Bischöfe, Papst, Doctor, Pfarrherr, Prediger, Schulmeister, Richter und Herr. Der Vater hat alle Namen und Amt GOttes über seine Kinder; und gleichwie für uns GOtt sorgt, nährt uns, schützt und schirmt, lehrt und unterweist uns, also auch der Vater lehrt das Kind, nährt und versorgt es.

186. So könnte auch das Kind den Eltern nicht näher sein, denn es ist, nämlich, Fleisch und Blut, ja, die Natur seiner Eltern. Darum hat ein fromm Kind zu seiner Creatur eine größere Zuversicht denn zu seinen Eltern, in welchem uns so trefflich sein abgemalt ist, wie GOtt gegen uns, und wir gegen ihm gesinnt sind. Denn wie sich ein Kind alles Gutes zu seinen Eltern versieht, also versieht sich ein Christ alles Gutes gegen GOtt; und wiederum, GOtt stellt sich gegen einen Christen wie ein Vater gegen seinem Kinde, und noch wohl freundlicher. Es weiß auch ein Christ, daß GOtt mehr für ihn sorgt denn alle Menschen und Creaturen, ja, mehr denn er selber.

187. Daß nun die Kinder solches an ihren Eltern erkannten, hat ihnen GOtt geboten, sie sollen Vater und Mutter ehren. Sie sollen nicht allein ansehen Fleisch und Blut an ihren Eltern; denn wo sie das allein ansehen, so finden sie nichts Köstliches an ihnen, und werden die Eltern bald verachten; sie sollen die Augen aber aufthun, und das Gesicht über Fleisch und Blut heben, da werden sie in ihren Eltern ein wunderlich Ding finden.

188. Denn man muß zwei Dinge an den Eltern sehen: zum ersten, daß sie Fleisch und Blut sind; zum andern, das Kleinod, das GOtt an die Eltern gehängt hat, nämlich sein Wort. Denn er hat sie also in sein Wort gefaßt, wie in eine Monstranz, und sie bekleidet mit seinem Willen. So muß man die Eltern nun ansehen, als die das Wort und den Willen GOttes tragen. Gleichwie vor Zeiten der Pabst die Todtenbeine gefaßt hat in köstliche Monstranzen, in Gold, Silber, Seide, Purpur und Edelgestein: das hatte denn ein groß Ansehen um solcher Pracht willen; sonst wäre es wohl in\footnote{„in“ fehlt in der Wittenberger.} Verachtung geblieben, wenn es nicht also geschmückt wäre gewesen. Aber das\footnote{Erlanger: da.} ist das rechte lebendige Heiligthum in Vater und Mutter. Denn GOtt hat da sein Wort hingelegt, darinnen die ganze göttliche Majestät ist. Und ist also sein göttlicher Wille, daß man sie ehren soll; darum man ja nicht sie allein nach Fleisch und Blut soll messen, sondern nach dem Worte Gottes.

189. Man ist viel zu St. Jakob gelaufen, und zu andern Heiligen, gen Rom, gen Jerusalem, gen Aach, das Heiligthum zu besehen; aber hier will niemand zu dem rechten Heiligthum laufen. Warum? Es ist GOttes Werk, darum wird man sein bald überdrüssig; der Teufel zeucht uns davon, Fleisch und Blut mag es nicht. Denn es gleibt nicht, darum gilt es auch vor
\quellseite{556}
der Vernunft nicht. Da läuft einer zu St. Jakob; der gelobt sich zu Unserer Frauen; ein anderer läuft in das Kloster; der aber keines hat dir GOtt befohlen, sondern will, daß du deinen Eltern dienest; da findest du ihn, nicht um der Eltern, sondern um seines Worts willen.

190. Aber wie die Kinder GOttes Gebot hierinnen verachten, und nehmen ihnen sonderliche Weise vor, GOtt zu dienen, also thun die Eltern auch. Da gelobt eins dies, das andere das; niemand aber gelobt, daß er, GOtt zu Ehren, seine Kinder wohl regiere und lehre; läßt die sitzen, die ihm GOtt befohlen hat, an Leib und Seele zu bewahren, und will GOtt an einem andern Orte dienen, davon ihm nichts befohlen ist. Solch verkehrt Wesen wehrt kein Bischof, straft kein Prediger, ja, um Geizes willen bestättigen sie es, und erdenken täglich nur mehr Wallfahrten, Heiligen-Erhebung, und ander solch Teufelsgespenst.

191. Willst du aber das rechte Heiligtum finden, so bleib daheim im Hause, du seiest Vater oder Kind; darfst nicht zu St. Jakob laufen, gen Rom, gen Aach,\footnote{In den Ausgaben fehlt das Komma hinter „Aach“. Nun ist zwar zu Aachen ein Gewand der heiligen Jungfrau, doch daß daselbst ein besonderer Mariencultus gewesen sei, ist uns nicht bekannt. Dagegen finden wir bei Luther besonders erwähnt: gen Lorett zur Maria, Maria zur Eiche, zur Maria gen Regensburg.} zu Unserer Frauen, zum heiligen Blut, in das Kloster. Siehe Vater und Mutter an, da wird dich das Wort GOttes lehren, wenn du Vater und Mutter ehrst, daß du nicht Fleisch und Blut, sondern GOtt selber ehrst, der sein Wort auf sie gelegt hat. Und ihr Eltern, steht ihr euren Kindern wohl vor, so thut ihr GOtt einen großen Gefallen daran, wie St. Paulus 1 Tim. 5, 4. 8. lehrt: „Eine Wittwe, die Kinder oder Neffen hat, solche laß zuvor lernen, daß ihre eigenen Häuser gottselig seien, und Gleiches vergelten ihren Vor-eltern; denn das ist angenehm vor GOtt. So aber jemand die Seinen, sonderlich seine Hausgenossen, nicht versorgt, der hat den Glauben verleugnet, und ist ärger, denn ein Ungläubiger.“ Wenn wir nicht so toll und thöricht, und ganz verstocket wären, und besessen von der Welt Fürsten, dem Teufel [Joh. 12, 31.], könnten wir wohl merken und greifen, daß dies die rechten Werke wären, die aus GOttes Wort und Befehl hergehen.

192. Dies sind aber bisher der Welt gute Werke gewesen: sich selber einschließen, wie ein wild Thier, in die Klöster, die Körner am Paternoster zählen, Kappen und Platten tragen, härene Hemde und graue Röcke anziehen, und dergleichen Affenspiel treiben. Aber dies löbliche Werk, der Eltern Ehre, und die Kinder ziehen, hat niemand können sehen, so hoch aus diesen Werken, und ihres Gleichen, im Glauben gethan, ein rechter Christ erkannt wird. Aber die Welt will nicht Werke haben, die gut sind, sondern die da gleichen und geschmückt sind; da läuft sie zu, gibt mit beiden Händen. Wo aber ein frommer und gehorsamer Sohn daher gegangen ist in einem rothen oder grünen Rocke, hat Vater und Mutter geehrt, pfui, der ist nichts gewesen.

193. Nun wäre hier anzuzeigen, wie man wider dies\footnote{So die Jenaer. Wittenberger und Erlanger: das.} Gebot sündiget, wie oben kurz begriffen; es würde aber allzu lange. Kurz davon zu reden: Der Papst hat dies Gebot gar zunichte gemacht und aufgehoben. Also ist es auch bei den Juden zugegangen, wie Christus anzeigt Matth. 15, 3–5. „Um eurer Satzungen\footnote{Jenaer: Aufsätze. Die Erlanger bezeichnet dies wieder als eine Note Walchs.} willen“, spricht Christus, „habt ihr aufgehoben GOttes Gesetz. GOtt hat geboten: Du sollst Vater und Mutter ehren; wer aber Vater und Mutter flucht, der soll des Todes sterben. Aber ihr sprecht zu den Kindern, daß sie zu den Eltern sagen sollen: Corban“; das ist, wenn ein erwachsen Kind Geld und Gut hat, und sein Vater und Mutter waren arm, dürften Hülfe von ihrem Kinde, da fahren die Phariseer zu, schlossen sich dahin, man soll es GOtt geben und zu Gottesdienst, da wäre es besser angelegt, denn wenn es den Eltern gegeben würde. Das ist denn nichts anders, denn GOtt Lügen strafen.

194. Also ist es bei uns auch zugegangen: Wenn man Vater und Mutter sollte geben, wär es eine schlechte Sache; wenn man aber Jahrstage, Seelmessen, Vigilien stiften sollte, den Pfaffen und Mönchen mit Haufen zutragen, ihnen zu ihren faulen Tagen helfen, das wär ein geistlich gut Werk, die waren unsere Fürbitter, huben uns in Himmel. Aus, zum Teufel mit den guten Werken! Ueber das verließen die Kinder die armen Eltern, die wohl ihrer Handreichung bedurft hätten, liefen in die Klöster, gaben vor, sie wollten GOtt dienen. Ja, dem Teufel. Denn,
\quellseite{557}
heißt das GOtt dienen, wider GOttes Gebot Vater und Mutter im Elend verlassen, so ist es mir ein seltsamer Gottesdienst. Darum ist aller Mönche und Nonnen Stand wider dies Gebot, davon genugsam geschrieben ist im Büchlein von den Klostergelübden.\footnote{Dies ist die Schrift: „Luthers Urtheil von den geistlichen und Klostergelübden“, Walch, St. Louiser Ausg., Bd. XIX, 1500.}

195. Von der Strafe der ungehorsamen Kinder steht also, 5 Mos. 21, 18–21.: „Wenn jemand einen eigenwilligen und ungehorsamen Sohn hat, der seines Vaters und Mutter Stimme nicht gehorcht, und, wenn sie ihn züchtigen, ihnen nicht gehorchen will, so soll ihn sein Vater und Mutter greifen, und zu den Aeltesten der Stadt führen, und zu dem Thore desselben Orts, und zu den Aeltesten der Stadt sagen: Dieser, unser Sohn, ist eigenwillig und ungehorsam, und gehorcht unserer Stimme nicht, und ist ein Schwelger und Trunkenbold; so sollen ihn steinigen alle Leute derselbigen Stadt, daß er sterbe“ etc.

\subsection[Was die Ehre gegen die Eltern sei]{Was die Ehre gegen die Eltern sei.\footnote{„gegen die Eltern“ fehlt in der Erlanger Ausgabe.}}

196. Die erste Ehre, die man Vater und Mutter soll erzeigen, ist, daß man ihnen gehorsam sei, folge ihrem Gebote, wie St. Paulus spricht Eph. 6, 1.: „Ihr Kinder, seid gehorsam euren Eltern.“ Die andere Ehre, wenn wir nun erzogen sind, und jemand selber Mann und Weib sind worden, wo es Vater und Mutter fehlt, daß sie arm, hungrig, durstig, nackend, krank und schwach sind, daß, wir ihnen die Hand reichen, ihnen helfen, dienen mit Speisen, Tränken, Kleiden und allerlei Nothdurft reichen, und sie für das größte Heiligthum halten, das auf Erden ist. Denn die Ehre steht nicht allein in Worten und Gebärden, sondern vielmehr in der That. Es wäre eine kleine Ehre, wenn ich den Hut vor meinen Eltern abzöge, und ließe sie gleichwohl daneben Hunger leiden. „Nicht allein mit dem Worte (spricht St. Johannes in seiner ersten Epistel Cap. 3, 18.) sollen wir lieben, sondern mit der That und Wahrheit“, also, daß das Werk der Liebe mit Willen, und aus einem rechten gläubigen Herzen gehe. Das heißt denn, mit der Wahrheit und mit der That lieben. Item, daß man sie groß achte, und viel von ihnen halte, und uns ganz geringe achten gegen ihnen, daß wir wissen, daß es GOtt also wohlgefalle, daß wir von unsern Eltern sollen geboren werden, sie seien hohes oder niedriges Standes, reich oder arm; auch was sie uns heißen, daß nicht wider GOtt ist, daß wir nicht anders gedenken, denn daß es GOtt heiße, und von uns fordere. Denn es ist also von GOtt versehen, daß wir ihre Söhne und Töchter worden sind. Darum wir ihre Worte und Werke also annehmen sollen, als schaffte sie GOtt selbst mit uns. Darum spricht er nicht: Du sollst sie lieben, sondern du sollst sie ehren.

197. Denn Ehre ist etwas Größeres, denn Liebe. Liebe ist gegen denen, die uns gleich sind; als, wenn zwei einander lieb\footnote{Erlanger: an einander lieb.} haben, achtet sich ein jeglicher nicht höher denn den andern. Aber die Ehre geht gegen einem Höhern, und hat eine Furcht mit sich, daß man den nicht beleidige, so wir ehren, und thut sich auch unter den, als unter einen Herrn, nicht um der Strafe willen, sondern daß er ihn lieb hat, den er ehrt. Also sollen wir auch Vater und Mutter nicht allein lieben, sondern auch ehren, sie hoch achten, uns vor ihnen, als unsern Oberherren, von GOtt verordnet, fürchten.

198. Und niemand ist hier ausgenommen, den Eltern solche Ehre zu leisten, er sei so hohes Standes er immer wolle. Ja, der Kaiser, Pabst, und alle Cardinäle sind nicht über das Gebot, denn sie sind ja nicht über GOtt. Darum ist der Pabst und Kaiser unter Vater und Mutter, nach Ausweisung dieses Gebots. Und solches bestätigt Christus auch Matth. 22, 45., da er spricht: „So David ihn einen HErrn nennt, wie ist er denn sein Sohn?“ Als wollte er sagen: Wenn er sein Herr wäre, so wäre der Sohn über den Vater. Das reimt sich nicht. Des Königs Sohn muß noch seinen Hut vor seinem Vater abziehen, wie auch Salomon that; wiewohl er König war, 1 Kön. 2, 19., stand er auf, und ging seiner Mutter entgegen, und betete sie an.

199. Darum übertrifft der Eltern Obrigkeit und Herrschaft alle andere Obrigkeit, des Pabsts, Kaisers, der Könige, Fürsten und Herren. Es ist der höchste Stand, schwebet hoch über alle Stände, und die andern sind nur Stück- und Flickwerke gegen ihm; ja, alle Stände kommen daher. Aber weil wir alle Vater und Mutter haben, ist der Stand ein gemein, verachtet Ding worden, es wird nicht viel davon gehalten, wie
\quellseite{558}
Gottes Werken geschieht; so sie gemein werden, wird man ihrer bald müde, sucht nur, was neu und seltsam ist.

200. Nun, solche Ehre soll den Eltern nicht allein mit Haupteigen und Hutabziehen, wiewohl dasselbe jetzt auch seltsam ist, geschehen, sondern von ganzem Herzen. Denn GOtt spricht nicht: Deine Hand, Mund, Zunge oder Knie solle Vater und Mutter ehren, sondern „du“. Was ist aber „du“? „Du“ heißt nicht deine Hand, Zunge, Herz, Leib oder Seele, sondern das alles zu Haufe mit deinem ganzen Wesen und Natur, wie du gehst und stehst, Leib und Seele, Sinn und Witz, was an, aus und in dir ist.

201. Wie, wenn sie wider GOtt etwas gebieten, soll ich gehorchen? Nein, da ist GOtt ausgenommen; da sollen wir nicht gehorchen, wenn sie gebieten wider die Gebote der ersten Tafel; die sollen den Vorzug haben. Da soll der Sohn sprechen zum Vater: Wir haben einen GOtt, der ist mehr denn du. Ich will dir gerne gehorsam sein, so ferne es nicht wider GOtt ist, wie Petrus spricht Apost. 5, 29.: Man muß GOtt mehr gehorsam sein denn den Menschen. Es wäre wohl fein, daß [es] in der Welt also zuginge, daß die Kinder die Eltern ehrten etc., so aber die Kinder nicht wollen gehorchen, befiehlt GOtt den Eltern die Ruthe. Und wo sie nicht stark genug sind, hat ihnen GOtt einen Gehülfen gegeben, der die frechen, ungezogenen Kinder soll strafen, nämlich, die weltliche Obrigkeit; wie im Alten Testamente GOtt gebot, 5 Mos. 21, 18–21., wie oben [§ 195] angezeigt ist.

\subsection{GOtt zeigt an durch die zehn Gebote, was er von der Welt hält}

202. Nun können wir wohl abnehmen, was GOtt von uns halte, daß er uns solche Gebote gibt, daß wir nicht sollen fremde Götter haben, seinen Namen nicht unnütz brauchen, seinen Sabbath heiligen. Damit uns angezeigt wird, daß wir vor seinen Augen verdächtig sind; ja, er hält nicht mehr von uns, denn daß wir verzweifelte Buben in der Haut sind, die voll Abgötterei stecken, seinen Namen lästern und schänden und den Sabbath brechen. Denn, wenn er uns für fromm hielte, was dürfte er uns das alles gebieten? Aber er läßt ein solch Gerücht von uns ausgehen, läßt von uns singen und sagen, daß wir so feine, fromme Kinder sind, die ihren eigenen Vater, ihren GOtt und Schöpfer nicht kennen, ja, verachten ihn, und schmähen und schänden [ihn] auf das allerhöchste.

203. Das ist ja verdrießlich, daß er niemand ausnimmt, daß er nicht der Phariseer und Schriftgelehrten, auch nicht unserer Mönche und Pfaffen verschont, der heiligen, frommen Leute; läßt das Gericht straks über Könige und Kaiser, Pabst, Cardinäle, Fürsten und Herren gehen etc., läßt von ihnen singen und sagen, daß sie GOtt nicht kennen, lieben, sondern schänden und lästern, ja, daß sie auch Mörder, Todtschläger, und Buben in der Haut sind; warum verschont er nicht der heiligen Leute? Es gelüstet ihn nicht. Er jagt es nicht zu mir und dir allein, sondern zu allen Leuten: „Ehre Vater und Mutter.“ Er weiß, was in uns ist.

204. Darum kommt er zuvor, erschreckt uns mit seinen Geboten, damit er anzeigt, was er von uns halte. Es ist gar eine große Schlappe; wir mögen wohl schamroth werden. Wenn zu mir einer sagte: Lieber Freund, ehre Vater und Mutter, brich deine Ehe nicht, stiehl nicht; Lieber, schlage deinen Vater und Mutter nicht; ich sollte es bald mit Unwillen annehmen, sollte sprechen: Bist du toll, für wen hältst du mich? Nun hält uns GOtt nicht besser. Sollte es uns nicht von GOtt verdrießen, daß er also mit Ernst mit uns handelt, hält uns für Ungehorsame, Mörder und Diebe etc.? Ach lieber HErr GOtt, hältst du mich mehr von mir, denn daß ich ein solcher soll sein, der einem andern bei seinem Weibe schlafe, nach des andern Gut stehe etc.? Wenn wir es mit Ernst betrachteten, was GOtt damit meinte, so fänden\footnote{Erlanger: finden.} wir nichts Anderes, denn daß GOtt uns für verzweifelte Buben hält. Wie wir denn auch alle sind, halten seine Gebote nicht, keiner ausgenommen; denn GOtt hat uns alle unter die Sünde beschlossen, Röm. 11, 32. Gal. 3, 22.

205. Also sind die zehn Gebote ein Spiegel, darin wir sehen, was wir für Leute sind. Ja, sprichst du, ich will die zehn Gebote wohl halten; ich will niemand das Seine nehmen, will nicht tödten, will niemand bei seinem Weibe schlafen, niemand betrügen, will auch nicht lügen etc. Ei, Lieber, wir reden hier nicht vom äußerlichen Thun oder Lassen, sondern daß du thuest, was
\quellseite{559}
die Gebote GOttes erfordern, oder lässet, was sie verbieten, von ganzem, reinem Herzen. Du sprichst wohl: Ich will Vater und Mutter ehren; ja, du ehrst sie wohl, aber nicht ferner, denn so es dir wohlgefällt, und weil es dir frommt, und [du] Nutzen davon hast. Wenn aber Vater und Mutter verderben, nichts haben, und du reich bist, wenn sie krank sind, und du gesund, dann findet es sich, wenn du sie versorgen sollst mit Kleidern, Essen und Trinken, wenn du ihr Alter, Schwachheit und Unlust an ihnen tragen sollst, da wirst du wohl gewahr, wie du sie von Herzen wirst ehren. Da wird dann erfüllt das gemeine Sprüchwort, daß ein Vater kann leichter zehn Kinder ernähren, denn zehn Kinder Einen Vater.

206. Man sagt ein Exempel von einem Vater, der übergab seinen Kindern alle seine Güter, Haus, Hof, Acker und Wiesen und alle Bereitschaft, versah sich des zu seinen Kindern, sie würden ihn ernähren. Da er nun bei einem Sohne eine Zeitlang war, war der Sohn sein überdrüssig, und sprach: Vater, wann willst du zu meinem Bruder gehen? Da er nun auch eine Zeitlang bei dem andern Sohne gewesen war, ward er sein auch müde, und schickte ihn zur Schwester. Die hatte seiner auch bald genug. Also ward der gute alte Vater ganz unwerth, und übel versorgt. Da er nun sterben sollte, schickte er einen alten Kasten zuwege, und legte einen versiegelten Brief drein, und rief alle seine Kinder zu sich, als wollte er ihnen ein Testament und einen Schatz bescheiden. Da kamen die Kinder, waren alle gehorsam. Sie meinten aber den Vater nicht. Da er nun gestorben war, fanden sie einen Brief, da stund also drinnen: Es soll kein Vater seinen Kindern sein Gut übergeben, oder seinen Stab aus der Hand geben, weil er ihn selbst halten kann. Darum ist es verloren. Wie diese gethan haben, also thun alle.

207. Darum hält keiner dies und andere Gebote; es kann’s auch niemand halten. Denn könnten wir’s halten, so hätte es GOtt nicht geboten; weil er es aber gebeut, so ist es ein Zeichen, daß niemand Vater und Mutter ehrt. Es wird wohl geboten, aber Vater und Mutter werden es wohl gewahr, wie es die Kinder halten. Darum bleibe unverdorben, behalte deine Güter, und siehe den Kindern nicht in das Maul \&c. Man findet wohl Kinder, die Vater und Mutter ehren, weil sie keiner Hülfe bedürfen; wenn sie ihnen aber helfen sollen, da ist niemand daheim, da ist die Ehre denn schon aus. So erkennst du nun, daß von Natur alle Kinder Vater und Mutter ungehorsam sind. Soll nun ein Kind Vater und Mutter ehren vom Grunde seines Herzens, wie dies Gebot erfordert, so muß [es] geschehen durch die Gnade des Heiligen Geistes; die Natur vermag es nicht.

208. Es verdienen jetzt die Eltern so großen Undank an ihren ungehorsamen Kindern, daß etliche sprechen, und meinen, sie thun wohl daran: Der Teufel soll ein Weib nehmen und Kinder zeugen, wenn es also zugeht. Denn, so Vater und Mutter arbeiten und lassen ihnen es blutsauer werden, ziehen in Mühe und Arbeit die Kinder auf, wagen für sie Leib und Leben, was verdienen sie damit? Nichts anders denn, so sie alt werden, daß ihnen die Kinder das Leben nicht gönnen, wollten lieber, daß sie todt wären, daß sie in die Güter kämen. Ist das nicht ein schändlich, verdrießlich Ding, daß die Kinder ihren Eltern nicht das Leben vergönnen? O Welt, du edles Kind! man sollte dich billig krönen, daß du so dankbar bist. Kannst du denen nicht Gutes thun, von denen du alles Gut, ja, das Leben hast? Wenn solltest du denn Gutes thun? Sie haben dir zwanzig, dreißig, vierzig Jahre gedient und dich ernährt, und du kannst ihnen nicht vierzehn Tage dienen und sie ernähren? Pfui dich, du schändliche Welt, kannst du Vater und Mutter nichts Gutes thun, die dich aus dem Dreck erzogen haben, wem willst du es denn thun?

209. Also geht es auch mit andern Geboten: „Du sollst nicht tödten“ \&c. Sprichst du aber: Schlage ich doch niemand todt. Ja, weil du deines Nächsten genießen kannst, thust du nicht wider ihn; wenn er dich aber antastet, da sage an, wie du gegen ihn gesinnt bist; da forsche dein eigen Herz, so wirst du finden, daß du ein Mörder und Todtschläger bist. Denn wenn du deinen Bruder hassest, so bist du ein Todtschläger, wie Johannes spricht 1. Ep. 3, 15. Und ob du ihn\footnote{„ihn“ fehlt in der Wittenberger.} gleich mit der Faust nicht tödtest, so gönnst [du] ihm doch das Leben nicht; geht es ihm übel, so lachst du in die Faust. Das ist uns nun angeboren, die Natur, Vernunft und freie Wille vermag nichts anders denn tödten \&c. Es steckt in allen Menschen ein Mörder, Bluthund; in dir gleichwohl, als in mir, daß also
\quellseite{560}
vor Gottes Angesicht niemand gerecht erfunden wird; ja, wir sind alle gleich Bösewichter, und ob die Faust stille hält, so ist doch der Fehl im Herzen. Denn was vom Fleisch geboren ist, das ist Fleisch [Joh. 3, 6.]. Darum sehen wir fein in den zehn Geboten, als in einem Spiegel, was GOtt von uns hält, nämlich, daß wir Buben in der Haut sind, und sein Urtheil fehlt nicht.

210. So haben wir nun genugsam gehört im vierten Gebote, wie man Vater und Mutter ehren soll, und was dies Gebot in sich halte und lehre, also, daß man wohl greifen kann, daß GOtt viel daran gelegen ist, daß dieser Gehorsam gegen Vater und Mutter im Schwange gehe. Und wo solches nicht geschieht, da sind keine guten Sitten, noch kein gut Regiment. Denn, wo in Häusern Gehorsam nicht gehalten wird, wird man es nimmermehr dahin bringen, daß eine ganze Stadt, Land, Fürstenthum oder Königreich wohl regiert werde. Denn da ist das erste Regiment, davon einen Ursprung haben alle anderen Regimente und Herrschaften. Wo nun die Wurzel nicht gut ist, da kann weder Stamm noch gute Frucht folgen.

211. Denn was ist eine Stadt anders, denn ein Haufen Häuser? Wie sollte denn eine ganze Stadt wohl regiert werden, wo in den Häusern kein Regiment ist, ja, da weder Kind, Knecht noch Magd gehorsam ist? Item, ein ganz Land, was ist es anders, denn ein Haufen Städte, Märkte und Dörfer? Wo nun die Häuser übel regiert werden, wie kann ein ganzes Land wohl regiert werden? Ja, da muß nichts anders draus werden, denn eitel Tyrannei, Zauberei, Morden, Dieberei, Ungehorsam. Denn ein Fürstenthum ist ein Haufen Länder und Grafschaften, ein Königreich ein Haufen Fürstenthümer, ein Kaiserthum ein Haufen Königreiche. Diese alle spinnen sich aus einzelnen Häusern. Wo nun Vater und Mutter übel regieren, lassen den Kindern ihren Muthwillen, da kann weder Stadt, Markt, Dorf, Land, Fürstenthum, Königreich noch Kaiserthum wohl und friedlich regiert werden. Denn aus dem Sohne wird ein Hausvater, ein Richter, Bürgermeister, Fürst, König, Kaiser, Prediger, Schulmeister etc.; wo er nun übel erzogen ist, werden die Unterthanen wie der Herr, die Gliedmaßen wie das Haupt.

212. Darum hat GOtt als am Nöthigsten angefangen, daß man im Hause wohl regiere. Denn wo das Regiment im Hause wohl und rechtschaffen geht, ist dem andern allen wohl geraten. Ursache, denn wir sehen, daß das ganze menschliche Geschlecht daher kommt. Denn es hat GOtt also gefallen, daß aus Vater und Mutter die ganze Welt herkomme. Er könnte aus Stein und Holz Menschen, ja, Kinder Abrahä, wie der heilige Johannes der Täufer spricht, Matth. 3, 9., wohl machen; er will es aber nicht thun, sondern er will, daß eines vom andern komme. Darum schafft er uns auch Kinder, und gebeut ihnen, daß sie ihren Eltern gehorsam seien, und uns, daß wir sie auferziehen und zum Besten halten. Denn was wären wir sonst GOtt nütze, wenn wir solches nicht thäten? Darum hat er uns die Kinder so nahe eingepflanzt, daß er sie nicht aus Stein oder Holz, sondern aus unserm eigenen Fleische und Blute spinnt; daß ja die Ehre und Gehorsam der Kinder gegen den Eltern, und die Sorge, Mühe und großer Fleiß der Eltern gegen den Kindern desto herzlicher und williger geschehe. Wenn wir nun die Kinder nicht\footnote{„nicht“ fehlt in der Wittenberger.} wohl regieren, ziehen und lehren, die aus unserm Fleische und Blute kommen, wie wollten wir uns ihrer annehmen, wenn sie aus Stein oder Holz herkämen.

213. Darum sehen die Eltern zu, daß sie nach GOttes Befehl den Kindern wohl vorstehen, und flugs in der Erste dazu thun, weil sie sich noch ziehen, biegen und leiten lassen, und nicht harren, bis sie erwachsen und in ihrem Muthwillen erhärtet worden, oder erharret,\footnote{„oder erharret“ fehlt in der Wittenberger und in der Jenaer.} bis sie andern Leuten in die Hände kommen. Denn man darf nicht denken, daß fremde Kinder so nahe zu Herzen gehen, als die eigenen; und ob es schon zu Zeiten geschieht, so geschieht es doch gar selten, also, daß unter hundert Kindern kaum eins ist, deß man sich so herzlich annimmt, als wäre es ein eigen, leiblich Kind. So nun die Kinder sich nicht selber regieren können (wie wir sehen, wenn Vater und Mutter den Kindern entfallen, wie sie so elend und weißlos\footnote{In der Jenaer: „weißlos“ = ohne Unterweisung. „weißlos“ wird wohl so viel sein als: verwahrlost. Erstere Lesart scheint uns den Vorzug zu verdienen. (Siehe § 214: „niemand, der sie recht unterweise“.)} hergehen, niemand sich ihrer recht annimmt), so hat GOtt Noth halben\footnote{Erlanger: auch vonnöthen.} so ein hartes Gebot darauf legen müssen.

\quellseite{561}
214. Wir sehen zwar auch, daß ein größerer Fehler und Mangel an den Eltern ist denn an den Kindern, daß sie ungehorsam und ungezogen sind. Denn die Eltern sind nachlässig, thun keinen Fleiß bei den Kindern; und solche Eltern sind nicht werth, daß ihnen ihre Kinder wohl gerathen. Nun, das Gebot steht wohl da, daß die Eltern ihre Kinder in der Furcht GOttes aufziehen sollen; wo sind aber die Eltern, die solches thun? Das thun sie wohl, daß sie nach dem Laufe dieser Welt die Kinder lieben und aufziehen, wie sie sich in die Welt sollen schicken; aber nach der Seele, in der Furcht GOttes ist niemand, der sie recht unterweise und lehre. Man sehe nur darauf, wie man sich zur Sache stellt, was man durch und durch im ganzen Lande für Schulen hält. Niemand ist, der seine Kinder recht lehrt beten, und die Stücke, so zur Seligkeit gehören. So will auch niemand so viel daran wagen, daß seine Kinder durch andere Leute erzogen, gelehrt und unterwiesen würden.

215. Es sind etliche Thiere, die ihre eigenen Jungen fressen, und verderben ihre eigenen Früchte. Also sind auch solche Menschen, die ihre Kinder nicht lehren und unterweisen. Ja, es ist kein Thier auf Erden, das gegen seine Jungen so hart ist, wie ein Mensch, wenn wir es nach der Seele wollen ansehen. Darum wären sie wohl werth, wenn GOtt nicht so fromm wäre und die Eltern vor den Kindern vertheidigte, daß die Kinder die Eltern über die Köpfe schmissen, ja, gar zu Tode schlügen, daß sie so wenig Achtung auf sie haben, und nicht wohl ziehen und unterweisen.

216. Daß aber das Regiment der Eltern noch ein wenig geht, und die Kinder noch ein wenig ihnen gehorsam sind, das kommt wahrlich nicht aus unserm Verdienste, sondern ist eitel Gnade GOttes, die ihnen solches zum Ueberfluß schenkt, und, wie andere Gaben, in die Kappusie wirft; sonst sollte es wohl alles zu Trümmern gehen; wie denn der Teufel Lust dazu hat, daß er es alles auf einen Haufen werfe, und in einen Kloß schlage.

217. Die Kinder wachsen auf, sind ihrer selbst nicht mächtig, sondern sind unter der Hand der Eltern; thun die nicht das Beste bei ihnen, so wird nichts Gutes aus ihnen. Was ist denn die Schuld? Ist es der Kinder? Was können sie dazu! Darum sollen Vater und Mutter über ihnen halten, sie mit Fleiß ziehen, unterweisen und lehren; nicht allein nach weltlicher Weise, sondern auch in geistlichen Sachen, die der Seelen Seligkeit belangen.

218. Was soll ich aber viel\footnote{„viel“ fehlt in der Erlanger.} sagen? Wie sollen sie die Kinder viel lehren und unterweisen in geistlichen Sachen, so sie doch sie\footnote{„sie“ fehlt in der Wittenberger und in der Jenaer.} nicht lehren recht äußerlich vor der Welt leben, haushalten etc. Es geschiehet doch nicht, denn es ist ein Gebot GOttes. Der Teufel hat die Herzen gefangen, führt sie wie er nur will; darum bleibt auch die Welt voller Wölfe, Bären, Löwen und unvernünftiger, wilder Thiere.

219. Aber die rechtschaffenen Menschen sehen wohl, was GOtt will angezeigt haben mit dem Regiment und Gehorsam der Eltern, nämlich, daß Vater und Mutter in ihren Häusern Bischöfe, Päpste, Doctores, Kaiser, Fürsten und Herren sind. Darum soll ein Vater sein Kind wie ein Richter strafen, lehren wie ein Doctor, ihm vorschreiben wie ein Pfarrherr oder Bischof. Thut nun ein Vater solches, so kann er vor GOtt bestehen; thut er es nicht, so wird er seinen Lohn von GOtt wohl bekommen zu seiner Zeit. Denn, wie droben [§ 184] gesagt ist, Vater und Mutter können das Himmelreich verdienen an ihren Kindern. Also wiederum, mögen die Eltern nicht leichter die Hölle verdienen, denn an ihren eigenen Kindern, in ihrem eigenen Hause, wo sie dieselben versäumen und nicht lehren die Dinge, wie gesagt ist. Was hülfe es, wenn sie sich gleich zu Tode fasteten, beteten, und alle Werke thäten, und ließen doch unterwegen, was ihnen von GOtt befohlen ist. GOtt wird sie von diesen Stücken nicht strafen am jüngsten Tage, sondern von den Dingen, die er ihnen befohlen hat.

220. Man soll aber wohl merken, wie in andern Geboten gesagt ist, daß dies Werk muß auch im Hauptwerke gehen, nämlich, im Glauben. Denn Zucht und Lehre, dazu man die Kinder hält, ist an ihm selber nichts vor GOtt, es sei denn, daß es geschehe im Glauben, daß der Mensch nicht daran zweifle, es gefalle GOtt wohl, daß die Eltern die Kinder zum Besten halten, und lasse ihm solch Werk nicht anders sein, denn eine Vermahnung und Uebung seines Glaubens, GOtt zu vertrauen, und [sich] alles
\quellseite{562}
Gutes zu ihm zu versehen, ohne welchen Glauben kein Werk lebt, gut und angenehm ist. Denn viel Heiden haben ihre Kinder hübsch und ehrlich vor der Welt erzogen; es ist aber alles verloren gewesen um des Unglaubens willen, wie Paulus spricht Röm. 14, 23.: „Was nicht aus dem Glauben gehet, das ist Sünde“; und wiederum: „Die GOtt lieben, denen ist es alles gut“, Röm. 8, 28.

221. Das ist nun ein herrlich Gebot. Wiewohl es klein ist in Worten, so ist es doch kräftig in der That, denn die ganze Welt wird in dem Gebote regiert. Denn wo dies Regiment der Eltern hinweggenommen würde,\footnote{Erlanger: hinweg ist.} so wäre es mit der ganzen Welt geschehen; denn ohne das Regiment kann sie nicht bestehen. Wer nun in diesem Regiment ist, dem ist befohlen ein groß Amt. Das kann der heilige Paulus fein ausstreichen und hoch anziehen, da er [1 Cor. 11, 7.] spricht: „Das Weib ist des Mannes Ehre.“ Denn Gott hat dem Manne ein Amt befohlen, sein Weib und Kinder zu regieren; nicht darum, daß sich deß\footnote{In der Jenaer und in der Erlanger: „sichs“ = sich des. Wittenberger: sich.} der Mann überhebe und ein Wohlgefallen darinne habe, sondern, daß es ein trefflicher Befehl ist von der göttlichen Majestät, darum ist der Mann ehrenreich.\footnote{In der Wittenberger und in der Jenaer: ehrenreich.} Wer nun nicht Weib und Kind hat, der hat diese Ehre und das Amt nicht. Also spricht er auch vom Manne, daß er Gottes Ehre sei, das ist, daß er unter GOttes Gewalt sei, und wird von GOtt regiert. Wie nun der Mann unter GOtt ist, so ist das Weib unter dem Manne. Also lobt und preist die Schrift überall den ehelichen Stand; bei uns aber ist er verächtlich. Darum sind wir Säue, und bleiben Säue, und sehen nicht, was GOtt wolle und gebiete. Also ist gehandelt, wie sich die Eltern gegen ihren Kindern, und die Kinder gegen ihren Eltern halten sollen.

222. Von den Eltern kommt das Regiment auf die weltliche Obrigkeit. Denn, wie die Eltern daheim im Hause Gewalt haben über ihre Kinder und Hausgesinde, also hat die Obrigkeit Gewalt über eine ganze Gemeinde. Die Eltern schaffen Recht und Frieden im Hause; die Obrigkeit schafft Frieden und Recht in einer ganzen Gemeinde und an allen Orten. Darum ist auch der Obrigkeit Amt, daß sie Vater sein soll, und allen Menschen wohlthun, und Gute erzeigen, auch niemand weder Gewalt noch Unrecht geschehen lassen. Darum sie auch in der hebräischen Sprache den Namen haben Nedibim [\hebraeisch{נדיבים}], von der Wohlthat; und im 51. Psalm, V. 14., wird der Heilige Geist genannt ein fürstlicher Geist, das ist, welcher nichts denn wohlthun kann, wie der Fürsten Amt sein soll; und Christus [Luc. 22, 25.] heißt sie beneficos, gnädige Herren, auch vom Wohlthun. Wie sich aber weltliche Obrigkeit halten solle gegen ihre Unterthanen, und die Unterthanen gegen ihre Obrigkeit, ist genugsam angezeigt im Büchlein von weltlicher Obrigkeit.\footnote{St. Louiser Ausgabe, Bd. X, 374.}

223. Hierher gehört auch der Gehorsam des Gesinds, der Werkleute und Tagelöhner gegen ihren Herren, Frauen, Meister und Meisterinnen, davon Paulus an Titum Cap. 2, 9. Col. 3, 22. und St. Petrus 1. Ep. 2, 18. sagen.

\section[Das fünfte Gebot, oder das andere Gebot der andern]{Das fünfte Gebot, oder das andere Gebot\footnote{„Gebot“ fehlt in der Erlanger.} der andern Tafel}

\begin{schriftzitat}
V. 13. Du sollst nicht tödten.
\end{schriftzitat}

\subsection{Summa dieses Gebots}

224. Das andere Gebot der andern Tafel lehrt, wie man sich soll halten gegen seinem Gleichen oder Nächsten, seiner eigenen Person halben, daß man dieselbige nicht beleidige, sondern, wo sie es [be]darf, fördere und helfe.

\begin{schriftzitat}
Du sollst nicht tödten.
\end{schriftzitat}

225. Das erste Gebot in der andern Tafel gehört dahin, daß man Vater und Mutter ehre, ist also das allernächste auf die Gebote, die GOtt betreffen; und geht darum auf die erste Tafel, denn in dem vierten Gebote ist begriffen die Obrigkeit, welche an GOttes Statt ist, wie Paulus Röm. 13, 2. anzeigt: „Welcher der Obrigkeit widerstrebet, der widerstrebet GOttes Ordnung.“ In den ersten dreien Geboten wird gesündigt wider GOtt; im ersten Gebot der andern Tafel sündigt man wider die Obrigkeit, welche GOtt selbst geordnet und eingesetzt hat, 1 Mos. 9, 6., und Christus und Paulus haben dieselbe nicht aufgehoben, sondern bestätigt.

\quellseite{563}
226. Nun, die Gebote, die hernach gehen, gehören nicht auf die Obrigkeit, sondern gegen unsern Gleichen, und gegen unsern Nächsten. Es fehlt und mangelt uns allen daran, daß wir dies, wie andere Gebote, verachten, und glauben nicht, daß es GOtt rede, und GOttes Ernst sei; achten es, als hätte es ein loser Schwätzer gesagt. Darum bleiben es wohl Gebote, denn der natürliche Mensch versteht ja nicht die Dinge, so den Geist betreffen [1 Cor. 2, 14.]. Darum hält kein Mensch die Gebote, er sei denn ein Christ und erleuchtet durch den Geist Gottes.

227. So lautet nun das fünfte Gebot also: „Du sollst nicht tödten“, und geht auf unsers Nächsten Person. Da sehen wir aber einmal, was GOtt von uns hält, wie viel Gutes er sich zu uns versieht, was er im Sinne hat mit diesem Gebote: „Du sollst nicht tödten.“ Er gedenkt also: Ich habe wilde, unvernünftige, tolle, rasende Thiere in der Welt, Wölfe, Bären, Löwen etc., darum muß ich sie verwahren, verriegeln, mit eisernen Gittern vergittern und mit starken Mauern verschließen, daß sie sich nicht unter einander würgen und großen Schaden thun. Denn wenn GOtt die Sorge nicht hätte, was dürfte er die Gebote geben? Also erkennt GOtt unser Herz und Natur aus der Maßen wohl, daß uns der Mord in das Fleisch eingeboren sei; darum gibt er auch dies Gebot, daß wir uns selber sollen erkennen; hat Sorge, wir würgen uns unter einander, wie die tollen, rasenden Hunde, Wölfe und Bären; hält uns also für verzweifelte Buben, die sich unter einander würgen und morden. Und die Historie, die Moses hernach beschreibt nach Adam, ist vom Mord und Todtschlag, wie ein Bruder den andern würget.

228. Nun gehe hin, lieber Mensch, rühme dich, wir wollen heilig sein. [Wir] rühmen uns unsrer Vernunft, Weisheit und des freien Willens; was hält aber GOtt von uns? Das hält er von uns, daß wir alle Mörder und Todtschläger sind, keiner ausgenommen. Und thut GOtt gleich wie ein Bürgermeister oder ein Fürst, wenn derselbige hörte, daß etliche Dräuworte vorgeben, sie wollten Schaden thun, bei der Nacht einbrechen, so ließe er wachen und hüten, damit ihnen gewehret würde. Also versieht sich GOtt auch nichts Gutes zu uns, sondern hält uns alle für Todtschläger; darum gebeut er uns, wir sollen nicht tödten.

229. Das ist aber verdrießlich, daß GOtt in die Gemein redet, und nimmt niemand aus, läßt das Urtheil über alle Menschen gehen, über Fromme und Böse, Arme und Reiche, hohes und niedriges Standes, es\footnote{Erlanger: er.} sei Fürst, Herr oder Knecht [Röm. 11, 32. Gal. 3, 22.]. Er möchte doch gesondert haben der heiligen Leute, der Pharisäer. Aber er schont niemand, er nimmt keinen Menschen aus, trifft sie alle, und wirft sie alle in einen Haufen. Als wollte er sagen: Sie sind allzumal Buben, Mörder und Todschläger, es ist keiner, er hat einen Bluthund im Busen. Das ist nun ein kurzer Beschluß: wie wir allzumal ungehorsam Vater und Mutter sind, also sind wir auch allzumal Mörder. Da lerne nun erkennen, was die Welt für ein Kind ist; welches die zehn Gebote gar fein weisen, darin du, als in einem Spiegel, siehst, wie wir außen und innen sind.

230. Du sprichst aber:\footnote{Die Worte: „Ich will ... todtschlagen“ fehlen in der Wittenberger.} Ich will niemand tödten, morden, noch todtschlagen; ich wollte auch einem Kinde nicht gerne ein Leid thun. Ei, Lieber, wenn man thut, was du gerne wolltest, weil man dich freundlich anlacht und dich lieber Freund heißt, so bist du wohl fromm; wenn man dir aber Unrecht thut, wenn man dich überrumpelt, mit dir zürnt, da findet es sich, da wird man deine Geduld wohl sehen, du solltest bald einen in das Maul schlagen, ja, erwürgen, wenn dir es ohne Gefahr wäre. Was ist die Ursache? Deine böse Natur, darinnen du geboren bist, dieselbe sieht GOtt in dir, der ein Erforscher der Herzen ist [1 Chron. 29, 9.], darum hat er dies Gebot gegeben. Wenn du aber denn ein süßes Herz könntest haben gegen deinem Feinde, und ihm alles Gutes gönnen, und nichts Arges, wenn er dein Herz erzürnt und erbittert, so möchtest du dich berühmen: „Ich tödte nicht. Wo ist aber einer, der es thut? Keiner thut es. Denn so ferne kannst du deinem Nächsten wohl freundlich sein, so ferne er dir keinen Schaden thut. Thut er dir Schaden, so ist deine Freundschaft gegen ihm schon aus, und kannst dich nicht enthalten, daß du wider dies Gebot nicht sündigest. Du tödtest ihn; thust du es mit der Faust nicht, so thust du es doch mit dem Herzen.“

\quellseite{564}
\subsection{Auslegungen über das zweite Buch Mosis}

231. Darum ist es nicht genug, daß du äußerlich kein Mörder bist. Denn, ob die Hand gleich still hält, so ist doch der Grund des Herzens vergiftet, denn du gönnst deinem Feinde nichts Gutes, lachst in die Faust, wenn es ihm übel geht, wenn er krank ist, verdirbst oder stirbt; und wenn es ihm wohl geht, gönnst du es ihm nicht. Das alles heißt morden und todtschlagen. Und diese hübsche Tugend steckt allen Menschen von Natur in ihrem Herzen.

232. Darum, wie ich gesagt habe, sind uns die zehn Gebote gegeben, daß wir uns selber sollen erkennen lernen, was wir für fromme Kinder, und unsere Herzen für Kräutlein sind, und was GOtt von uns halte, auf daß wir uns in unsere Herzen schämen und erschrecken vor uns selber, wenn wir in den zehn Geboten, wie in einem Spiegel, sehen, daß wir alle abgöttische Menschen, Lästerer göttliches Namens, ungehorsam unsern Eltern, Mörder und Bluthunde sind.

233. Christus legt dies Gebot aus, Matth. 5, 21. 22., da er also spricht: „Ihr habt gehöret, daß zu den Alten gesagt ist: Du sollst nicht tödten. Wer aber tödtet, der soll des Gerichts schuldig sein. Ich sage aber euch: Wer mit seinem Bruder zürnet, der ist des Gerichts schuldig.“ Da zeigt Christus an, daß du wider dies Gebot sündigst, nicht allein, wenn du mit der Faust einen erwürgst, sondern auch, wenn du mit deinem Nächsten zürnst. Und Christus macht des Zorns etliche Grade und Unterschiede. Am ersten soll man nicht zürnen im Herzen; zum andern den Zorn mit Zeichen und Gebärden nicht erzeigen; zum dritten mit Worten nicht schelten; und zum vierten mit der Faust still halten.

234. Da ist niemand ausgenommen, sondern wir werden alle schuldig. Denn wiewohl die Faust, das Maul und andere Glieder still halten, so ist doch das Herz voll Zorn, Neid und Haß. Das alles findet sich bei uns. Es sehe nur ein jeglicher in sein Herz, so darf man ihm nicht viel Schrift anzeigen, sein eigen Gewissen wird ihm wohl sagen und ihn überzeugen, daß er hierinnen schuldig ist. Man sieht es auch zwar äußerlich an unsern Früchten und Wandel wohl, was wir für Gesellen sind; da kann keiner dem andern ein Wort vorbrechen; wie soll er denn seinem Feinde von ganzem Herzen freundlich sein?

\section{Das dritte Gebot der andern Tafel}

\begin{schriftzitat}
V. 14. Du sollst nicht ehebrechen.
\end{schriftzitat}

\subsection{Summa}

235. Das dritte Gebot der andern Tafel lehrt, wie man sich halten soll gegen des Nächsten höchstes Gut nach seiner eigenen Person, das ist, sein ehelich Gemahl, Kind oder Freund, daß man dieselbe nicht schände, sondern bei Ehren behalte.

\begin{schriftzitat}
Du sollst nicht ehebrechen.
\end{schriftzitat}

236. Das ist aber eine große und schöne Ehre, die GOtt der Welt zugewiget, nämlich, daß sie ein Stall voll Ehebrecher und Ehebrecherinnen ist! GOtt verdient es wohl an uns, daß wir ihm feind würden, weil er uns also schändet, höhnt und schmäh[t], und niemand dazu ausnimmt, auch unsere Mönche nicht, wenn sie gleich noch einmal Keuschheit gelobt hätten. Da siehst du, daß uns GOtt nicht vertraut, daß ein Ehemann wäre, der sich an seinem Eheweibe,\footnote{Erlanger: Weibe.} ließe benügen, und wiederum, ein Eheweib, die sich an ihrem Manne ließe benügen. Es sollte ja noch einen verdrießen, wenn man ihn einen Ehebrecher schälte, und spräche zu ihm: Lieber, halte deine Ehe, gehe keinem andern zu seinem Weibe, schände keinem seine Tochter, und so fort an;\footnote{Wittenberger: und so fort an. Erlanger: zufortans. Jenaer: zuvor aus.} wenn es ein frommer Mann wäre. Er sollte bald sprechen: „Wofür siehst du mich an, hältst du mich für einen solchen? Aber GOtt schont keines nicht, schilt uns alle zugleich in diesem Gebote Ehebrecher und Ehebrecherinnen; will damit anzeigen, was wir für Gesellen sind.

237. Also ist dies Gebot auch ein Lasterbüchlein und ein Schandtitel, schilt uns alle, niemand ausgenommen, daß wir Huren treiber sind; ob wir es gleich nicht vor der Welt öffentlich sind, so sind wir es doch im Herzen, und wo wir Raum, Zeit, Statt und Gelegenheit hätten, brächen wir alle die Ehe. Die Art ist allen Menschen eingepflanzt, es wird keiner ausgenommen, er heiße Mann oder Frau, Alt oder Jung, sie liegen allzumal in diesem Spitale krank. Und diese Seuche hängt uns nicht an wie ein rother Rock, daß wir es könnten ausziehen
\quellseite{565}
\footnote{Erlanger: ausschlagen.} oder weglegen, sondern wir haben es aus Mutterleibe gebracht, und ist uns durch Fell und Fleisch, Mark und Bein, und durch alle Adern durch und durch gezogen. Sind doch viel, die nicht huren, sondern ein fein Leben führen? Ei, Lieber, ich sage nicht vom Thun, sondern von der Art. GOtt läßt sich nicht mit den Werken äffen; die Schrift nennt ihn einen Herzenserkenner [1 Kön. 8, 39.], er sieht tiefer denn wir.

238. Daß nun etliche fromm sind, und wider dies Gebot nicht sündigen, das richtet seine göttliche Gnade aus, oder Meister Hans mit dem Schwert und Staupebesen,\footnote{Wittenberger: staupe und besem.} der\footnote{„der“ fehlt in der Erlanger.} treibt ihnen eine Furcht ein, daß sie öffentlich solche Sünde vermeiden; wo solches nicht wäre, so würden wir wohl an Tag geben, was uns im Herzen steckt, und also leben, wie etliche Heiden, die Hurerei außerhalb der Ehe nicht gestraft haben. So spricht nun GOtt: „Du sollst nicht ehebrechen“, sondern sollst keusch sein; als wollte er sprechen: Ihr seid alle der Art, daß ihr Hurer und Ehebrecher seid, einer wie der andere.

239. Weil nun die Art in uns ist, so hat GOtt zugegeben einem jeglichen, ein Weib zunehmen, und daß eine jegliche Frau ihren Mann habe, damit Hurerei und Ehebruch vermieden würde [1 Cor. 7, 2.]. Denn dies Laster ist ganz weitläufig in der ganzen Welt, wie ein groß Feuer, hilft schier weder Schwert noch Geist. Darum ist sehr vonnöthen, daß man solche Sünde äußerlich strafe und nicht gestatte, zu einer Warnung den andern. Und darauf sollen die Regenten fleißiglich sehen, und hierinnen nicht versäumig sein.

240. Und zwar unsere Vernunft sagt uns, daß Hurerei, Ehebruch und andere Sünden unrecht seien. Denn einem jeglichen sagt sein eigen Herz, er solle Vater und Mutter ehren, weil wir von ihnen herkommen, und haben Fleisch und Blut von ihnen; [sie] wagen um unserwillen Leib und Leben, strecken alle ihr Hab und Gut für uns dar. So auch lehrt\footnote{In den Ausgaben: lernet.} uns die Natur, daß wir niemand morden sollen, nicht unsere eigenen Richter sein. Und Summa, in der Natur sind alle diese Gebote verfasset; aber wir sind so toll und so\footnote{„so“ fehlt in der Jenaer.} voll Blindheit, daß wir sie nicht sehen noch erkennen. Also lehrt uns auch die Natur, daß wir niemandes Weib oder Kinder sollen schänden. Denn die Natur sagt: Was du willst, daß man dir nicht thue, das sollst du andern Leuten auch nicht thun [Matth. 7, 12.]. Nun will keiner, daß man ihm Weib und Kind schände, darum so lasse deinem Nächsten sein Weib und Kind auch ungeschändet.

241. Wenn man aber in die Brunst kommt, so vergißt man alles, Gesetz, Natur, Schrift, Bücher, GOttes und seines Gebots, ist nur lauter Gesuch da, böse Lust zu büßen.

\section{Das vierte Gebot [der andern Tafel]}

\begin{schriftzitat}
V. 15. Du sollst nicht stehlen.
\end{schriftzitat}

\subsection{Summa}

242. Das vierte Gebot der andern Tafel lehrt, wie man sich soll halten gegen des Nächsten zeitliches Gut, daß man es nicht nehme noch hindere, sondern fördere.

\begin{schriftzitat}
Du sollst nicht stehlen.
\end{schriftzitat}

243. Da kommt wieder ein groß Register, aus welchem niemand ausgenommen wird; der geringste Theil der Diebe wird mit dem Strange gehängt. Denn wenn man alle Diebe hängen sollte, die in der Welt sind, wo wollte man Stricke genug nehmen? Man könnte ihrer nicht genug bekommen, es müßten alle Gürtel und Riemen zu Strängen gemacht werden.

\subsection{Die Ordnung der Gebote der andern Tafel}

244. Nun merke die Ordnung der Gebote der andern Tafel. Das erste geht die Obrigkeit an. Das andere die Person des Nächsten: Du sollst deinen Nächsten nicht verletzen an seiner eigenen Person. Das dritte gehört die Personen an, die deinem Nächsten zugehören, als: Du sollst nicht schänden deines Nächsten Hausfrau, Tochter, Schwester, Magd etc. Das vierte geht die Güter deines Nächsten an, daß du sie ihm nicht entwendest, noch verrückest.

245. Dies Gebot zeigt uns wieder einmal, was GOtt von uns halte, nämlich, daß wir allesamt Diebe sind, keinen ausgenommen, vor GOtt und vor der Welt. Daß aber der Diebstahl etlichermaßen nachbleibt, ist das Ursache, daß man sich vor dem Henker und dem
\quellseite{566}
Galgen fürchtet; auch, wo GOttes Gnade und der heilige Geist wohnt, wird wider dies Gebot nicht gesündiget; sonst liegt der Dieb begraben im Herzen, und wo GOtt nicht hält, oder der Henker nicht schreckt, so geht der Dieb heraus in das Werk, auf das allergröbste.

246. Du sollst aber nicht meinen, daß das allein gestohlen heiße, wenn du deinem Nächsten das Seine ausführst, sondern, wenn du siehst deinen Nächsten Noth, Hunger, Durst leiden, keine Herberge, Schuh und Kleider haben, und hilfst ihm nicht, so stiehlst du gleich sowohl, als wenn einer dem andern das Geld aus dem Beutel oder Kasten stähle; denn du bist ihm schuldig zu helfen in seiner Noth. Denn deine Güter sind nicht dein, du bist allein ein Schaffner darüber gesetzt, daß du sie austheilest denen, so es bedürfen. Darum gehören die auch an den Reigen\footnote{So die Jenaer. Wittenberger: reien; Erlanger: Reiben.} und in dies große Register, darin die Diebe geschrieben stehen, so Güter haben, und geben nicht denen, so es bedürfen, nehmen sich\footnote{Erlanger: sie.} ihrer Nächsten Nothdurft nicht an, gehen vorüber.

247. Darum ist selten ein reicher Mann, der in dem Fall nicht ein Dieb, ja, ein großer Dieb ist; es müßte auch ein großer Strick sein, daran man solche großen Diebe sollte hängen. Solcher Diebe aber sind nicht allzuviel; denn der gemeine Mann hat nicht so große Güter, als die reichen, großen Hansen, die nur zu sich scharren, schinden, schaben, und schatzen die Armen, sammeln große Schätze. Dies sind die rechten und größten straßbrüderischen Diebe, die hängt man nicht an Galgen, sondern werden geehrt von jedermann, sitzen obenan; sie werden aber dem Galgen nicht entlaufen, der Teufel wird selbst Meister an ihnen werden, der wird sie in der Hölle recht anbinden, daß sie ihm nicht entkommen werden.

248. Die kleinen Diebe aber, so kaum zehn Groschen stehlen, die müssen flattern. Das haben auch die Römer wohl gesehen, welche gar weise und kluge Leute waren. Denn einer aus ihnen, Cato genannt, spricht: Die kleinen Diebe hängt man an Galgen, aber die großen Diebe gehen in mardern Schaub(en) herein. Item, ein anderer weiser Mann hat gesagt: daß das Recht gleich sei einem Spinnengewebe; wenn die kleinen Fliegen darein kommen, so müssen sie herhalten, wenn aber die großen Hummeln darein kommen, so fahren sie hindurch, zerreißen mit Gewalt das Spinnengewebe. Und ist auch wahrlich also; denn so ein Armer kaum fünf Groschen gestohlen hat, so muß er hangen. Nur hin an Galgen mit ihm, da hilft keine Bitte,\footnote{Wittenberger und Erlanger: Gebet.} da ist das Recht strenge. Die aber Tag und Nacht schinden, schaben, als die Wucherer, trügen und lügen, böse Waare verkaufen, falsche Maße geben, die heißt man noch Gnad-Henker dazu, welche wohl dreimal mehr verdient hätten, daß man sie an den Galgen hängte, denn die geringen Diebe. Doch sie werden wohl gehängt werden, und schändlicher, denn wenn sie von dem Henker angeknüpft würden.

249. Item, die in einer Stadt oder Gemeinde sitzen, als Bierbrauer, Weinschenken, Fischer, Fleischer, Schneider, Schuster und alle Handwerksleute etc. [Es] ist gar selten einer unter ihnen, der nicht einen Dieb im Busen trage. Der gibt falsche Maß, der andere hat falsch Gewicht, Elle etc.,\footnote{„Elle etc.“ allein in der Jenaer.} der betrügt seinen Nächsten sonst, der andere so. Mit seiner falschen, geringen Waare treibt ein jeglicher seinen Muthwillen, wie er nur will. Und ob sich\footnote{Wittenberger und Erlanger: sie.} etliche enthalten, so ist doch die Art da; wo sie also könnten oder dürften, führen sie auch heraus. Also ist es mit den Kaufleuten auch. Es ist ihrer keiner, der dem Diebe könnte entlaufen. Sehet doch, ob ihrer einer auch ein Gewissen oder eine Vernunft habe? Sie verkaufen ihre Waare nach ihrem Gefallen, und wollen es noch für recht haben; sprechen noch dazu: Es ist mein Gut, ich mag mit handeln, wie ich will. Lobe dich, Ruckuf, mit deinem Gesange; man hört am Geschrei wohl, was du für ein Vogel bist.

250. Es können’s auch die Bauern wohl, die zu Märkte kommen. Bringt einer faule, stinkende Eier oder garzige Käse zu Markte, der macht sich stolzer damit, denn etwa einer mit einem Kram, der etliche hundert Gulden werth ist. Ein anderer kommt mit anderer Hümpelei, und kann sich nicht ekel genug damit machen. Niemand gedenkt, daß solches Dieberei sei; ja, es meint einer, es sei eine Kunst, mit Vortheil wissen umzugehen in unserer Handthierung, und also unsern Nächsten [zu] betrügen. Du darfst aber nicht gedenken
\quellseite{567}
, daß du damit wolltest entschuldigt sein; du bist gleich sowohl ein Dieb, wo du solcherücke gebraucht, als hättest du es einem aus dem Beutel gestohlen. Dazu weißest du, daß du nicht recht handelst; denn dein Gewissen straft dich. Und willst doch kein Dieb sein; ja, du lässest dich von andern Leuten für einen frommen, redlichen und aufrichtigen Mann schelten, gibst unterweilen einem armen Menschen einen Heller um GOttes willen, damit soll es bezahlt sein. Es wird sich aber wohl anders finden. Du magst wohl ein gut Lob und Namen vor der Welt haben; siehe aber darauf, wie du vor GOtt daran bist. Da sollen Fürsten und ihre Regenten zusehen, daß man darinnen ein Mittel treffe etc.

251. Item, mir sehen wohl, wie es zuging in der theuren Zeit, wie man mit dem Korn handelte; jedermann gab es nur, wie es ihm gefiel, so mußten es auch denn die Armen so theuer erkaufen, wie sie es boten, denn sie mußten ja essen. Das hält man für keinen Diebstahl. Ja, solche Verkäufer rühmen sich\footnote{Erlanger: sie.} noch dazu, und sprechen: So viel habe ich jetzt gewonnen. Lieber, was du daran gewinnst, schmiere die Schuhe damit. Item, solche Gesellen, wenn man einen Dieb ausführt zum Galgen, lachen in die Faust, und sprechen: Es geschieht ihm recht; vergessen ihrer selber, daß sie größere Diebe sind, denn derselbe, und besser hängen verdient hätten, denn er.

252. Das ist nun vom groben, äußerlichen Diebstahl gesagt, den man greifen kann, und täglich vor Augen sieht; also, daß gar selten ein Mensch ist,\footnote{Wittenberger und Erlanger: sei.} der nicht also einen Dieb im Busen trage, und vor der Welt ein Dieb sei. Und dies Laster ist gemeiner und läufiger in der Welt, denn Ehebruch oder Mord. Denn weil man den Mord und Ehebruch straft, gehen sie nicht also sehr im Schwange, als solch Wuchern und Scharren. Denn es mag jedermann das Seine verkaufen, wie er will, ohne alle Strafe; darum so geht es über und über, und des Betrügens ist kein Ende. Aber die Obrigkeit sollte hier drein sehen, daß es recht zuginge, und die strafen, so die Armen also übersetzen. Aber die solches strafen sollten, sind hierin selbst sträflich; darum fragt keine Krähe der andern ein Auge aus.

253. Nun, GOtt läßt solches wohl geschehen, er kann es aber wohl wieder hereinbringen, nämlich, kann wohl ergotzen die, so betrogen sind worden, und die strafen, so es mit Unrecht erschunden haben, mit ihrer Nächsten Schaden. Denn wenn einer lange gescharrt und gesammelt hat, so kommt ein anderer, der es verzehrt, und unnützlich umbringt. Es geschiehet auch also, wie man vor Augen sieht und erfährt. Wie oft kommt es, daß einer so reich ist, daß er seines Guts kein Ende nicht weiß, und zerstäubt und zerfliegt doch so bald, daß man nicht weiß, wo es hinkommt. Daher auch ein gemein Sprüchwort aus solcher täglichen Erfahrung gekommen ist: De male quaesitis non gaudet tertius haeres, böse gewonnen Gut kommt an den dritten Erben nicht; und: Male partum male disperiit,\footnote{In den Ausgaben: disperit.} übel gewonnen, übel zerschlungen.

254. Also stehst du, daß uns GOtt alle für Diebe hält; sind wir es alle äußerlich nicht, so fehlt es doch am Herzen nicht. Denn wer seines Nächsten Gut begehrt, das ist vor GOtt gestohlen, wie die letzten zwei Gebote ausweisen. Daß wir aber nicht alle herausfahren mit der That und stehlen, das wehrt Junker Henker mit dem Stricke; der ist ein feiner Mann, macht viel fromme Kinder. Denn viel mehr stehlen nicht, aus Furcht des Galgens, denn aus Liebe und Frömmigkeit.

255. So ist nun das nicht allein ein Diebstahl, so man zu Nachts in die Häuser bricht, und einem das Seine heimlich bei Nacht nimmt, sondern ein jeglicher Abbruch oder Schade, der da geschieht deinem Nächsten am Gute, es sei mit Kaufen, Verkaufen, mit Arbeiten, unrechtem Handel, also, daß der Diebstahl ein gemein Laster in der ganzen Welt ist, und welche es nicht können zuwege bringen mit der That, die thun es doch mit den Gedanken; welche aber Raum haben, die thun es auch mit der Faust. Heutiges Tages wird viel Betrug mit Kaufen und Verkaufen gebraucht von allen Handwerkern, diejenigen werden alle in großer Dieberei erfunden. Wir können nicht mehr, denn daß wir dawider predigen; die weltliche Obrigkeit aber soll solches wehren, und mit Ernst dazu thun. Doch, das Herz bleibt gleichwohl voll Diebstahl, es sei denn, daß GOtt das Herz sonderlich erleuchte. Der Richter wird es nicht zwingen; GOtt aber, der
\quellseite{568}
ein rechter Richter ist [Ps. 7, 12.], kann das sein vergelten, wie oben gesagt ist.

256. Wie soll man ihm aber wehren? Dem Herzen kann niemand wehren, denn allein GOtt durch seine Gnade. Aber äußerlich wäre es gut, und ja wohl vonnöthen, daß eine rechte Ordnung gemacht würde in einer Stadt, und solch Schinden und Schaben abgestellt, also, daß der arme Mann zukommen könnte. Wo die Kaufleute oder Handwerksleute sich nicht weisen wollten lassen, so wollte ich einen solchen Rath geben, daß ein Bürgermeister die Fleischer, Bäcker, Brauer, Schenken etc. vor sich forderte und sie ermahnte, daß sie recht mit der Sache umgingen und ohne Falschheit ihre Waare verkauften, und also machten, daß es nicht zuletzt an ihren Kindern ausginge. Denn wo sie also wollten schinden und schaben, sollen sie wissen, daß GOtt solch unrecht, böse gewonnen Gut nicht läßt gedeihen, noch den Nachkommen zugute kommen, wie oben [§ 253] gesagt ist. Denn wie es herkommt, so geht es wieder hin.

257. Wo nun solcher Rath und treue Ermahnung nicht helfen wollte, sondern wollten also nach ihrem Gefallen stolzieren und die Leute trotzen, müßte man thun, wie in etlichen Städten Sitte und Gewohnheit ist, daß man einen redlichen, frommen Mann aufwürfe, dem ein Rath zwei- oder dreihundert Gulden vorstreckte, auf daß er eine ganze Stadt mit Fleisch oder Brod versiehe, und der Rath ihm vergönnte, daß er allein in der Stadt schlachtete, damit solche stolzen Gesellen gedemüthigt würden, und nicht also stolzieren, wie wir hören und erfahren müssen. Mit der Weise könnte man sie Mores lehren, daß sie nicht so trotzig sagten: Willst du es nicht, so lasse es liegen.

258. Also müßte man auch andern Handwerken mitfahren, die in gleicher Gestalt wollten stolzieren und eine ganze Gemeinde pochen, so würden sich andere vorsehen, und nicht so trotzig sagen: Man muß es wohl haben. Weil man ihnen ihren Muthwillen gestattet,\footnote{Erlanger: ihren Muth läßt.} machen sie es mit ihrer Hantierung nach all ihrem Gefallen, schinden und schaben ohn alle Scheu, und meinen noch dazu, sie habens Recht. Ich frage dich: Bist du nicht ein ärgerer Dieb, wenn du also nach deinem Willen fährst mit deinem Verkaufen, denn ein öffentlicher Straßenräuber? Ja freilich; vor einem solchen kann man sich doch vorziehen und ihn fliehen, oder sich desselben wehren; dazu steht derselbige Straßenräuber alle Stunden in Fährlichkeit seines Lebens, denn wenn er ergriffen wird, muß er herhalten; vor dir aber kann sich niemand hüten. Dazu stehst du in keiner Fährlichkeit, sondern lebst, mucherst und stiehlst ohne Furcht, und willst noch dazu ein redlicher Mann gescholten sein. Darum bist du ein Dieb und Räuber, ja, viel ärger denn ein Strauchdieb. Gedenke doch auch an dein Gewissen und Seelen Seligkeit.\footnote{„und Seelen Seligkeit“ fehlt in der Erlanger.}

259. So ist nun dieses ein weitläufigig Gebot, daran alle Menschen schuldig sind. Darum sehe sich ein jeglicher eben vor, wie er handelt. Denn ob einer gleich nicht also stehle, daß man mit Fingern auf ihn zeigt, so ist doch genugsam gesagt, wie weitläufigig dies Gebot sei, und wie die ganze Welt darüber sündigt.

\section{Das fünfte Gebot [der andern Tafel]}

\begin{schriftzitat}
V. 16. Du sollst nicht falsch Zeugniß geben wider deinen Nächsten.
\end{schriftzitat}

\subsection{Summa dieses Gebots}

260. Das fünfte Gebot der andern Tafel lehrt, wie man sich halten soll gegen des Nächsten zeitliche Ehre und gut Gerücht, daß man das nicht schwäche, sondern mehr schütze und erhalte.

\begin{schriftzitat}
Du sollst nicht falsch Zeugniß reden wider deinen Nächsten.
\end{schriftzitat}

Dieses Gebot ist vornehmlich zu verstehen von dem, wie es vor Gericht zugeht, und quillt aus dem Diebstahl heraus. Und wie das andere Gebot der andern Tafel verbiet den Schaden am Leibe, das dritte an der Person, dem Nächsten zugehörig, das vierte den Schaden am Gute deines Nächsten, daß du ihn nicht übervortheilst oder vortheilest: also verbietet dies fünfte Gebot den Schaden an der Ehre und Namen deines Nächsten, daß man\footnote{„man“ fehlt in der Wittenberger.} niemand seinen Namen soll verletzen und beschädigen.

261. Dieses Gebot quillt aus den dreien vorigen. Denn wenn einer verletzet wird an seinem Leib, an seinen zugehörigen Personen
\quellseite{569}
und an seinem Gute, so erhebt sich denn Zank und Hader, und am allermeisten um das Gut. Und das geschieht in der ganzen Welt, wo nicht der Heilige Geist wohnt.

262. Wie nun alle Menschen Diebe und Ehebrecher sind, also sind sie auch an diesem Gebote schuldig. Weil nun die Uebertretung dieses Gebots so groß ist, muß man öffentliche Gerichte halten, daß man wehre, auf daß [es] nicht noch ärger zugehe. Darum ist dies Gebot gut für den Unschuldigen, daß ihm nicht Gewalt und Unrecht vor Gericht geschehe.

263. Darum soll man in solchem Falle beide Theile öffentlich verhören, und dem Unschuldigen nicht Unrecht thun oder Unrecht lassen thun, weder Gaben noch Geschenke nehmen, nicht Freundschaft oder Gunst ansehen. Aber das Laster ist in der Welt gemein, und geht in vollem Schwange. Da nimmt man Gut und Geld, macht aus Recht Unrecht. Man findet auch selten einen frommen Fürsten, Richter oder Juristen, der darinnen nicht strauchelt, ja, der nicht zu einem Buben drüber wird. Denn man hat mehr Aufsehens auf die großen, gewaltigen Hansen denn auf die armen Leute.

264. Es gehört ein großer Muth dazu, und ein kühner Mann, der in dem Falle sein Amt recht\footnote{„recht“ fehlt in der Wittenberger.} ausrichtete. Denn wo GOttes Gnade in einem Richter nicht wohnt, so thut er seinem Amte nimmermehr genugsam, fällt seinem Freunde und guten Gönner, oder sonst einem großen Hansen zu; sieht also durch die Finger, und spricht ein falsch Urtheil wider das andere Theil, da kein Ansehen, Gewalt und Freundschaft ist. Als, wenn ein Armer kommt, davor man sich nicht darf fürchten, keines Schadens gewarten, der muß herhalten. Das haben die Heiden fein angezeigt durch ein Gleichniß eines Spinngewebs: Wenn die kleinen Fliegen drein kommen, bleiben sie darinnen behangen, die großen Hummeln aber fahren hindurch und zerreißen das Gewebe. Also ist es auch vor Gericht: Wenn ein armer Mann kommt, da man keines Schadens besorgt, der muß sich leiden; wenn man sich aber besorgt, es möge einer sich rächen, der fährt hindurch, ob er gleich siebenmal Unrecht hätte.

265. Also geht es zu mit falschem Zeugnisse, das ein gemein Laster ist in der Welt. Denn der Mensch, der nicht den Heiligen Geist hat, hänget am Gute, will das nicht verlieren; da hebt sich denn der Hader, da schwüre einer wohl zehn Eide, ehe er wiche. Den Gebrechen hat GOtt an uns allen gesehen, daß wir falsche Zeugen sind.

266. Darum sollen die Regenten keine Sache urtheilen, und zum Ende bringen, sie haben denn beide Theile gehört. Es sei der Kläger so gewaltig er immer wolle, er schmücke auch seine Sache, wie hoch er kann, da soll ein Richter sprechen: Ich habe zwei Ohren, und du hast Einen Mund. Was du redest, das fasse ich mit dem einen Ohre; was aber der Mund deß, den du anklagst, redet, fasse ich mit dem andern Ohre.

267. So weiß nun GOtt, wer wir sind, hält uns nicht für besser, denn wenn wir vor Gericht kommen, daß wir unsere Sache schön machen, schmücken und vertheidigen, daß recht sei, des andern aber muß stinken. In Summa: Es ist niemand so fromm, der seine Sache schände. So wir nun von Natur also gesinnt sind, daß ein jeglicher das Seine sucht, und fragt niemand nach dem andern, darum soll man nicht einem Theile glauben, man habe denn den andern Theil auch gehört. Also ist ein jeglicher, als ferne er ein Mensch ist, ein falscher Zeuge, Verräther, Lügner, und geht kein wahr Wort aus seinem Maule, nämlich, so es zum Treffen kommt. Es kann wohl geschehen, daß du nicht falsch Zeugniß redest wider deinen Nächsten, weil er mit dir lacht, und ist guter Dinge; wenn es aber zum Ernst kommt, daß dein Nächster einen Zuspruch zu dir hat, oder zu deinem Gute etc., da findet es sich, da wirst du nicht ohne Lügen sein, es gebe dir es denn GOtt in das Herz; sonst machst du dich fromm, gerecht, und schmückst deine Sache auf das allerbeste, es sei der andere Theil wie fromm und gerecht\footnote{Erlanger: wie.} er immer wolle, und deine Sache sei so faul und böse als sie immer kann.

268. Das ist nun aber einmal ein schöner Titel und Ruhm, damit uns GOtt abmalt, schilt uns Lügner, falsche Zeugen, Verräther und Bösewichter. Darum wir billig vor ihm erschrecken sollen, daß er von uns nichts Besseres hält, denn daß wir allzumal Lügner und falsche Zeugen sind. Denn wo er uns nicht dafür hielte, so hätte er nicht ein solch Gebot gegeben. Also
\quellseite{570}
lernen wir aber hier, was GOtt von uns hält; ja, es fühlt’s ein jeglicher bei sich in seinem Herzen.

269. Darum ist auch das gemeine Sprüchwort wohl wahr, daß allenthalben in der Welt geht: Es ist keine Treue in der Welt. Item: Treue ist mißlich. Es darf der Vater dem Sohne, der Sohn dem Vater, die Tochter der Mutter, die Mutter der Tochter nicht vertrauen. Und wie ein Heide spricht: Die Brüder sind selten eins. Darum ist es nicht neu, daß einer betrogen und belogen wird, es ist allewege also gewesen, wir bringen es aus Mutterleibe. Das ist gewiß. Derhalben, so du mit einem Menschen handelst, es sei mit Kaufen oder Verkaufen, ist GOtt mit ihm, so ist er fromm, und fährt redlich mit dir; wo ihn aber GOtt läßt, und ist lauter Mensch da, ohne GOttes Gnade, so ist er lügenhaftig und betrügt dich, wo er kann. Wenn du nun mit ihm willst handeln, so gedenke also: GOtt wird ihm Gnade geben, daß er mir Glauben halten wird. Denn die Liebe hält das Beste vom Nächsten. Aber gleichwohl siehe dich vor; Treue ist mißlich, wie ein griechischer Fürst,\footnote{„ein griechischer Fürst“ fehlt in der Wittenberger.} Epicharmus, gesprochen hat: Gedenke, daß du niemand glaubest. Dieser\footnote{Erlanger: Der Fürst.} sahe, was der Mensch, was die Welt für ein fromm Kind war, nämlich darinnen keine Treue und Glauben ist.

270. Summa Summarum: Wir müssen trauen, und auch nicht trauen, und wissen, zum ersten, daß die Welt voller Teufel, böser Buben und Schälke ist, und wir unter Wölfen und wilden, unvernünftigen Thieren wohnen, die uns betrügen und täglich anlügen. Welcher nun das nicht leiden will, der gehe an einen andern Ort, da keine Welt ist, denn er findet es allenthalben also in der Welt, wie auch Paulus 1 Cor. 5, 10. spricht: Wo ihr nicht wollet unter bösen Buben wohnen, müßtet ihr aus der Welt gehen. Wir dürfen uns anderes und Besseres nicht versehen; der Mensch bleibt Mensch, die Welt bleibt Welt, der Teufel bleibt Teufel. Wer sich nun darein schicken kann, das ist eine große Gnade; wer es aber nicht kann, da ist Jammer und Noth. Zum andern, daß wir wissen, daß man dem Nächsten frei dahin traue, und es in die Schanze schlage, es treffe, oder nicht. Trifft es, daß er dir Glauben hält, so ist es eine Gnade von GOtt; trifft es nicht, daß er dir nicht Glauben hält, und betrügt dich, so hast du vorhin gewußt, daß [es] von Natur besser mit ihm nicht wäre.

271. So redet nun dies Gebot: Du sollst nicht falsch Zeugniß geben, vornehmlich davon, daß, wenn man vor Gerichte mit dem Nächsten zu schaffen hat, es treffe an Gut, Ehre, oder was es sei, man\footnote{Wittenberger und Jenaer: daß man.} rechtschaffen handle. Dahin gehören auch alle Lügen her, die da dem Nächsten Schaden bringen. Aber die Lügen, damit man scherzt und narrt, das sind nicht rechte Lügen; denn man weiß, daß dem Nächsten kein Schade daraus folgt. Wenn ich aber Schaden will mit Lügen meinem Nächsten, es sei am Leibe, Gut, Ehre, oder Namen, das heißen rechte Lügen.

272. Item, daher gehören auch die unnützen, wäschigen Mäuler, die die Leute gerne zur Bank hauen und austragen, wenn es schon nicht vor Gericht geschieht, sondern gegen andern Leuten in geheim, oder in Collation.\footnote{So die Jenaer: Collation = Gesellschaft, Gasterei. Wittenberger und Erlanger: Collatien.} Als die thun, so man heißt Afterredner, die da Lust haben, andern Leuten Uebels nachzureden, haben sonst nichts zu schicken, denn daß sie die Leute ausrichten, und ist niemand hübscher, denn sie; sie sind allein rein, und jedermann stinkt vor ihnen. Wenn ein anderer in Sünde fällt, so muß es jedermann wissen, tragen es hin und her, verleumden also ihren Nächsten, schänden und lästern auf das allerärgste.

273. Wo man von andern Leuten Laster reden wollte, wenn es recht zuginge, so sollte solches geschehen dem Nächsten zur Besserung, nicht zu Schaden und Nachtheil; wo ich es aber nicht bessern kann, soll ich schweigen, und meines Nächsten Gebrechen zudecken; wo ich es aber bessern kann, so soll ich mit ihm allein davon handeln, wie Christus lehrt Matth. 18, 15.: „Strafe deinen Bruder zwischen dir und ihm allein“; das ist denn eine heimliche Besserung. Die andere Besserung ist öffentlich; so du was weißt von deinem Nächsten; so gehe an den Ort, da man es bessern kann. Als, wenn deines Nächsten Kind, Sohn oder Tochter sündigt, strafe es zum ersten; darnach, wo du nichts ausrichtest, sage es seinen Eltern, oder dem Pfarrherrn, Bürgermeister, Richter, oder denen, die
\quellseite{571}
Macht haben zu strafen, so ist es nicht unrecht; ja, du thust recht daran, so du auf die Weise deines Nächsten Sünde meldest; nicht, daß du Lust daran habest, sondern allein zu seiner Besserung. Aber die bösen Mäuler thun das Widerspiel, reden nicht von ihres Nächsten Fehl und Gebrechen, daß er dadurch gebessert soll werden, sondern, daß sie sich damit kitzeln, Lust und Freude darin haben, gedenken nicht einmal, daß sie ihres Nächsten Sünde damit strafen oder bessern wollten.

274. Solche Leute sind rechte Säue und Dreckfresser. Denn des Nächsten Sünde ist nicht anders, denn Koth und Dreck. Welche nun also Lust haben an anderer Leute Sünde, die sind wie die Säue, warten auf den Dreck, lassen also den Koth durch die Zähne und Maul laufen. Denn sie fassen hin und her der Andern Sünde auf, daß sie ihr Maul damit spülen, wie die Säue auf der Gasse Dreck auffressen; haben also ihre Lust und Freude daran, ist ihnen nicht leid, daß ihr Nächster also in Sünde oder Schande fällt, afterreden und malen darnach die Sünde auf das schärffste ab, können der Sache fein helfen, und viel ärger machen, denn sie an ihr selber ist. Es wäre ihnen auch leid, daß ihr Nächster fromm bliebe, denn sie könnten ihre Lust nicht büßen, und das Maul nicht waschen. Solche Leute sind allein schön und rein, die andern alle sind unrein. Ihr Ding muß recht sein, und lauter Balsam, aber anderer Leute Ding ist bei ihnen lauter Teufelsdreck.

275. Nun, ihrer sind wenig, die nicht schuldig sind an diesem Gebote, denn es ist weitläuftig, und ist das Laster gemein allenthalben. Und ob schon einer nicht von anderer Leute Sünde redet, so hört er es zum wenigsten gerne, oder hindert es nicht, wo er es hört.

276. Also habt ihr kürzlich den Begriff dieses Gebots. Es verbeut allen Schaden, der dem Nächsten geschieht mit dem Maule, oder mit der Zunge, damit man ihm Schaden thut an seinen Ehren oder gutem Gerüchte. Auf dies Gebot haben sie auch gezogen die Fälschung der Briefe, aber es gehört mehr zum andern Gebot. Dies Gebot ist so weitläuftig, daß alle Welt darin ersoffen ist.

277. So sind nun die zehn Gebote ein Spiegel, darinne wir sehen, daß die Welt nichts anders ist, denn ein böser, verzweifelter Haufe, der GOtt nicht vertraut, nicht glaubt, wenig

von ihm hält, strebt Tag und Nacht wider ihn mit allen Kräften, lästert GOtt, und schmäht seinen heiligen Namen, achtet der Werke GOttes nicht, sondern vielmehr wirft sein eigen Werth auf, verachtet Vater und Mutter und alle Obrigkeit, ist voll Mords und Würgens, voll Dieberei und Ehebruchs, voll falsches Zeugniß und Lügen. Solches läßt GOtt also von uns singen und sagen. Darum sollte billig unser Herz erschrecken, und im Leibe erbeben, wenn wir hören, daß GOtt, die höchste Majestät, ein solch Urtheil über uns schleußt.

\subsection{Beschluß}

278. Also ist verboten, allen Gütern des Nächsten zu schaden, und geboten, denselben zu frommen. Wenn wir nun das natürliche Gesetz ansehen, so finden wir, wie billig und gleich alle diese Gebote sind. Denn nichts ist hier geboten gegen GOtt und den Nächsten zu halten, das nicht ein jeglicher ihm wollte gehalten haben, wenn er an GOttes und seines Nächsten Statt wäre.

\section{Das sechste und siebente Gebot der andern Tafel}

\begin{schriftzitat}
V. 17. Du sollst dich nicht lassen gelüsten deines Nächsten Haus. Du sollst dich nicht lassen gelüsten deines Nächsten Weibes, noch seines Knechts, noch seiner Magd, noch seines Ochsen, noch seines Esels, noch alles, das dein Nächster hat.
\end{schriftzitat}

\subsection{Summa}

279. Die letzten zwei Gebote lehren, wie böse die Natur sei, und wie rein wir von allen Begierden des Fleischliches und der fremden\footnote{„der fremden“ fehlt in der Erlanger.} Güter sein sollen. Aber da bleibt Krieg und Arbeit, dieweil wir leben.

\begin{schriftzitat}
Du sollst dich nicht lassen gelüsten etc.
\end{schriftzitat}

280. Etliche theilen die zwei Gebote. Es liegt nicht viel an der Theilung, St. Paulus faßt es in eines zum Röm. Cap. 7,\footnote{Erlanger: „Röm. am 17.“} 7., da er spricht: „Du sollst dich nicht lassen gelüsten.“ So wird nun da die Lust verboten.

281. Die Juden haben das Gebot also glossirt und ausgelegt: Die Gedanken, wie man pflegt
\quellseite{572}
in einem Sprüchworte zu sagen, sind zollfrei. Als, wenn einer also wünscht: Wäre der Ochse, die Kuh, der Esel, das Haus, der Acker, das Gut mein. Diese Gedanken haben sie für keine Sünde gerechnet; die andern Sünden der andern Gebote haben sie\footnote{„für keine Sünde ... haben sie“ fehlt in der Wittenberger.} dahin gedeutet, so sie äußerlich mit Werken und mit der That geschehen. Als, so einer einen Abgott aufrichtet, oder ein Bild macht, den Namen GOttes lästert, ein leiblich Werk am Sabbath thut, Vater und Mutter schlägt, mit dem Schwert seinen Nächsten erwürgt, bricht bei Nacht in das Haus und nimmt ihm sein Gut heraus, schläft ihm bei seinem Weibe, und gibt falsch Zeugniß wider ihn vor Gericht.

282. Aber die letzten zwei Gebote haben sie also gedeutet: Wenn sich einer mit einem äußerlichen Zeichen erzeigt, und vollbringt es doch nicht. Als, so einer Vater und Mutter nicht gehorsam wäre; wenn er ihnen gleich nicht flucht, oder sie nicht schlüge, höbe doch einen Knüttel auf, oder ergriffe einen Stein, oder höbe die Faust wider sie auf, und schlüge sie nicht, thäte doch dergleichen, als ob er wollte schlagen oder werfen; das Zeichen wäre wider das Gebot, nach der Juden Auslegung. Also, des Nächsten Weib schänden, ist wider das sechste Gebot; aber wenn es nicht ins Werk gebracht wird, sondern wird nur dergleichen gethan mit einem Zeichen, als: so einer bei eines andern Weibe ergriffen würde in einer Kammer, säße bei ihr auf dem Bette, das Geberd und Zeichen haben sie auf dies Gebot geführt, ob gleich das Werk nicht geschieht. Also, stehlen mit der Faust ist wider das siebente Gebot; wenn aber einer ergriffen wird auf einem Boden, im Keller, in der Kammer, bei dem Kasten, läßt sich merken mit Zeichen, als wollte er stehlen, und des Nächsten Gut entwenden, und hat es doch noch nicht hinweg, der sündigt wider dies Gebot.

283. Also haben es die Juden gedeutet; aber sie haben weit gefehlt, und unrecht glossirt. Denn es wird hier nicht also, und soll auch also nicht gedeutet werden. Moses war nicht so ungelehrt, daß er nicht wußte, wie er reden sollte. Darum gilt die Glosse gar nichts. So wissen wir auch wohl, was Lüsten und Begehren heißt. Es steht nicht in der Faust, oder in den äußerlichen Geberden, sondern im Herzen; wie hin und wieder die Schrift die Lust oder Begierde nicht auf die Faust oder andere Glieder zeucht, sondern auf das Herz. Als, wenn geschrieben steht, die Begierde ihres Herzens. Also deutet es Christus auch recht Matth. 5, 28.: „Welcher eines andern Weib siehet, und begehret ihr (nicht mit dem Munde, Augen, Hand oder Füßen, sondern im Herzen), der ist ein Ehebrecher“, ob er gleich kein Zeichen der Lust von sich gibt äußerlich. Also steht das Gebot im Herzen, ob es gleich nicht heraus bricht mit Zeichen und Werken. Denn diese Lust steckt in unsern Herzen, Mark und Beinen.

284. So er nun spricht: „Du sollst dich nicht lassen lüsten deines Nächsten Weib“, soll also verstanden werden, wenn es einer also begehrt, daß er es ihm selber wollte haben, GOtt gebe, was sein Nächster hätte; vergönnt also seinem Nächsten nicht, daß er ein fromm, ehrbar Weib, ehrsame und züchtige Kinder habe. Also soll es verstanden werden mit andern Dingen, als, mit Knechten und Mägden, Haus und Hof, Ochsen, Kühen, Eseln und andern Gütern; wenn es einem leid ist, daß sein Nächster so ein fromm Gesinde, so ein schön Haus hat, so gute Acker und Wiesen, fette Ochsen und Schafe hat; er wollte, daß er es allein hätte, ob gleich sein Nächster nichts hätte. Also, daß es nicht steht allein in dem, daß einem leid ist, daß sein Nächster so köstlich Ding hat, sondern, daß er es auch wollte selber haben und besitzen, wo er mit Fug möchte dazu kommen.

285. Also haben wir die zehn Gebote, die uns alle zu Schälken und Buben machen vor GOtt. Und ob gleich das weltliche Schwert äußerlich wehrt, und einen Unterschied macht in der Welt, daß man spricht: Der ist fromm, der ist ein Schalk, so sind gleichwohl alle Menschen mit einander Sünder im Herzen. Darum darf sich keiner berühmen, daß er fromm sei vor GOtt, sondern muß bekennen, daß er an allen Geboten schuldig sei, wie sie lauten. Es greife ein jeglicher in seinen Busen, so wird er es also erfinden. Darum sollen wir schreien und heulen zu GOtt, daß er uns den Geist gebe, daß wir nicht allein äußerlich fromm werden vor der Welt, sondern auch\footnote{„auch“ fehlt in der Wittenberger.} vor GOtt im Herzen. Amen.
